\documentclass[12pt,a4paper]{article}

\usepackage{mathtools}
\usepackage{mathrsfs}
\usepackage{float}
\usepackage{caption}
\usepackage{tikz}
\usepackage{tikz-cd}
\usetikzlibrary{arrows,positioning,calc,decorations.pathmorphing}

%% Macros-v2_198: current project v3.072 III-D and bibliography binding; mathematical macros unchanged.
%% Current project v3.070 source-reference bindings; previous macro bodies unchanged.
% Current II-D reference root for PapersMinimal v3.063; mathematical macros otherwise unchanged.
% Macros-v2_188: re-based current paper identities and A--G family; no change to geometry.
% Macros-v2_184: active Paper III-E v1.200/III-F v1.267 binding for framework v3.031; mathematical macros unchanged.
% Macros-v2_183: updated shared cross-reference bindings for Paper1B/III-E/III-F v3.030; mathematical macros unchanged.
% Macros-v2_182: active Paper II-F/V-A cross-references for framework v3.029; definitions unchanged.
% Paper-series v2.381: new I-B torsion/modular results, bibliography and current source roots.
% Paper-series v2.380: new exceptional bibliography and synchronized paper references.
% Paper-series v2.379 current 31-paper macro/aux registry.
% Paper-series v2.378: current 31-paper AUX/bibliography registry updated for Paper I-B v1.2 literature additions.
% Framework v2.371: current-source reference registry refreshed; old theorem labels retained.
% Framework v2.370: current 30-paper source/AUX root rebinding; native mathematical macros unchanged.
% Framework v2.368: synchronized merged 30-root registry; mathematical notation unchanged.
% Framework v2.358: active 30-paper source registry synchronized; existing mathematical macros unchanged.
% Framework v2.357: selective new source-owned certificates; active reference roots updated for five revised papers.
% Framework v2.356: active 30-root registry synchronized to the 30 papers in this release.
% Framework v2.275: exact positive dense-twin word cubature, finite repeatability and eighth-moment classification, nested sampling frames, joint action and POVM-jet exports; preserves v2.274.
% Framework v2.274: filename and root-source synchronization.
% Framework v2.273: restored missing Problem 3 theorems from v2.268; exact exceptional frames, twin geometry, convolution repeatability and projector/occurrence firewalls; merged on v2.272 without overwriting independent branches.
% Framework v2.272: merged v2.270 covariant native-readout mathematics with v2.271 Problem 4 finite-action Hopf hypercharge/holonomy corrections; no universal bulk-gauge uniqueness claim.
% Independent branches: v2.270 readout/guard exports and v2.270 K16 action-consumed charge closure; all prior scoped assumptions retained.
% Framework v2.269: synchronized 30 active AUX roots; mathematics and bibliography preserved from Macros-v2_161.
% Framework v2.266: Paper II-C converse-derived outcome resources; II-E, III-C, IV-A/B, V-A and VI-B downstream mathematical closures; 30-root registry synchronized.
% Framework v2.265: II-C coordinate-stage constitution and exact first-symbol host inheritance; prior mathematical domains retained.
% Framework v2.264: paper II-B reverse/audit and IV-A observable-minimal exchange; synchronized external references to 30 active roots.
% Framework v2.263: Paper I-C adds bridge-controlled root-carry parity and source-relative old-cap cancellation; other root mathematical bodies retained; active 30-root references synchronized.
% Framework v2.262: integrates Paper II-F radius, Casimir and readout results; active 30-root references synchronized.
% Framework v2.261: merged Paper II-B GP2-level/geometric repairs; preserves v2.260 II-D repairs; active references synchronized.
% Framework v2.260: Paper II-D signed RS, finite-symbol typing, and conditional curvature/Jordan repairs; references synchronized.
% Framework v2.259: conditional native SU(2) atlas and explicit later discharge; v2.258 mathematical bodies retained elsewhere.
% Framework v2.258: preserves parallel II-A/dynamics update and adds complete readout/frequency closures.
% Framework v2.257: exact minimal complete position readout and arbitrary-preparation frequency law.
% Framework v2.256: exact native inertia and sharp bounded ordered-source closure; preserves concurrent v2.255 work.
% Framework v2.255: preserves parallel detector integration and adds exact finite readout closures.
% Framework v2.254: Pasterski bibliography and synchronized 30-paper external-reference registry.
% Framework v2.253: typed raw coordinate defect in I-A and synchronized VI-A import; active references synchronized.
% Framework v2.252: synchronized source registry; optional III-G shell-frame calculation.
% Framework v2.251: forward common-action execution, joint source jets and scoped coupled control.
% Framework v2.251: exact full-heat variations, fermion source jets and explicit physical completion frontier.
% Framework v2.251: certified interacting source jets, pole stability and scoped completion proofs.
% Framework v2.251: synchronized active source registry; mathematical notation unchanged.
% Framework v2.251: synchronized active source registry; mathematical notation unchanged.
% Framework v2.251: synchronized active source registry; mathematical notation unchanged.
% Framework v2.251: synchronized source registry; no mathematical notation changes.
% Framework v2.244: synchronized 30-paper AUX registry and rooted phase-source action-order tests.
% Framework v2.242: signed rooted phase lift, complete transverse return and exact family cubic.
% Framework v2.241: proved rooted conormal geometry, conditional cap realization and scoped action tests.
% Framework v2.239: all 30 active source roots synchronized; no mathematical macro changes.
% Framework v2.238: Paper IV A/B/C split and 30-root reference registry.
% Framework v2.237: Paper VII-D technical split and 28-root reference registry.
% Framework v2.236: Paper VI-C technical-volume reference added; all current root bindings advanced.
% Framework v2.235: 26-root AUX and bibliography bindings synchronized; mathematical macros unchanged.
% BEGIN INLINE SOURCE: Paper00-TeXBuildCompatibility-v1_0.tex
% Rendering-only repairs for inherited long appendices and bitmap text fonts.
% Integer counters through 26 retain the original labels; later values use AA,
% AB, ... . This file changes no action coefficient, state or theorem premise.
\RequirePackage{alphalph}
\makeatletter
\let\BHPrimitiveAlph\@Alph
\newalphalph{\BHExtendedAlph}{\BHPrimitiveAlph}{26}
\let\@Alph\BHExtendedAlph
\AtBeginDocument{%
  \@ifpackageloaded{microtype}{\microtypesetup{expansion=false}}{}%
}
\makeatother

% END INLINE SOURCE: Paper00-TeXBuildCompatibility-v1_0.tex

% Framework v2.164: current reference/bibliography profile; mathematical definitions retained.
% Framework v2.162: finite metric, functional Hilbert and scoped Ward integration.
% Framework v2.161: inherits Macros-v2_57.tex exactly except the current external-reference binding.
% Macros v2.57 / framework v2.156: current reference registry only; mathematical definitions unchanged.
% Framework v2.143 / Macros v2.55: current external-reference bindings; mathematical notation unchanged.
% Framework v2.110 / Macros v2.54: current external-reference registry and bibliography; mathematical notation inherited.
% Framework v2.67 / Macros-v2_53: adds Born-action-cell and joint-EH readout symbols; absolute unit realization is metrological; MATCH_{J/scal}^{EH} remains the sole G24 normalization bridge.
% Framework v2.35 / Macros-v2_52: grouped paper numbering, active AUX registry, and Born-v6_20 bibliography binding synchronized; notation unchanged.
% Framework v2.26: active-source AUX and bibliography bindings synchronized; notation unchanged.
% Framework v2.09: reviewed early-reciprocity integration; scoped physical imports.
% Framework v2.08: Born--Euler mathematics and conditional complete-action potential criterion.
% Framework v2.07: active AUX registry includes Paper I-B v1.2; shared bibliography Born-v6_18.
% Macros-v2_46 / framework v2.04: mathematical atlas and downstream realization correction.
% Macros-v2_46 / framework v2.03: scoped structure integration and current dependency bindings.
% Macros-v2_46 / framework v2.01: shared source bindings; prior revision comments are historical.
% v2.42 / framework v1.99: Synchronized RG/AS owner and bibliography bindings.
% v2.41 / framework v1.99: active manuscript AUX bindings and geometric-closure bibliography synchronized.
% Release-v1.80 centralized source synchronization for scalar-circle/comparison-band integration.
% Release v1.79: Hurwitz hierarchy and kinetic-normalized action response integration; physical mass matching remains conditional.
% Release v1.78: approved M8 closure-batch integration; physical and analytic domains retained.
% Current v1.78 registry: all active TeX roots are synchronized below.
% Earlier release comments are historical scope records, not current status.
% Release-v1.77 synchronization: native same-owner Noether closure, generator--flow/history typing, analytic-certification firewalls, scalaron/K16, LF/QCD, quantum-Ward, and RAR boundary-owner transport are registered; shared mathematical notation is unchanged.
% Release-v1.76 synchronization: registry points to the v1.76 active paper roots.
% Release-v1.76 synchronization: registers the local A+Phi+g Noether/constraint closure and dynamics-frontier propagation; shared mathematical notation is unchanged.
% Release-v1.75 synchronization: registry points to the v1.75 active paper roots.
% Release-v1.75 synchronization: integrates the joint-Dirac Einstein bulk real form and the local odd-scalar heat-coefficient action on explicit analytic first-jet/interior domains; v1.74 Yang--Mills closure is retained. Fermionic reality, pure-metric curvature-squared/exact-nonlocal terms, noncanonical boundaries, physical vacuum selection, and actual solution-history tokens remain open.
% Release-v1.74 synchronization: registry points to the v1.74 active paper roots and retains the v1.73 theorem labels.
% Release-v1.74 synchronization: canonical Macros-v2_34.tex; generated local nonlinear Yang--Mills action, constrained connection-field phase space, source-free Gauss momentum map, and Legendre closure are integrated on their proved domains. The remaining non-gauge ACT_cl sectors and actual solution-history selection remain open.
% Release-v1.71 synchronization: closes the former JOINT_x/JOINT_5
% connection-occurrence gap by canonical structure-group extension of the
% nonflat Born connection, faithful associated representation, and the local
% cotangent--Clifford operator constructor.  Actual fluctuation members,
% actions, vacua, solutions, calibrations, and analytic QFT domains remain
% separately typed.  The active-AUX registry is advanced once.
% Release-v1.71 canonical macro and active-AUX registry.
% ======================================================================
% Macros-v2_39.tex -- canonical publication macro file
% Release v1.71 retains the v1.70 null-ray/QCD boundaries and additionally
% removes the duplicate structural connection actuality premise.  It does not
% claim global bundle topology, select a gauge/Higgs member, or promote the
% finite construction to an interacting continuum theory.
% Release v1.68 constructs the finite light-front total-fraction operators and
% parton number/momentum operator measures, proves their simplex and Hausdorff-limit
% theorems, and separates positive cut measures from scheme-renormalized
% light-ray distributions.  It also separates finite real Schur poles, stable
% norm-resolvent/Riesz pole-residue limits, and continued resonance
% pole/residue limits.  The Paper-VII-C scattering package is claim-indexed:
% local QFT and spectral claims import no S-matrix or infrared premise, while
% asymptotic amplitudes, LSZ, S-matrix poles, completeness, and cross sections
% import only their named massive or massless/long-range clauses.
% Release v1.67 factorizes BND48_H into a same-action Hessian/Schur
% descendant and still-open typed boundary-trace, critical-lift, and
% refinement-limit conditions.  It proves the finite-collar Dirichlet-to-
% Neumann law sqrt(y)tanh(L sqrt(y)), its nonuniform critical limit, and the
% finite-return ultra-infrared turnover.  The Paper-I-A Hardy constant mode is
% an exact geometric near-hit, not yet the physical metric boundary mode.
% The current local f(R) Hessian remains massive and therefore still refutes
% a RAR prediction on the minimal action branch.
% Release v1.66 derives the transverse depth-six normalizer bridge from the
% already generated C6 Reynolds image and Paper-I-A binary-dihedral normalizer.
% The bridge is the unique least 2D3-stable scalar-containing realization of
% the faithful quotient 2D3/C6=C2, has trace defect six and hence K=48, but is
% transverse to the nested active carrier.  Exact K=44 and K=52 contraction,
% semigroup, and energy-midpoint identities are registered as structural
% readouts rather than projector events.  The active E48 carrier gap remains
% unchanged.  A same-action boundary-symbol/Calderon intertwiner is isolated
% as the necessary and still unproved owner bridge to the static DtN sort;
% representation equivalence alone neither derives RAR nor creates a K=48
% action junction, pole, occupation, or cosmological transition.
% Release v1.65 derives the native E_48-centered carrier mass gap from the
% GP1--GP3-generated active spectrum, distinguishes its carrier-free interval
% from a particle located at E_48, and withdraws the unsupported promotion to
% a positive continuum Yang--Mills mass gap.  The latter remains a valid
% conditional transfer theorem on an enlarged continuum branch, but it is not
% an active framework prediction or an unfinished native-gap proof.
% Release v1.64 synchronizes Derivation Batches 14--22.  It registers the
% GP1--GP3 heat-profile orbit, same-Hessian electroweak response and pole
% ownership, K16/K48 action-response corrections, perturbation state and
% non-Gaussianity typing, finite-vacuum/gap/QCD limit repairs, factorized
% position-area ownership, the actual static metric response, and the ordered
% Wilson/string-tension limits.  Strict continuum convergence, positive
% physical gap, surviving higher cumulants, and positive tension remain open.
% Release v1.61 closes TIMEWIDTH16 as an independent direct-production datum.
% JUNC_{2|7} continues to own the finite intrinsic microscopic profile/width,
% while the actual same-action background/production history determines its
% physical-time pullback, support, and duration.  The current Paper-I-A--VII-C AUX
% registry is advanced once; shared mathematical notation is unchanged.
% Release v1.60 derives the final-octet primitive-remnant depletion, canonical
% Hopf-current representative, exact Hopf descent to dOmega/(4pi), and
% closed-FLRW a^{-3} neutral-current shape.  It retypes K=48 as the onset
% of the final one-unit geometric transfer and records cDM/RAR as distinct
% physical projections of that common neutral parent while preserving
% occurrence, Hilbert-owner, stress-matching, and static-kernel premises.
% Release v1.59 transports the v1.58 exceptional-shell/K16 closure through
% Papers III-F and VI-B--VII-C.  It replaces Paper VI-B's obsolete mixed-Hessian
% output typing by the third-variation pair vertex and Bogoliubov occupation
% support, derives the post-seed Higgs-pair feedback vertex and its exact
% alpha--delta kernel factors, propagates the unique K16/K48 action-junction
% split, identifies the QCD delta count with the upstream 2O shell atom only
% as a common-scalar numerical compatibility, and narrows THERM16 to dynamical
% approach/beta selection because the finite Gibbs/KMS target is already
% unique on a supplied simple observable algebra once its Hamiltonian and
% inverse temperature are supplied; direct-sum central weights remain typed.
% Release v1.58 registers the exceptional nearest-shell closure, the derived
% transverse K=16 action event, the unique interior action junction, the
% type-correct ordered pair vertex, the thin-layer squeeze law, and the
% radiation-like/matter-like kinematic split.  REG_0, transfer magnitude,
% physical layer width, and thermalization remain explicitly conditional.
% Release v1.57 registers the active-depth-six obstruction, the native
% K=48 carrier desert, the typed 7|9 owner reduction, and the dynamical
% reheating handoff across Papers III-B--VI-B.
% Release v1.56 registers the least-support n_2 inflation dressing of the
% persistent n_7 Laurent response, the analytic global-convexity and positive
% endpoint-measure proofs, the derived inflation core and K=16 coefficient
% freeze, and the no-particle dark-sector split into static RAR plus the
% conditional homogeneous effective-curvature response at reflected K=48.
% There is no n_6 action, scalaron-DM mass, or second Einstein--Hilbert term.
% Release v1.55 registers the protected-octet resolved/residual reflection
% q_sp+Z_W=8, its j<->8-j and K<->64-K scale covariance, and the derived
% rank-one-exit quadratic grade C_1^2/U=C_2 in Papers III-C, III-E, and VIII.
% The reflection is a typed trace/scale theorem, not a time reversal or a
% physical event-occurrence theorem; reheating transfer remains conditional.
% Release v1.54 registers the exact positive-part Laurent indices, the
% active-domain/terminal-endpoint action split, the no-second-EH firewall, and
% the q_DM=6/K_DM=48 scalaron readout typing across Papers III-D/E/F, VI,
% VIII, and IX.
% Release v1.53 registers the common nonnegative Laurent obstacle projection
% for n_1^+ and n_7^+, discharges the three former exit/freeze assumptions,
% and advances the Paper-III-D/E/F, VIII, and IX consumers.  Shared notation
% and bibliography ownership are unchanged.
% Release v1.51 registers the PQTI derivation batch and its updated
% Paper-III-A/C/D/E/F, VIII, and IX consumers.  Shared notation and
% bibliography ownership are unchanged.
% Release v1.50 advances the Paper-III-B AUX target to v1.40 after the
% rigorous audit repair; shared notation and bibliography ownership are unchanged.
% Release v1.49 advances the Paper-II-E/F AUX target to v3.14 after the
% minimal audit repair; shared notation and bibliography ownership are unchanged.
% Release v1.46 synchronizes Paper II-A's minimal frozen constructor-origin
% theorem and its five dependent AUX targets.  Shared notation is unchanged.
% Release v1.45 updates the Paper-II-B AUX target to v2.23 after the
% GP3-rigidity and appendix-placement repair; shared notation is unchanged.
% Release v1.44 advances the centralized Paper-I-A--VII-C AUX registry after the
% Paper-I-A v5.16 protected-completion scope repair.  Every current
% paper imports this one file; no paper-level opt-in or \externaldocument list
% is permitted.  The bibliography remains centralized here through Born-v6_13.
% Shared bibliography, theorem infrastructure, colors, notation alignment,
% and current Paper-I-A--VII-C external-reference registration.
% This file centralizes common imports, bibliography sorting, colors, theorem
% infrastructure, and the notation shared by the complete series.
% v2.9 (release v1.43) made the Paper-I-A--VII-C AUX registry automatic and
% centralized the bibliography and all external-reference ownership.
% v2.8 release-registry refresh (v1.33) registers the rooted generated-parent
% cascade and the POL/REG/free-law/11+4/Born-group closure revisions.  No
% shared mathematical notation or Paper-II-A fundamental principle is changed.
% v2.8 release-registry refresh (v1.32) registers the universal generated-
% continuation revisions in Papers II-A, III-A, III-C, IV, and VIII.  No
% shared mathematical notation or Paper-II-A fundamental principle is changed.
% v2.8 release-registry refresh (v1.30) registers the CRP4/projective-state
% source revisions and their direct registry-only consumers.  No shared
% mathematical notation or Paper-II-A axiom is changed.
% v2.7 release-registry refresh (v1.27) updates only the current external
% Paper-I-A--VII-C basenames after the placement/assumption synchronization; no
% shared mathematical notation or Paper-II-A axiom is changed.
% v2.7 synchronizes the geometric defect-repair revision: Paper-II-A v3.42,
% Paper-III-C v1.26, Paper-III-D v1.27, Paper-III-E v1.32, and Paper-III-F
% v1.34.  It registers the common trace-defect/spectral-filtration theorem,
% its owner-typed response pullback, and the reduced GRED and TERM interfaces.
% No shared mathematical notation is changed.
% v2.6 synchronizes Paper-III-E v1.31's cone-type repair, minimal-response
% classification, and explicit classical action-role boundary, together with
% Paper-III-F v1.32's downstream response-atlas synchronization.  No shared
% mathematical notation is changed.
% v2.5 synchronizes the Paper-III-D v1.26 faithful-trace-duality proof,
% exact ordered-owner joint locus, and typed alternative-presentation
% bookkeeping repair, together with its Paper-III-E and closure-frontier
% references.  No shared mathematical notation is changed.
% v2.4 synchronizes the Paper-III-C v1.25 orbit/process/owner repair, its
% downstream Papers III-B--F, V, VI, VIII, and IX revisions, and Born-v6_13.
% v2.3 synchronizes the Paper-II-A v3.41 witness/support/principal inheritance
% layer, the Paper-III-A v1.24 joint-realizability/readout firewall, all dependent
% host-rate revisions, and the current Paper-I-A--VII-B source registry.
% v2.2 synchronizes the generated active-realization, determinant-depth,
% dimension-weighted hierarchy, and typed cosmological-event interfaces.
% v2.1 synchronizes the exceptional-support interface revisions: Paper I-A
% v5.9, Paper II-A v3.30, Paper II-E/F v3.7, and Paper IV v4.16.
% v2.0 introduces the stage-safe G0 occurrence notation, synchronizes the
% active Paper-I-A--IV source registry (current rigorous Paper-III closure:
% III-A v1.17, III-B v1.22, III-C--E v1.17, III-F v1.14, IV v4.14), and makes
% external-reference self-skip depend on logical paper role rather than a
% versioned job name.
% ======================================================================

% ======================================================================
% macros.tex — Shared macros for the Born Geometry series
% ======================================================================

% ======================================================================
% 1. PACKAGES, BIBLIOGRAPHY, COLORS, AND UNIT DEFINITIONS
% ======================================================================
% Core packages shared across the complete Paper-I-A--VII-C series.  Loading them here makes
% the macro file a genuine series preamble while remaining compatible with existing
% paper preambles that load the same packages before \input.
\usepackage{amsmath}
\usepackage{amssymb}
\usepackage{amsthm}
\usepackage{mathrsfs}
\usepackage{array}
\usepackage{booktabs}
\usepackage{xcolor}
\makeatletter
\@ifpackageloaded{geometry}{}{\usepackage{geometry}\geometry{margin=1in}}
\makeatother
\usepackage{mdframed}

% Shared support packages required by the macros below and by most papers.
\usepackage{ragged2e}
\usepackage{longtable}
\usepackage{tabularx}
\usepackage{enumitem}
\usepackage[most]{tcolorbox}

% Series-wide numeric bibliography convention using the portable BibTeX path.
% All Papers I-A--VII-C use the ordinary \cite interface.
\usepackage[numbers,sort&compress]{natbib}
\providecommand{\printbibliography}{%
  \bibliographystyle{plainnat}%
  \bibliography{Born-v6_60}%
}

% Shared colors.  \providecolor avoids clobbering a paper-specific override.
\providecolor{postulatecolor}{RGB}{230,242,255}
\providecolor{compare}{RGB}{230,245,230}
\providecolor{definitioncolor}{RGB}{230,242,255}

\providecolor{openriskcolor}{RGB}{255,235,235}
\providecolor{bridgecolor}{RGB}{240,240,255}
\providecolor{gold}{RGB}{255,215,0}

% ======================================================================
% 2. CUSTOM ENVIRONMENTS
% ======================================================================
% Common table column types used in all papers.  Guarded so that a paper may
% define its own versions before loading this file.
\makeatletter
\@ifundefined{NC@find@L}{%
  \newcolumntype{L}[1]{>{\RaggedRight\arraybackslash}p{#1}}%
}{}
\@ifundefined{NC@find@C}{%
  \newcolumntype{C}[1]{>{\Centering\arraybackslash}p{#1}}%
}{}
\@ifundefined{NC@find@Y}{%
  \newcolumntype{Y}{>{\RaggedRight\arraybackslash}X}%
}{}
\makeatother

% ======================================================================
% 3. THEOREM STYLES AND ENVIRONMENTS
% ======================================================================
% Custom underlined theorem style (useful for remarks or notes)
\newtheoremstyle{underlined}
  {3pt}   % Space above
  {3pt}   % Space below
  {}      % Body font (empty means normal/roman text, typical for remarks)
  {}      % Indent amount
  {}      % Theorem head font (can be \itshape, \bfseries, etc.)
  {.}     % Punctuation after theorem head
  { }     % Space after theorem head
  {\underline{\thmname{#1}\thmnumber{ #2}}\thmnote{ (#3)}} % Theorem head formatting

% Standard plain theorem style environments.  Each environment is guarded so
% the shared macro file can be loaded after a paper-specific theorem setup.
\theoremstyle{plain}
\makeatletter
\@ifundefined{theorem}{\newtheorem{theorem}{Theorem}[section]}{}
\@ifundefined{lemma}{\newtheorem{lemma}{Lemma}[section]}{}
\@ifundefined{proposition}{\newtheorem{proposition}{Proposition}[section]}{}
\@ifundefined{corollary}{\newtheorem{corollary}{Corollary}[section]}{}
\makeatother

% Definition style environments (typically non-italicized body)
\theoremstyle{definition}
\makeatletter
\@ifundefined{definition}{\newtheorem{definition}{Definition}[section]}{}
\@ifundefined{postulate}{\newtheorem{postulate}{Postulate}[section]}{}
\@ifundefined{closure}{\newtheorem{closure}{Closure}[section]}{}
\makeatother

% Applying the custom underlined style
\theoremstyle{underlined}
\makeatletter
\@ifundefined{remark}{\newtheorem{remark}{Remark}[section]}{}
\@ifundefined{convention}{\newtheorem{convention}{Convention}[section]}{}
\@ifundefined{principle}{\newtheorem{principle}{Principle}[section]}{}
\makeatother

% Framed proof environment.
% Every \begin{proof} ... \end{proof} is automatically framed/shaded.
% This avoids mdframed/trivlist nesting, which can drop the top border after
% theorem-like environments.
\makeatletter
\renewenvironment{proof}[1][\proofname]{%
  \par\addvspace{6pt}%
  \noindent
  \begin{tcolorbox}[%
    enhanced,
    breakable,
    colback=black!4,
    colframe=black!35,
    boxrule=0.35pt,
    arc=1pt,
    left=6pt,
    right=6pt,
    top=4pt,
    bottom=4pt,
    boxsep=0pt,
    before skip=0pt,
    after skip=6pt,
    parbox=false
  ]%
  \pushQED{\qed}%
  \normalfont\itshape #1\@addpunct{.}\enspace
  \normalfont\ignorespaces
}{%
  \unskip\nobreak\hfill\popQED
  \end{tcolorbox}%
  \par\addvspace{2pt}%
}
\makeatother

\mdfdefinestyle{postulatebox}{
  backgroundcolor=yellow!10,
  linecolor=black,
  linewidth=0.4pt,
  skipabove=\baselineskip,
  skipbelow=\baselineskip,
  nobreak=true
}

% ======================================================================
% PROOF-READABILITY HELPERS
% ======================================================================

\newcommand{\ProofPlan}[1]{%
  \par\smallskip
  \noindent\textbf{Proof plan. }#1%
  \par\smallskip
}

\newlist{proofsteps}{enumerate}{2}
\setlist[proofsteps,1]{
  label=\textbf{(\alph*)},
  ref=\alph*,
  leftmargin=2.2em,
  itemsep=0.45\baselineskip,
  topsep=0.35\baselineskip
}
\setlist[proofsteps,2]{
  label=\textbf{(\roman*)},
  ref=\roman*,
  leftmargin=2.0em,
  itemsep=0.25\baselineskip,
  topsep=0.2\baselineskip
}


% ======================================================================
% 4. NUMBER SYSTEMS AND ALPHABETS
% ======================================================================
\providecommand{\RR}{\mathbb{R}} % Real numbers
\providecommand{\PP}{\mathbb{P}} % Projective space / Probability
\providecommand{\CC}{\mathbb{C}} % Complex numbers
\providecommand{\ZZ}{\mathbb{Z}} % Integers
\providecommand{\NN}{\mathbb{N}} % Natural numbers
\providecommand{\HH}{\mathbb{H}} % Quaternions / Hyperbolic space
\providecommand{\LFRayGeom}{\mathsf{LF\!\mbox{-}RAY\!\mbox{-}GEOM}}
\providecommand{\LFRep}{\mathsf{LF\!\mbox{-}REP}}
\providecommand{\LFState}{\mathsf{LF\!\mbox{-}STATE}}
\providecommand{\LFLim}{\mathsf{LF\!\mbox{-}LIM}}
\providecommand{\LFScat}{\mathsf{LF\!\mbox{-}SCAT}}
% Number systems / projective spaces
\providecommand{\CP}{\mathbb{CP}}
\providecommand{\RP}{\mathbb{RP}}

% ======================================================================
% 5. GROUPS AND LIE THEORY
% ======================================================================
\providecommand{\SO}{\ensuremath{\mathrm{SO}}}     % Special Orthogonal group
\providecommand{\SU}{\ensuremath{\mathrm{SU}}}     % Special Unitary group
\providecommand{\Spin}{\ensuremath{\mathrm{Spin}}} % Spin group
\providecommand{\CS}{\ensuremath{\mathrm{CS}}}     % Chern-Simons (capitalized)
\providecommand{\cs}{\ensuremath{\mathrm{cs}}}     % Chern-Simons (lowercase)

% ======================================================================
% 6. COMMON MANIFOLDS AND TOPOLOGY
% ======================================================================
\providecommand{\Sone}{S^1}                 % Standard circle

% ======================================================================
% 7. MATH OPERATORS (Declared via \DeclareMathOperator)
% ======================================================================
\DeclareMathOperator{\End}{End}   % Endomorphisms
\DeclareMathOperator{\Aut}{Aut}   % Automorphisms
\DeclareMathOperator{\spec}{spec} % Lowercase spectrum
\DeclareMathOperator{\tr}{tr}     % Lowercase trace (if you already have \Tr)
\DeclareMathOperator{\supp}{supp} % Support
\DeclareMathOperator{\TrA}{Tr_{A}} % Partial trace over subsystem A
\DeclareMathOperator{\TrB}{Tr_{B}} % Partial trace over subsystem B
% Lie/algebra operators
\DeclareMathOperator{\Ad}{Ad}
\DeclareMathOperator{\Sym}{Sym}
\DeclareMathOperator{\su}{su}
% ======================================================================
% 8. STANDARD OPERATOR-LIKE SYMBOLS (Safe definitions via \providecommand)
% ======================================================================
\providecommand{\vol}{\operatorname{vol}}   % Volume form
\providecommand{\Vol}{\operatorname{Vol}}   % Total volume
\providecommand{\Ric}{\operatorname{Ric}}   % Ricci tensor
\providecommand{\Rm}{\operatorname{Rm}}     % Riemann tensor
\providecommand{\Tr}{\operatorname{Tr}}     % Uppercase trace
\providecommand{\Spec}{\operatorname{Spec}} % Uppercase spectrum
\providecommand{\diag}{\operatorname{diag}} % Diagonal matrix
\providecommand{\rank}{\operatorname{rank}} % Rank
\providecommand{\scal}{\mathcal R} % Scalar curvature
\providecommand{\Cl}{\operatorname{Cl}}     % Clifford algebra
\providecommand{\diam}{\operatorname{diam}} % Diameter

% ======================================================================
% 9. FUNCTIONS, TRIGONOMETRY, AND CALCULUS
% ======================================================================

% Codifferential naming conventions
% (Avoid \ddag: commonly defined by physics.sty and symbol packages)
\providecommand{\codiff}{\mathrm{d}^{\dagger}} % Codifferential operator

% Quantum mechanics bra-ket notation (needed for basis tables)
\providecommand{\ket}[1]{\lvert #1 \rangle}
\providecommand{\bra}[1]{\langle #1 \rvert}
\providecommand{\braket}[2]{\langle #1 \mid #2 \rangle}

% ======================================================================
% 10. PHYSICAL SCALES AND CONSTANTS (Paper 1c / Born Geometry)
% ======================================================================
\providecommand{\Rp}{\ensuremath{R_{\mathrm P}}}     % matched common EH curvature scale (not defined from observed G)
\providecommand{\hBorn}{\ensuremath{h_{\mathrm B}}} % physical representative 2\pi\hbar of one primitive Born scalar action cell
\providecommand{\UEHjoint}{\ensuremath{U_{J}}} % joint-Dirac EH curvature-coordinate scale 8 kappa_J
\providecommand{\Gread}{\ensuremath{G_{\rm scal}}} % scalar-action coordinate c_op^3/(h_B U); physical matching explicit
\providecommand{\GreadJoint}{\ensuremath{G_{J,\mathrm{raw}}}} % raw ordinary-trace coordinate = G_scal/96
\providecommand{\Rds}{R_{\mathrm{dS}}}    % de Sitter Curvature
\providecommand{\Mpl}{M_{\mathrm{Pl}}}    % Planck mass
\providecommand{\lp}{\ell_{\mathrm{Pl}}}  % Planck length
\providecommand{\Pmax}{P_{\mathrm{max}}}  % Maximum momentum
\providecommand{\Mplbar}{\bar{M}_{\mathrm{Pl}}} % Reduced Planck mass
\providecommand{\eff}{\mathrm{eff}}       % Subscript for effective quantities
\providecommand{\Cosmo}{\ensuremath{\Lambda}}                % Cosmological constant (use \Cosmo in source)
\providecommand{\LCDM}{\ensuremath{\Cosmo\mathrm{CDM}}}  % \LambdaCDM (typesets with \Cosmo)
\providecommand{\Lcut}{\ensuremath{\Lambda_{\mathrm{cut}}}}  % Cutoff scale (brace when used as script)
\providecommand{\Rcs}{\ensuremath{R_{\mathrm{cs}}}}      % Configuration-space radius
\providecommand{\CosmoEff}{\ensuremath{\Cosmo_{\eff}}}   % Effective cosmological constant

% Disambiguation: \Rmom is the momentum-space (S^3) radius. 
% Reserve plain R for the 4D curvature scalar in f(R).
\providecommand{\Rmom}{\ensuremath{R_{\mathrm{mom}}}} 
\providecommand{\Punit}{P_{0}} % Momentum unit set by \Rmom (= \hbar/\Rmom)
\providecommand{\Rvec}{\ensuremath{R_{\mathrm{vec}}}} % invariant vector-space radius in the reduced round sector
% Specific numeric values and expressions
\providecommand{\Vlogcap}{L_{\mathrm{Hopf}}^{(1)}}
\providecommand{\Hopf}{\operatorname{Hopf}}
% Scalar phase (S^1) sector + universal cover notation (Paper I-A)
\providecommand{\tper}{T_\star}      % Time period/scale
\providecommand{\Eper}{E_\star}      % Energy period/scale
\providecommand{\ttilde}{\tilde t}   % Universal cover time
\providecommand{\Etilde}{\tilde E}   % Universal cover energy

% ======================================================================
% 10A. SERIES-WIDE LOGICAL / NOTATIONAL DISAMBIGUATIONS
% ======================================================================
% Use these symbols to avoid collisions among: (i) the Paper--III
% no-insertion/geometric-closure class, (ii) Paper--I-A trimer diagonals,
% (iii) Paper--V-A/V-C family-space diagonals, (iv) scalar-index running,
% and (v) protected shell-clock labels.
\providecommand{\Rprim}{\ensuremath{\mathsf R_{\mathrm{prim}}^{\Sigma}}} % primitive finite structural readouts
\providecommand{\Rloop}{\ensuremath{\mathsf R_{\circ}^{\Sigma}}} % closed-loop readout role
\providecommand{\Rbdry}{\ensuremath{\mathsf R_{\partial}^{\Sigma}}} % boundary-conditioned readout role
\providecommand{\Rgen}{\ensuremath{\mathsf R_{\mathrm{fin}}^{\Sigma}}} % finite generated closure of primitive roles
\providecommand{\PhysOcc}{\operatorname{PhysOcc}}
\providecommand{\Actual}{\operatorname{Actual}}
\providecommand{\GeomLic}{\operatorname{GeomLic}}
\providecommand{\InvDist}{\operatorname{InvDist}}
\providecommand{\Dcl}{\ensuremath{\mathsf{D}_1}}      % unmarked positive Dirac-even closed spectral-readout form class
\providecommand{\Dheat}{\ensuremath{\mathsf D_{\mathrm{heat}}}} % one-parameter heat-semigroup subclass
\providecommand{\DtriI}{\ensuremath{D^{\mathrm{tri}}_1}} % Paper--I-A cyclic-trimer diagonal
\providecommand{\DtriII}{\ensuremath{D^{\mathrm{tri}}_2}}
\providecommand{\DfamI}{\ensuremath{D^{\mathrm{fam}}_1}} % Family-space diagonal selector
\providecommand{\DfamII}{\ensuremath{D^{\mathrm{fam}}_2}}
\providecommand{\alpharun}{\ensuremath{\alpha_{\mathrm{run}}}} % scalar spectral-index running
\providecommand{\Ksh}{\ensuremath{K_{\mathrm{sh}}}} % protected shell-clock coordinate
\providecommand{\ddef}{\ensuremath{\delta_{\mathrm{def}}}} % universal defect scalar after Paper--III realization

% ======================================================================
% 10B. SERIES-WIDE LOGICAL STATUS TAGS
% ======================================================================
% Upright small-caps wrapper for audit/status labels.  It prevents LaTeX from
% trying to use unavailable small-caps italic font shapes inside theorem bodies
% such as TU/lmr/m/scit.
\providecommand{\AuditLabel}[1]{}
\renewcommand{\AuditLabel}[1]{{\normalfont\scshape #1}}

% These tags are bookkeeping labels for consistency audits. They do not
% create new mathematical hypotheses; they record how a datum is being used
% in the proof chain.
\providecommand{\StatusPrimitive}{}
\renewcommand{\StatusPrimitive}{\AuditLabel{Primitive}}
\providecommand{\StatusDerived}{}
\renewcommand{\StatusDerived}{\AuditLabel{Derived}}
\providecommand{\StatusImported}{}
\renewcommand{\StatusImported}{\AuditLabel{Imported}}
\providecommand{\StatusConditional}{}
\renewcommand{\StatusConditional}{\AuditLabel{Conditional}}
\providecommand{\StatusBridge}{}
\renewcommand{\StatusBridge}{\AuditLabel{Bridge-derived}}
\providecommand{\StatusComparison}{}
\renewcommand{\StatusComparison}{\AuditLabel{Comparison-only}}

% ======================================================================
% 10C. SERIES-WIDE CLAIM-WORDING AUDIT LABELS
% ======================================================================
% These are textual audit labels used to keep modal language consistent
% across the series.  They introduce no new mathematical assumptions.
\providecommand{\ClaimForced}{}
\renewcommand{\ClaimForced}{\AuditLabel{forced}}
\providecommand{\ClaimUnique}{}
\renewcommand{\ClaimUnique}{\AuditLabel{unique}}
\providecommand{\ClaimCanonical}{}
\renewcommand{\ClaimCanonical}{\AuditLabel{canonical}}
\providecommand{\ClaimMinimal}{}
\renewcommand{\ClaimMinimal}{\AuditLabel{minimal}}
\providecommand{\ClaimClosed}{}
\renewcommand{\ClaimClosed}{\AuditLabel{closed}}
\providecommand{\ClaimPrediction}{}
\renewcommand{\ClaimPrediction}{\AuditLabel{prediction}}

% ======================================================================
% 10D. SERIES-WIDE THEOREM-CHAIN AUDIT LABELS
% ======================================================================
% These labels mark the direction of load-bearing theorem chains. They are
% audit annotations only; they introduce no new mathematical assumptions.
\providecommand{\ChainInput}{}
\renewcommand{\ChainInput}{\AuditLabel{Chain input}}
\providecommand{\ChainOutput}{}
\renewcommand{\ChainOutput}{\AuditLabel{Chain output}}
\providecommand{\ChainGuard}{}
\renewcommand{\ChainGuard}{\AuditLabel{Chain guard}}
\providecommand{\ChainNoBackflow}{}
\renewcommand{\ChainNoBackflow}{\AuditLabel{No backflow}}

% ======================================================================
% 10E. SERIES-WIDE LOCAL THEOREM-AUDIT LABELS
% ======================================================================
% These labels mark a statement-level dependency audit around individual
% theorem/definition nodes. They are bookkeeping labels only.
\providecommand{\LocalAudit}{}
\renewcommand{\LocalAudit}{\AuditLabel{Local audit}}
\providecommand{\UsesOnly}{}
\renewcommand{\UsesOnly}{\AuditLabel{Uses only}}
\providecommand{\Exports}{}
\renewcommand{\Exports}{\AuditLabel{Exports}}
\providecommand{\DoesNotUse}{}
\renewcommand{\DoesNotUse}{\AuditLabel{Does not use}}

% ======================================================================
% 10F. SERIES-WIDE NUMERICAL-ROW AUDIT LABELS
% ======================================================================
% These labels classify individual rows in numerical, benchmark, status,
% and comparison tables.  They are audit annotations only and do not change
% any theorem, value, formula, or comparison number.
\providecommand{\RowAudit}{}
\renewcommand{\RowAudit}{\AuditLabel{Numerical-row audit}}
\providecommand{\RowTheorem}{}
\renewcommand{\RowTheorem}{\AuditLabel{Theorem-fixed row}}
\providecommand{\RowBridge}{}
\renewcommand{\RowBridge}{\AuditLabel{Bridge-fixed row}}
\providecommand{\RowRepresentative}{}
\renewcommand{\RowRepresentative}{\AuditLabel{Representative row}}
\providecommand{\RowComparison}{}
\renewcommand{\RowComparison}{\AuditLabel{Comparison row}}
\providecommand{\RowCalibrated}{}
\renewcommand{\RowCalibrated}{\AuditLabel{Calibrated row}}
\providecommand{\RowEstimate}{}
\renewcommand{\RowEstimate}{\AuditLabel{Estimate row}}
\providecommand{\RowOpen}{}
\renewcommand{\RowOpen}{\AuditLabel{Open row}}
\providecommand{\RowNoSelector}{}
\renewcommand{\RowNoSelector}{\AuditLabel{No selector}}

% ======================================================================
% 10G. TABLE HELPERS
% ======================================================================
% Table helpers used by Paper V-C.  The unused readout-label preset has been
% removed; only active table helpers remain.
\providecommand{\gpsection}[1]{%
  \addlinespace[0.55em]
  \multicolumn{3}{@{}l}{\textsc{#1}}\\[-0.25em]
  \midrule
}
\providecommand{\mathcell}[1]{%
  \(\begin{gathered}#1\end{gathered}\)%
}

\newcommand{\DerBy}[1]{%
  \begin{tabular}[t]{@{}l@{}}#1\end{tabular}%
}

\newcommand{\TableNotes}[1]{%
  \par\medskip
  \begin{minipage}{0.94\linewidth}
  \footnotesize\emph{Notes.} #1
  \end{minipage}\par
}
% ======================================================================
% 10H. PAPER VI-A / QCD NOTATION
% ======================================================================
% Shared symbols used by Paper VI-A and downstream QCD/Yang--Mills handoffs.
\providecommand{\Pqcd}{\mathfrak P_{\rm VIII}^{\rm QCD}}
\providecommand{\Iqcd}{\mathfrak I_{\rm QCD}^{(7)}}
\providecommand{\muQCD}{\mu_{\rm QCD}^{\rm int}}
\providecommand{\LambdaQCD}{\Lambda_{\rm QCD}}
\providecommand{\LambdaBH}{\Lambda_{\rm BH}^{\rm QCD}}
\providecommand{\HQCD}{\mathcal H_{\rm QCD}^{\rm BH}}
\providecommand{\HQCDfin}{H_{\rm QCD}^{\rm fin}}
\providecommand{\Hraw}{\mathcal H_{\rm raw}}
\providecommand{\HphysQCD}{\mathcal H_{\rm phys}^{\rm QCD}}
\providecommand{\Vcol}{V_C}
\providecommand{\Ps}{P_{\rm sing}}
\providecommand{\PiOct}{\Pi_8}
\providecommand{\MSbar}{\overline{\rm MS}}
\providecommand{\YM}{\mathrm{YM}}
\providecommand{\Ktri}{K_{\triangle}}
\providecommand{\KYM}{K_{\YM}^{\rm BH}}
\providecommand{\aBH}{a_{\rm BH}}
\providecommand{\epsBH}{\varepsilon_{\rm BH}}
\newcommand{\RowConstructive}{\textnormal{\textsc{Constructive}}}
\newcommand{\RowObligation}{\textnormal{\textsc{Bridge obligation}}}
\newcommand{\RowTarget}{\textnormal{\textsc{Target row}}}
\newcommand{\RowValidation}{\textnormal{\textsc{Validation row}}}
\providecommand{\slashed}[1]{\not{#1}}

% ======================================================================
% 10I. PAPER VI-B / CONTINUOUS-QFT NOTATION
% ======================================================================
% Shared symbols used by Paper VI-B and any downstream continuous-QFT or
% Yang--Mills handoff paper.  All definitions are guarded so individual
% papers may override them locally before loading this file if needed.
\providecommand{\A}{\mathfrak A}
\providecommand{\Hphys}{\mathcal H_{\rm phys}}
\providecommand{\QCD}{\mathrm{QCD}}
\providecommand{\BH}{\mathrm{BH}}
\providecommand{\loc}{\mathrm{loc}}

\providecommand{\RBH}{\mathcal R_{\rm BH}}
\providecommand{\DG}{\mathfrak D_G}
\providecommand{\PW}{\mathcal{P\!W}} % Peter--Weyl spectral-symbol subspace
\providecommand{\Toe}{\operatorname{Toe}} % Toeplitz compression/representative
\providecommand{\Hol}{\operatorname{Hol}} % Holonomy

% ======================================================================
% 11. LATE-LOADED PACKAGES (Must be at the end of the preamble)
% ======================================================================
% Load late packages only if a paper has not already loaded them, avoiding
% option clashes in paper-specific preambles.
\makeatletter
\@ifpackageloaded{hyperref}{}{\usepackage{hyperref}}
\makeatother
% --- Series-wide disambiguations (Pass 3) ---
\providecommand{\diracdelta}{\delta_{\rm D}} % Dirac delta distribution (delta reserved for hierarchy)
\providecommand{\Hg}{H_{\mathrm g}} % Hodge harmonic projector (keep H for Hubble)
\providecommand{\Gr}{G_{\mathrm r}} % Green operator (keep G for Newton)
\providecommand{\FD}{\mathrm{D}} % functional/variational derivative symbol (delta reserved)
\providecommand{\Id}{\mathbf{1}} % identity/Kronecker tensor (avoid \delta_{ij})


% BEGIN INLINE SOURCE: Paper00-ActiveExternalReferences-v2_130.tex
% Framework v2.191: current reference binding; existing definitions retained.
% Framework v2.190: current22-root registry with native clock-evolution proofs.
% Framework v2.189: exact current22-root external registry, including retained v2.188 hardening owners.
% Framework v2.187: current registry synchronization only; source release v2.186 preserved.
% Framework v2.186: authoritative current 22-paper registry.
% Framework v2.182: all current 22 logical paper AUX bindings.
% Framework v2.170: current I-B v1.43 external AUX binding.
% Framework v2.167: exact current 22-root registry after CP--BL2 integration.
% Framework v2.166: exact current 22-root compact-release AUX registry.
% Framework v2.165: authoritative current root bindings.
% Framework v2.164: exact current 22 manuscript AUX bindings.
% Framework v2.162: finite metric, functional Hilbert and scoped Ward integration.
% Framework v2.161: current 22 manuscript AUX bindings; source only.
% Framework v2.156 / registry v2.113: synchronized current 22 manuscript roots; source-only registry.
% Framework v2.143 / registry v2.111: active manuscript AUX basenames; no generated AUX or PDF included.
% ======================================================================
% Central external-reference imports, numbers only, no PDF jumping
% ======================================================================

% xr-hyper is loaded centrally so that papers may retain their existing
% \usepackage{xr-hyper} lines without redefining \externaldocument after this
% registry is installed.  If plain xr was loaded earlier, retain it.
\makeatletter
\@ifpackageloaded{xr-hyper}{}{%
  \@ifpackageloaded{xr}{}{\usepackage{xr-hyper}}%
}
\makeatother

% Make external-document registration idempotent.  The registry identifies a
% logical paper by its Paper<number><optional letter> prefix.  Prefix-qualified
% imports remain independent.  Current paper roots contain no local
% \externaldocument calls; this guard also makes accidental duplicates harmless.
\ExplSyntaxOn
\prop_new:N \g_bh_external_refs_prop
\cs_new_eq:NN \BHOriginalExternalDocument \externaldocument

\cs_new_protected:Npn \bh_external_document:nnnn #1#2#3#4
 {
  \tl_set:Nn \l_tmpa_tl {#3}
  \regex_replace_once:nnN
   { \A (?: xref/ )? (Paper[0-9]+[A-G]?) .* \Z }
   { \1 }
   \l_tmpa_tl
  \tl_set:Nx \l_tmpb_tl { #1 | \tl_use:N \l_tmpa_tl }
  \prop_if_in:NVTF \g_bh_external_refs_prop \l_tmpb_tl
   {
    \PackageInfo{BornHopf}
     {Skipping~duplicate~external~refs~for~#3}
   }
   {
    \prop_gput:NVn \g_bh_external_refs_prop \l_tmpb_tl {#3}
    \IfNoValueTF{#4}
     { \BHOriginalExternalDocument[#1][#2]{#3} }
     { \BHOriginalExternalDocument[#1][#2]{#3}[#4] }
   }
 }

\RenewDocumentCommand{\externaldocument}{O{} O{} m o}
 { \bh_external_document:nnnn {#1}{#2}{#3}{#4} }

\NewDocumentCommand{\BHExternalRefs}{m}
 {
  % Do not import the aux file of the logical paper currently being compiled,
  % even when the job and registry use different version suffixes.
  \tl_set:Nx \l_tmpa_tl { \c_sys_jobname_str }
  \regex_replace_once:nnN
   { \A (Paper[0-9]+[A-G]?) .* \Z }
   { \1 }
   \l_tmpa_tl
  \tl_set:Nn \l_tmpb_tl { #1 }
  \regex_replace_once:nnN
   { \A (Paper[0-9]+[A-G]?) .* \Z }
   { \1 }
   \l_tmpb_tl
  \str_if_eq:eeTF { \tl_use:N \l_tmpa_tl } { \tl_use:N \l_tmpb_tl }
   {
    \PackageInfo{BornHopf}
     {Skipping~external~refs~for~#1~because~this~is~the~current~document}
   }
   {
    \IfFileExists{xref/#1.aux}
     { \externaldocument[][nocite]{xref/#1}[#1.pdf] }
     {
      \IfFileExists{#1.aux}
       { \externaldocument[][nocite]{#1} }
       {
        \PackageInfo{BornHopf}
         {External~aux~file~#1.aux~not~found;~no~labels~imported}
       }
     }
   }
 }
\ExplSyntaxOff

% External references are automatic.  This file is the sole owner of both the
% bibliography and the current Paper-I-A--VII-D AUX registry (30 roots including I-C).  Missing auxiliaries
% are informational on the first pass and are populated by the source-only
% two-pass series build.  Basenames deliberately omit temporary download/copy
% suffixes such as ``(1)'' or ``-1''.
  % Paper I-A: round S3 compendium.
  \BHExternalRefs{Paper1A-Geometry-v5_181}

  % Paper I-B: complete exceptional geometry and finite spectral structures.
  \BHExternalRefs{Paper1B-ExceptionalGeometry-v1_25}
  % Paper I-D: technical companion of I-A through I-C.
  \BHExternalRefs{Paper1D-TechnicalMaterial-v1_69}
  % Paper I-C: reciprocal M8 compendium.
  \BHExternalRefs{Paper1C-M8Geometry-v1_146}

\ifdefined\BHMathOnly
  % Papers I-A and I-B import only the mathematical atlas labels.
\else
  % Paper II: five physical-mathematical main papers plus one technical proof volume.
  \BHExternalRefs{Paper2A-GeometricConstitution-v3_227}
  \BHExternalRefs{Paper2B-MomentumCarrier-v2_147}
  \BHExternalRefs{Paper2C-ReciprocalLocalization-v2_176}
  \BHExternalRefs{Paper2D-SymmetrySpectra-v2_145}
  % The former II-F export labels remain in II-E; current II-F is the technical volume.
  \BHExternalRefs{Paper2E-KinematicExports-v3_140}
  % Paper II-F: complete technical proof and calculation owner for II-A--II-E.
  \BHExternalRefs{Paper2F-TechnicalMaterial-v1_71}

  % Paper III: Dynamics and Intrinsic Geometric Constitution (six components).
  \BHExternalRefs{Paper3A-GeneratedDynamics-v1_181}
  \BHExternalRefs{Paper3B-SpectralAction-v1_192}
  \BHExternalRefs{Paper3C-ProtectedHierarchy-v1_196}
  \BHExternalRefs{Paper3D-RunningResponse-v1_168}
  \BHExternalRefs{Paper3E-GravityAction-v1_203}
  \BHExternalRefs{Paper3F-CompleteDynamics-v1_279}
  % Paper III-G: complete preserved technical proof owner.
  \BHExternalRefs{Paper3G-TechnicalMaterial-v1_57}

  % Paper IV-A: generated finite quantum mechanics.
  \BHExternalRefs{Paper4A-FiniteQuantumMechanics-v4_217}
  % Paper IV-B: complete finite quantum processes and physical response.
  \BHExternalRefs{Paper4B-QuantumResponse-v1_55}
  % Paper IV-C: extended proofs and finite technical calculations.
  \BHExternalRefs{Paper4C-TechnicalMaterial-v1_50}

  % Paper V-A: The Standard Model.
  \BHExternalRefs{Paper5A-FiniteMatterGeometry-v3_218}

  % Paper V-B: Finite Spectrum and Flavor Closure.
  \BHExternalRefs{Paper5B-TerminalSpectrum-v5_202}

  % Paper V-C: Finite Interaction Boundaries of the Standard Model.
  \BHExternalRefs{Paper5C-InteractionResponse-v2_184}
  % Paper V-D: preserved technical proof and calculation owner.
  \BHExternalRefs{Paper5D-TechnicalMaterial-v1_58}

  % Paper VI-A: Gravity and Closed Cosmology.
  \BHExternalRefs{Paper6A-GravityCosmology-v2_268}

  % Paper VI-B: Geometric Dark Sector and Cosmological Readouts.
  \BHExternalRefs{Paper6B-CosmologyReadouts-v1_215}

  % Paper VI-C: complete technical proof and comparison owner.
  \BHExternalRefs{Paper6C-TechnicalMaterial-v1_53}

  % Paper VII-A: Finite QCD Kernel and Nonperturbative Closure.
\BHExternalRefs{Paper7A-QCDKernel-v1_183}

  % Paper VII-B: Hadronic, Thermal, and Nuclear QCD Readouts.
  \BHExternalRefs{Paper7B-QCDReadouts-v1_167}

  % Paper VII-C: Continuous Quantum Field Theory.
  \BHExternalRefs{Paper7C-ContinuousQFT-v2_208}

  % Paper VII-D: complete technical material for VII-A through VII-C.
  \BHExternalRefs{Paper7D-TechnicalMaterial-v1_60}

\fi

% END INLINE SOURCE: Paper00-ActiveExternalReferences-v2_130.tex


\hypersetup{hidelinks}

\usepackage{xr-hyper}

\title{\textbf{Born--Hopf Geometry III-D: Rank Valuation, Running Response, and Laurent Jets}\\[0.5em]
\large Hopf contact residues, positive mathematical providers, and incidence limits}

\author{Peter Kugelmann\thanks{\texttt{pkugelmann@proton.me}}}
\date{Version 1.168 --- 20 November 2026 \\ Framework v3.072}
\setlength{\emergencystretch}{3em}
\begin{document}
\maketitle
\begin{abstract}
Paper III-D establishes the Hopf contact-residue rank valuation,
reciprocal running response and fixed-depth Laurent-jet relations. Positive
measure/Schur providers are sharply separated from physical source incidence
and from a complete metric variation. Longer auxiliary calculations and
negative countermodels are in Paper~III-G.
\end{abstract}
\tableofcontents
\clearpage

\section{Inputs and protected contact valuation}
\label{paper3d:sec:organized-scope}
\subsection{Introduction and Paper-III series allocation}
\label{p3d:sec:series-allocation}
Paper~III-D imports the protected filtration and matched contact construction
of Paper~III-C and the loop-derived sectorwise spectral form class \(\Dcl\) of
Paper~III-B.  It proves native rank and quotient identities inside the
restricted branch \((\mathrm{RED})\), records ultraviolet and infrared
calibrations separately, and derives the exact two-root response together with
its Riccati/secant and neutral entropy consequences.  Paper~III-E owns any
later signed metric-action placement and distinguishes the auxiliary
secant-matched state from the co-realized hierarchy-active state.

\subparagraph{Imported fixed momentum seed and reciprocal role.}
The theorem-level generated momentum-role seed is the one-copy
six-type carrier \(\mathcal H_1\) imported from Paper~II-B.
Its structural role is co-realized with the upstream transfer-operation
root; this paper needs no additional generated-seed witness.  The
stronger identification of every complete physical momentum endpoint
with \(\mathcal H_1\) is conditional on the GP3 role-faithful
whole-candidate comparison and is not used in the reduced-response proof.
Paper~III-D does not consume \(\mathsf P_{\rm sli}\), momentum-side
TWR/SPEX/FIT, a component ceiling, or a fixed-action selector.  Alternative
momentum-selector routes are external to this reduced-response proof and are treated in
the Paper~II-B constitutional momentum theorem.  Physical reciprocal-role statements
consume the premise-minimal Paper-II-C intrinsic-evaluation/GP2 closure, which
supplies the Born/GP2-generated reciprocal role, universal relation-free product, and
Born exchange involution.  Paper~III-A's metric-curry construction is an
optional stronger G0 full-profile test realization.  The upstream certification tags
\((\mathrm{BR}_{\rm ind})\) and \((\mathrm{IND}_{xp})\) therefore record
upstream certification status rather than independent III-D hypotheses.
Neither this reciprocal constitution nor any III-D response construction
reopens the fixed momentum-seed theorem or supplies the separate local
position--momentum comparison or dimensionful reciprocal calibration.

\subparagraph{Imported G0 actualization boundary.}
All physical interfaces below specialize Paper~II-A's stage-safe G0.  Each
asserted occurrence requires a contingent token of a geometrically licensed,
claim-relative invariant distinction with result-independent provenance; its
tag alone supplies no witness, and no instance reopens frozen upstream
geometry.  A new role or standard enters only through a declared downstream
frontier.
Any later ``physical momentum'' language uses that co-realized role.  A
particular momentum process or readout remains a token-level occurrence.
An alternative or physically actualized support profile
\((\mathrm{ACT}_{\rm Cl})\), the conditioned response readout
\((\mathrm{UNIT})\), and a
realized reduced leaf \((\mathrm{RED})\) are G0 process, occurrence, or
readout certificates.  The tag \((\mathrm{NEUT}_7)\) denotes the generated
visibility family; when it is a constituent of an occurring response,
co-realization supplies its occurrence.  The endpoint roots
\((\mathrm{UV})+(\mathrm{IR})\) and every later dimensionful hierarchy placement are G0
calibrations of already constituted invariants.  The smooth completion retains the separate analytic realization family
\[
 \mathfrak A_{\rm REG}
 =\{\mathrm{REG}_{\sigma}:\sigma_{\rm fin}\in\mathbb R\};
\]
G0 neither selects \(\sigma_{\rm fin}\) nor derives the regularity assumptions
of a scheme member, and it does not select the occurring spectral process
represented by the loop-derived class \(\Dcl\).  The source member
\(\mathrm{REG}_0\) is a declared prescription, not a GP-selected endpoint.

\subparagraph{Protected Clifford and infrared boundary.}
The normalized physical protected real Clifford block requires only the
Paper~II-C nondegenerate stratum
\[
 (\mathrm{ND}_{\rm sp}):\qquad \varrho_{\rm sp}>0,
\]
not a particular point of the unresolved coordinate-completion moduli.  On a
G0-occurring Hurwitz-even Hopf/contact branch, Paper~III-C constructs the
equivariant orbit family of finite active-depth operators and proves that the
common terminal spectral projector is the unique scalar Haar projector.
Denote this orbit-level output by
\[
 \mathsf C_{\rm IR}^{\rm der}
 :=\{(\mathrm{ND}_{\rm sp}),\ \text{occurring active orbit family},\
       \ P_{\rm term}=P_0\}.
\]
The displayed braces abbreviate derived provenance; they are not a new
assumption package.  Finite functional calculus, equivariant descent, and GP2
co-realization pass the active composite's G0 witness to the orbit-invariant
terminal projector without selecting a diagonal.  Thus the native
protected endpoint
\[
 A_{\rm prot}\cong\Cl(0,3),
 \qquad
 \Cl(0,3)=\mathbb R\mathbf1\oplus W,
 \qquad
 \dim_{\mathbb R}W=7,
\]
is understood on \(\mathsf C_{\rm IR}^{\rm der}\).  The tag
\((\mathrm{IR}_{O(3)})\) merely certifies this derived endpoint.  Parity
evenness and full \(O(3)\)-equivariance are conclusions.  The native active
support has depths \(q=0,4,5,7\), so its final step is the rank jump
\(3\to1\) from \(q=5\) to \(q=7\); Paper~III-C,
Theorem~\ref{p3c:thm:active-depth-six-obstruction}, proves that no nested
rank-two carrier occurs.  Separately, an \(SO(3)\)-only average would retain
the pseudoscalar and have a rank-two endpoint at \(q=6\).  That is a
datum-enlarged counterfactual, not the missing native active stage.
The real rank-one scalar-corner tag \((\mathrm{CORN})\) records Paper~III-B's
derived ambient-constituent theorem and is used only when that particular
\(1/8\) coefficient is invoked; it is not needed for the normalized regular
support trace used in the \((\mathrm{UNIT})\) response below.

\subparagraph{Hierarchy and response interfaces.}
Write \((\mathrm{HT})\) only for certification that the relevant physical
readout has already been placed, by a proved generated ancestry morphism, in
Paper~III-C's determinant-dual owner.  It adds no normalization or
scale-by-scale occurrence premise, but it neither proves nor replaces that
ancestry morphism.  The neutral dual character itself requires no physical
owner.  The topological rank valuation does not require any member of the smooth
scheme family.  When the smooth determinant-shell correction is included,
retain the complete current Paper-III-C realization family
\[
 \boxed{\mathfrak A_{\rm REG}
 =\{\mathrm{REG}_{\sigma}:\sigma_{\rm fin}\in\mathbb R\}},
 \qquad
 \delta_\sigma:=\frac1{6\sqrt2}+\sigma_{\rm fin}.
\]
For a frozen member write
\[
 \alpha_\sigma:=\alpha_{\rm hier}(\sigma_{\rm fin}),
 \qquad
 H_{56}^{(\sigma)}:=\alpha_\sigma^{\,7}.
\]
The source prescription is the declared member
\((\mathrm{REG}_0):\sigma_{\rm fin}=0\), with
\(\delta_0=1/(6\sqrt2)\),
\(\alpha_0=\alpha_{\rm hier}(0)\), and
\(H_{56}^{(0)}=\alpha_0^7\).  No bare \(\alpha_{\rm hier}\) denotes a
scheme-independent number, and no generated discriminator selects
\(\mathrm{REG}_0\) from \(\mathfrak A_{\rm REG}\).
Write
\((\mathrm{FLAG})\) for an explicitly supplied ordered rank-one refinement of
\(W\).  The split \(1\oplus7\) itself is derived on
\(\mathsf C_{\rm IR}^{\rm der}\); it is not produced by an
\(O(7)\) symmetry.  Parity/full-\(O(3)\) status is already certified by the
rank-one terminal projector and is not an additional convention.
Paper~III-C now derives the canonical decreasing, left-continuous upper
spectral family
\[
 \mathsf P_u(s)=\mathbf1_{[s,7]}(\mathcal Q_{\rm act}(u))
\]
and its pullback along the smooth trace-defect coordinate.  Write
\((\mathrm{ACT}_{\rm Cl})\) only for an alternative protected-projector law or
for adoption of such a law as an actual external-scale process.  It is
unnecessary for the canonical dimensionless filtration or when a theorem
conditions on one fixed support.  Orbit-invariant generated supports inherit their
occurrence, whereas a representative-dependent member inherits occurrence
only together with its mark; an arbitrary additional support remains a G0
boundary datum.  Write \((\mathrm{UNIT})\) for the adopted
readout that uses the canonical conditioned trace state.  Write
\((\mathrm{NEUT}_7)\) for the already-proved Paper~III-C common-rate
\(O(7)\)-neutral semigroup and its generated visibility/entropy family on the
constituted complement \(W\); it does not reconstruct or constitute the
\(1\oplus7\) split.  Its occurrence follows with any occurring composite that
contains it in a frozen generated-constituent presentation.  The reduced response further uses
\((\mathrm{RED})\) and its two abstract endpoint calibrations
\((\mathrm{UV})+(\mathrm{IR})\).  Native III-D does not identify those roots
with an absolute dimensionful hierarchy reference.  Once a later occurring
response contains a generated ancestry morphism from its upper-root owner to
the weight-one determinant-dual owner, the lower root and every intermediate
ratio are generated, not separately matched.
The residual absolute placement is recorded in
Appendix~\ref{p3d:app:hierarchy-placement} of Paper~III-G.  No MON
occurrence tag is needed: tensor-product independence, when used, remains a
separate composition-domain statement and is not required for the one-shell
active state.

\subparagraph{Derived and conditional interface classification.}
For the present paper the structural outputs
\(\mathcal H_1\) (generated one-copy seed), the protected active orbit and its invariant
supports, the rank valuation, the positive determinant character, its neutral
dual depth character, the
Born role scalarizer, and the \((\mathrm{NEUT}_7)\) visibility family are
derived or generated/co-realized.  G0 occurrence data include an arbitrary
additional fixed support, \((\mathrm{UNIT})\), the realized
\((\mathrm{RED})\) response class, and its endpoint tokens.  Calibrations
include \(C,U\) and any absolute hierarchy reference; a physical hierarchy
placement additionally contains the typed owner-ancestry port.  The unanchored native family is \(\mathfrak A_{\rm REG}\); the occurring Hurwitz branch instead imports the derived \(\mathrm{REG}_0\) of III-C; the stronger
\((\mathrm{FLAG})\) interface and any noncanonical or physically actualized
\((\mathrm{ACT}_{\rm Cl})\) profile remain conditional.  If owner inheritance and the strictly ordered response roots
are jointly imposed, Proposition~\ref{p3d:prop:ordered-owner-joint-locus}
derives \(\mathrm{REG}_\sigma\in\mathfrak A_{\rm REG}^{+}\); membership in
that positive-character subfamily is not a second selector on that branch.
A later independently constituted discriminator may discharge
one of these families only at that later stage; it does not retroactively make
the earlier family unique.

\subparagraph{Terminology.}
A \emph{primitive rank event} in this paper means a real rank-one quotient of
the protected complement.  It is a local dynamical term and is unrelated to
the Paper-II notion of a primitive constructor atom.


\paragraph{Technical proof ownership and constituted scope.}
Every extended auxiliary derivation, explicit matrix verification,
and optional comparison moved during this reorganization is
preserved in Paper~III-G, Section~\ref{p3g:sec:laurent}.
The generated momentum seed \(\mathcal H_1\) is inherited from
Paper~II-B, without later scalar, action, regulator, or source data
reselecting it.  Universal full-endpoint multiplicity-one remains a
separate GP3 comparison obligation.

\section{Hopf contact residues and rank-multiplicative descent}
\label{paper3d:sec:organized-contact}
\subsection{Ordering and Hopf contact-residue valuation}
\label{p3:sec:reeb-universality}

Shell-dependent damping data may determine sharp projector thresholds; they
are separate from the integrated rank valuation derived below.  The
normalized Hopf contact residue supplies the additive rank-proportional
weight, exponentiation supplies its unique positive character, and the
generated common-rate depth supplies its ordered dual scale action.

\begin{definition}[Effective damping rate as ordering parameter]
\label{p3:def:effective-damping-ordering}
A shell flow may assign distinct effective rates to distinct active supports.
These rates determine when a support crosses its threshold, but not the
integrated geometric weight assigned to the completed removal.
\end{definition}

\begin{definition}[Reeb phase and Hopf contact residue]
\label{p3:def:reeb-phase-contact-residue-distance}
For the normalized Hopf connection,
\[
S_{\rm Reeb}^{(0)}=\oint_{S^1}\eta_0=2\pi,
\]
whereas the primitive hierarchy residue is
\[
\Vlogcap
=\int_{S^3}\eta_0\wedge d\eta_0
=\left(\oint_{S^1}\eta_0\right)\left(\int_{S^2}\omega\right)
=(2\pi)^2=4\pi^2.
\]
The metric-scaled orbit action is \(2\pi\Rmom\), but the hierarchy valuation
uses the dimensionless principal connection and its unweighted contact
residue.
\end{definition}

\subsubsection{The Hopf Structure Provides an Intrinsic Rank Valuation}
\label{p3:sec:hopf-rank-valuation}

The previous endpoint result gives the rank drop \(8\to 1\).  The matched
contact residue makes the integrated weight additive in the unresolved
protected rank, with the rank-one atom normalized by the Hopf contact residue.
Paper~III-C's determinant-character theorem and generated dual-depth theorem
give its multiplicative character and exact neutral orientation trichotomy.
The point is
deliberately not that the seven non-scalar directions
are seven distinguished equal channels.  After the scalar coarse-graining of
Proposition~\ref{p3:prop:reeb-selection}, such channel labels are not present in the declared
reduced dynamics.  The physically visible datum is the rank of the quotient removed at a threshold.

\begin{definition}[Unresolved protected rank quotient]
\label{p3:def:protected-rank-quotient}
On \(\mathsf C_{\rm IR}^{\rm der}\), let
\(W=\ker\mathbb E_{\mathrm{IR}}\) be the seven-dimensional non-scalar
complement of Theorem~\ref{p3:thm:hopf-decomposition}.  A completed
shell-removal event is represented by a
finite-dimensional quotient
\[
Q=F_{\mathrm{before}}/F_{\mathrm{after}}
\]
of the active protected non-scalar space.  Its \emph{protected rank} is
\[
\rank_{\mathrm{prot}}(Q):=\dim_{\RR}Q.
\]
A primitive event is an event with \(\rank_{\mathrm{prot}}(Q)=1\).
\end{definition}

\begin{remark}[Abstract orthogonal rank diagnostics deferred]
A general full-orthogonal rank-only valuation theorem is mathematically
available, but the physical protected complement retains generated
Clifford/\(O(3)\) structure and the positive proof does not need that stronger
abstract symmetry assumption.  The diagnostic lemmas are deferred to
Appendix~\ref{p3d:app:supplementary-valuations} of Paper~III-G.
\end{remark}

\begin{lemma}[Matched tensor residue for protected rank atoms]
\label{p3:lem:rank-atom-hopf-trace-localization}
On \(\mathsf C_{\rm IR}^{\rm der}\), let
\(Q\) be any finite real protected quotient, and use the matched tensor
extension exported by Paper~III-C,
\[
 A_Q:=A_{\rm H}\otimes I_Q,
 \qquad A_{\rm H}=8\Delta_{b,0}^{-2}.
\]
Then
\[
 \boxed{\operatorname{Res}_{H}(A_Q)
 =\dim_{\mathbb R}(Q)\,\Vlogcap},
 \qquad \Vlogcap=4\pi^2.
\]
This is a tensor-factorized smooth contact statement, not an identification of
a Clifford line with a Hopf fibre and not an application of the finite
spectral-action class \(\Dcl\).
\end{lemma}

\begin{proof}
Paper~III-C supplies the matched operator and its factorized residue.  The
Heisenberg residue acts on the geometric factor and the finite trace on the
multiplicity factor, hence
\[
 \operatorname{Res}_{H}(A_Q)
 =\operatorname{Res}_{H}(A_{\rm H})\Tr_Q(I_Q)
 =\Vlogcap\dim_{\mathbb R}Q.
\]
The normalization \(\operatorname{Res}_{H}(A_{\rm H})=4\pi^2\) is the
imported Hopf contact-residue theorem.  No finite-Hilbert \(\Dcl\)
classification is used.
\end{proof}

\begin{lemma}[Principal unit-Chern Hopf residue and scheme-labelled smooth completion]
\label{p3:lem:primitive-hopf-cell-normalization}
The proven topological logarithmic weight of one primitive protected rank atom is
\[
L_{1,{\rm top}}=\Vlogcap=(2\pi)^2=4\pi^2.
\]
For any frozen member \(\mathrm{REG}_{\sigma}\in\mathfrak A_{\rm REG}\), the
scheme-labelled completed logarithmic weight is
\[
 L_{1,\sigma}=\Vlogcap+\delta_\sigma,
 \qquad
 \delta_\sigma=\frac1{6\sqrt2}+\sigma_{\rm fin}.
\]
The source determinant-shell prescription is the member
\(\mathrm{REG}_0\), for which the local calculation and
Theorem~\ref{p3:thm:D1-primitive-smooth-uniqueness} give
\[
 \delta_0=\delta_{\rm loc}=\frac1{6\sqrt2}.
\]
No statement here selects \(\sigma_{\rm fin}=0\) from the fixed-shell scheme
family.
\end{lemma}

\begin{proof}
A primitive protected rank atom is, by
Definition~\ref{p3:def:protected-rank-quotient}, a real rank-one nonzero
quotient of the already constituted protected complement on the fixed momentum
seed of Theorem~\ref{p3:thm:factor-local-momentum-shape}.
It is not itself a Hopf cell.  By Lemma~\ref{p3:lem:rank-atom-hopf-trace-localization},
its topological logarithmic threshold contribution is one trace copy of the
principal unit-Chern Hopf contact residue,
\[
\mathcal L_{\rm top}(Q)=\Vlogcap.
\]
By Theorem~\ref{p3:thm:Vlog}, on the provenance-fixed principal unit-Chern
Hopf contact branch,
\[
\Vlogcap=\int_{S^3}\eta_0\wedge d\eta_0=(2\pi)^2=4\pi^2.
\]
This proves the primitive topological step.  Paper~III-C separately retains
the finite-part coordinate \(\sigma_{\rm fin}\), so a frozen smooth member
adds \(\delta_\sigma=1/(6\sqrt2)+\sigma_{\rm fin}\).  The local calculation
in Theorem~\ref{p3:thm:delta_rigorous} and
Theorem~\ref{p3:thm:D1-primitive-smooth-uniqueness} identify the source member
\(\mathrm{REG}_0\) with \(\delta_0=1/(6\sqrt2)\); they do not collapse the
larger fixed-shell scheme family.
\end{proof}


\begin{theorem}[Matched protected-rank residue and hierarchy readout]
\label{p3:thm:reeb-universal}
On \(\mathsf C_{\rm IR}^{\rm der}\), let
\(Q\) be a completed protected quotient with matched tensor operator
\(A_Q=A_{\rm H}\otimes I_Q\).  The contact logarithmic weight is
\[
 \boxed{\mathcal L_{\rm top}(Q)
 =\dim_{\mathbb R}(Q)\,\Vlogcap},
 \qquad \Vlogcap=4\pi^2,
\]
independently of the smooth scheme family.  For a frozen
\(\mathrm{REG}_{\sigma}\in\mathfrak A_{\rm REG}\), the normalized
scheme-labelled smooth completion is
\[
 \boxed{\mathcal L_{\sigma}(Q)
 =\dim_{\mathbb R}(Q)
 \left(\Vlogcap+\frac1{6\sqrt2}+\sigma_{\rm fin}\right)}.
\]
The source member \(\mathrm{REG}_0\) reduces this to
\(\mathcal L_0(Q)=\dim_{\mathbb R}(Q)(\Vlogcap+1/(6\sqrt2))\).
Exponentiation gives the unique positive determinant character for each frozen
scheme member, and algebraic duality gives the generated neutral ratio
\[
 \chi_{\sigma,w}(q)=\alpha_\sigma^{-wq}.
\]
For every \(w>0\), this ratio is strictly decreasing, constant, or strictly
increasing according as
\(L_1^{(\sigma)}=4\pi^2+1/(6\sqrt2)+\sigma_{\rm fin}\) is positive, zero,
or negative.  Hence decrease holds precisely on
\(\mathfrak A_{\rm REG}^{+}\), equivalently
\(\sigma_{\rm fin}>-4\pi^2-1/(6\sqrt2)\).  A physical scale law
\(X_w/X_0=\chi_{\sigma,w}\) additionally requires the owner-ancestry
hypothesis of Paper~III-C's owner-typed corollary.  For
\(\dim_{\mathbb R}W=7\),
\[
 H_{56}^{\rm top}=\alpha_{\rm top}^{7},
 \qquad
 H_{56}^{(\sigma)}=\alpha_\sigma^{7}
 \quad\text{for the frozen }\mathrm{REG}_\sigma.
\]
No dimensionful radius or curvature scale is assigned here.
\end{theorem}

\begin{proof}
Lemma~\ref{p3:lem:rank-atom-hopf-trace-localization} gives the matched
rank-proportional contact residue.  Paper~III-C's scheme-labelled smooth completion is separately rank-additive
for each frozen \(\mathrm{REG}_\sigma\).  Paper~III-C's positive-character
theorem exponentiates the additive logarithmic weight.  Its generated
dual-depth theorem constructs the neutral dual ratio and proves the sign
trichotomy by differentiating its logarithm.  The physical formula is then
exactly the conclusion of Paper~III-C's owner-typed inheritance corollary when
its ancestry premise is met.  The certification \((\mathrm{HT})\) label adds neither
that premise nor a normalization.
The endpoint exponent is seven because \(\dim_{\mathbb R}W=7\).
\end{proof}

\begin{definition}[Fixed protected support and conditioned trace state]
\label{p3:def:active-support-response}
On \(\mathsf C_{\rm IR}^{\rm der}\), let
\(A_{\rm prot}\cong\Cl(0,3)\) be the protected real regular module, and let
\[
 P\in\End_{\mathbb R}(A_{\rm prot})
\]
be one fixed nonzero orthogonal projection of rank
\(k\in\{1,\ldots,8\}\).  If \(P\) is one of the orbit-invariant generated
supports \(\Pi_8,\Pi_4,\Pi_1=P_0\), or the generated complement
\(P_W=1-P_0\), its occurrence is inherited from the III-C active composite by
finite functional calculus and generated-constituent co-realization.  A
rank-three member \(\Pi_3(u)\) is generated only in the equivariant orbit
family and can be used as this particular \(P\) only after the diagonal
\(u\) is marked.  Any
other fixed support used here is an explicit G0 boundary/support occurrence.
No general physical scale-dependent provenance of \(P\) is asserted.
Paper~III-C's canonical profile
\[
 P_{\rm spec}(\tau;u)
 =\mathbf1_{[q_{\rm sp}(\tau),7]}(\mathcal Q_{\rm act}(u))
\]
is nevertheless a generated dimensionless spectral filtration.  Its complete
orbit family requires no \((\mathrm{ACT}_{\rm Cl})\), while a particular
rank-three member still requires \(u\), and an alternative or physically
actualized external-scale family still requires \((\mathrm{ACT}_{\rm Cl})\).

Use the normalized regular trace
\[
 \tau_{\rm reg}(P):=\frac18\Tr_{A_{\rm prot}}(P).
\]
The conditioned regular trace state is
\[
 \omega_P(A):=\frac{\Tr_{A_{\rm prot}}(PAP)}{\Tr_{A_{\rm prot}}(P)},
\]
with density \(h_P\) relative to \(\tau_{\rm reg}\).  The optional readout
\((\mathrm{UNIT})\) adopts this canonical conditioned state as the reduced
scalar response on the fixed support.  Its support projector and unit-trace
normalization make it an \(\Rbdry\) readout, not the extensive unmarked
\(\Dcl\) loop from which the underlying trace operation is imported.
\end{definition}

\begin{lemma}[Unique conditioned density; basis neutrality is automatic]
\label{p3:lem:unit-trace-forces-reciprocal-active-fraction}
For every fixed support projection \(P\) of protected rank \(k\), the conditioned regular trace
state has the unique density relative to the faithful normalized regular trace,
\[
h_P=\frac{P}{\tau_{\rm reg}(P)}=\frac{8}{k}P,
\qquad
\theta(P)=\frac8k.
\]
This density is automatically invariant under every orthogonal change of basis
that preserves the support.  No basis-neutrality premise is needed for its
existence or uniqueness.
\end{lemma}

\begin{proof}
If \(P\) has rank \(k\), then
\[
 \tau_{\rm reg}(P)=\frac18\Tr_{A_{\rm prot}}(P)=\frac{k}{8}.
\]
Set
\[
 h_0:=\frac8kP.
\]
For every \(A\in\End_{\mathbb R}(A_{\rm prot})\), cyclicity and \(P^2=P\)
give
\[
 \tau_{\rm reg}(h_0A)
 =\frac1k\Tr(PA)
 =\frac1k\Tr(PAP)
 =\omega_P(A).
\]
Thus \(h_0\) represents the conditioned state.  If another density \(h\)
represents the same state, faithful trace duality gives
\[
 \tau_{\rm reg}\bigl((h-h_0)A\bigr)=0
 \qquad\text{for every }A.
\]
Taking \(A=(h-h_0)^*\) yields
\[
 \Tr\bigl((h-h_0)(h-h_0)^*\bigr)=0,
\]
so \(h=h_0\).  This proves uniqueness without a commutant assumption.

Finally, if \(O\) is orthogonal and preserves the support, then
\(OPO^{-1}=P\), hence
\[
 Oh_PO^{-1}=\frac8kOPO^{-1}=h_P.
\]
Basis neutrality is therefore a consequence of the unique density rather than
an additional premise.  The uniqueness asserted here is not linearity of the
conditioning map \(P\mapsto h_P\).
\end{proof}

\begin{corollary}[Canonical rank coefficient, support-indexed state, and \((\mathrm{UNIT})\) specialization]
\label{p3:thm:eight-mode-harmonic-response}
Assume \(\mathsf C_{\rm IR}^{\rm der}\) and let \(P\) be a fixed
rank-\(k\) support.  Before any response-role assumption is made, the unique
conditioned trace-state density has scalar support coefficient
\begin{equation}
\label{p3:eq:eight-mode-harmonic-response}
 \boxed{\theta_k=\frac8k},
 \qquad k=1,\ldots,8.
\end{equation}
If \((\mathrm{UNIT})\) is additionally supplied, this already derived density
is adopted as the reduced scalar response on that support.  Thus
\((\mathrm{UNIT})\) supplies the role identification, not the coefficient or
its normalization.  The coefficient \(\theta_k=8/k\) is rank-indexed, whereas the density
\(h_P=\theta_kP\) remains support-indexed.  Equal rank therefore does not
select or identify a support.  No probability distribution over different
ranks and no scale-dependent support family is asserted.
\end{corollary}

\begin{proof}
The protected regular module has real rank eight.  For the fixed projection
\(P\) of rank \(k\),
Lemma~\ref{p3:lem:unit-trace-forces-reciprocal-active-fraction} gives
\(h_P=(8/k)P\) from faithful trace duality, without using basis neutrality or
\((\mathrm{UNIT})\).  The optional
\((\mathrm{UNIT})\) readout supplies only the response interpretation.
\end{proof}

\begin{corollary}[Equivariant active-support inheritance]
\label{p3d:cor:generated-active-support-inheritance}
On the occurring III-C active branch, the invariant supports and the
rank-three orbit family
\[
 \Pi_8,\qquad \Pi_4,\qquad
 \{\Pi_3(u):u\in\mathcal O_{\rm diag}\},\qquad \Pi_1=P_0,
 \qquad P_W=1-P_0
\]
are generated constituents in their native operator and orbit-family sorts.
Hence every invariant member, and the rank-three formula equivariantly for
each \(u\), may be
inserted into Lemma~\ref{p3:lem:unit-trace-forces-reciprocal-active-fraction}
without a second constructor witness, giving
\[
 h_P=\frac8{\operatorname{rank}P}P.
\]
A physical use of one particular \(\Pi_3(u)\) nevertheless requires the
marked diagonal \(u\); orbit generation alone does not select it.
For an arbitrary orthogonal support outside this generated family, the generic
fixed-support theorem remains parameterized by the separately occurring
support.  In particular, equality of ranks never selects an ungenerated
projector.
\end{corollary}

\begin{remark}[Uniform-rank averaging is supplementary]
A uniform average over the eight possible rank values gives
\(H_8=761/280\), but no such probability measure is generated by the
fixed-support theorem.  The calculation and its status are deferred to
Appendix~\ref{p3d:app:supplementary-valuations} of Paper~III-G.
\end{remark}

\begin{remark}[Conditioning, not an inverse-trace action law]
\label{p3:rem:no-hidden-inverse-trace-axiom}
The reciprocal factor is canonical after a support is conditioned:
\[
 h_P=\frac{P}{\tau_{\rm reg}(P)}.
\]
Conditioning preserves linearity in the observable \(A\mapsto\omega_P(A)\),
but it is nonlinear in \(P\).  For orthogonal rank-one supports \(P_1,P_2\),
\[
 h_{P_1}=8P_1,\qquad h_{P_2}=8P_2,\qquad
 h_{P_1+P_2}=4(P_1+P_2).
\]
Thus the factor \(8/k\) is not inferred as a block-additive coefficient of
the uninserted \(\Dcl\) action.  Its physical response use is exactly the
separately displayed \(({\rm UNIT})\) readout.
\end{remark}

\begin{corollary}[Boundary normal form of \((\mathrm{UNIT})\)]
\label{p3d:cor:unit-boundary-normal-form}
For a fixed rank-\(k\) projection \(P\), put
\[
 \rho_P:=\frac{P}{\operatorname{Tr}P}=\frac{P}{k}.
\]
Then \(\rho_P\ge0\), \(\operatorname{Tr}\rho_P=1\), and
\[
 \boxed{\omega_P(A)=\operatorname{Tr}(\rho_P A).}
\]
Thus \((\mathrm{UNIT})\) is the state/effect specialization of the
boundary-conditioned role \(\Rbdry\).  The extensive loop
\(A\mapsto\operatorname{Tr}A\) remains in \(\Rloop\); dividing by
\(\operatorname{Tr}P\) changes the provenance type by retaining the prepared
support state.
\end{corollary}

\begin{proof}
Orthogonal projection gives \(P=P^*=P^2\ge0\) and
\(\operatorname{Tr}P=k\).  Cyclicity and \(P^2=P\) yield
\[
 \operatorname{Tr}(\rho_PA)
 =\frac1k\operatorname{Tr}(PA)
 =\frac1k\operatorname{Tr}(PAP)=\omega_P(A).
\]
This is exactly the finite state boundary normal form; the projector is
retained before evaluation and therefore cannot be erased by rewriting the
coefficient as a trace.
\end{proof}

\begin{remark}[Scope of the universal primitive step]
The theorem does not require \(\mathrm O(3)\) to act transitively on the seven non-scalar directions;
it does not.  Nor does it claim that vector, bivector, and pseudoscalar components have identical
linearized damping rates.  The equality of primitive step weights follows from a different fact:
the matched tensor extension \(A_Q=A_{\rm H}\otimes I_Q\) makes every
real rank-one multiplicity factor contribute the same finite trace
\(\Tr_Q(I_Q)=1\) to the same Hopf residue.  This does not assert equal damping
rates or other channel-resolved dynamics.  The proven analytic value is the
rank-one residue \(\Vlogcap\).  The topological factor
\(\alpha_{\rm top}=e^{\Vlogcap}\) follows directly from the positive
determinant character and does not use \((\mathrm{HT})\).  A scheme-labelled
completed step \(\alpha_\sigma\) additionally uses only the frozen member
\(\mathrm{REG}_\sigma\).
\end{remark}

\begin{proposition}[Ordering versus step size]
\label{p3:prop:ordering-versus-step-size}
On \(\mathsf C_{\rm IR}^{\rm der}\), different non-scalar shells may cross at different response parameters, but
every unresolved primitive rank-one quotient has the same contact residue.
For a frozen \(\mathrm{REG}_\sigma\), this gives the common
scheme-labelled completed factor \(\alpha_\sigma\).  Ordering and the
conditional step readout are logically distinct.  The neutral character needs
no \((\mathrm{HT})\); physical inheritance needs an ancestry map that the
certification label may only certify.
\end{proposition}

\begin{proof}
Different supplied shells or supports may cross at different response
parameters; a channel-resolved model with independently labelled thresholds
would be a datum enlargement.  Theorem~\ref{p3:thm:reeb-universal} gives the common rank-one
residue \(\Vlogcap=4\pi^2\); the positive determinant character turns it into
a multiplicative factor.  A frozen \(\mathrm{REG}_\sigma\) supplies only the
scheme-labelled smooth additive completion.
\end{proof}


\subsubsection{Rank-Multiplicative Protected Descent}

\begin{theorem}[Rank-multiplicative protected endpoint descent]\label{p3:thm:sequential}
On \(\mathsf C_{\rm IR}^{\rm der}\), Theorem~\ref{p3:thm:hopf-decomposition}
and the rank valuation of
Theorem~\ref{p3:thm:reeb-universal}, the ultraviolet-to-infrared hierarchy
carried by the protected non-scalar endpoint has the topological form
\begin{equation}
\text{Total topological hierarchy}
=
\alpha_{\rm top}^{\dim_{\RR}W}
=
\alpha_{\rm top}^7
=
\exp(28\pi^2).
\end{equation}
For any frozen \(\mathrm{REG}_\sigma\), the scheme-labelled completed
hierarchy is
\begin{equation}
\text{Total completed hierarchy}(\sigma_{\rm fin})
=
\alpha_\sigma^{\dim_{\RR}W}
=
\alpha_\sigma^7.
\end{equation}
The source determinant-shell member \(\mathrm{REG}_0\) is obtained by setting
\(\sigma_{\rm fin}=0\); it is not selected by the hierarchy theorem.
Paper~III-C's generated common-rate depth hits the first six ordered scalar
levels \(q=1,\ldots,6\), equivalently
\(\Ksh=8,16,\ldots,48\), at finite flow parameter and approaches the seventh
level \((q,\Ksh)=(7,56)\) only as its asymptotic endpoint.  If
\((\mathrm{FLAG})\) and a cutoff intertwiner realizing that flag are
additionally supplied, the same labels can represent seven successive
primitive projector losses.  Without those data, the scalar hitting
hierarchy, endpoint rank drop, exponent seven, and the active Clifford
landmarks \((q,\Ksh)=(0,0),(4,32),(5,40),(7,56)\) remain canonical, but no
rank-one projector chronology is asserted.
\end{theorem}

\begin{proof}
Theorem~\ref{p3:thm:hopf-decomposition} identifies the protected
ultraviolet-to-infrared endpoint rank drop as
\[
        \dim_{\RR}W=8-1=7.
\]
Theorem~\ref{p3:thm:reeb-universal} proves that each primitive rank-one
protected quotient carries the same topological logarithmic weight
\(\Vlogcap=4\pi^2\), equivalently the same multiplicative factor
\(\alpha_{\rm top}\), under the matched tensor extension.

By the matched tensor residue and the positive determinant character, a
rank-\(r\) protected non-scalar quotient has \(r\) times the primitive
additive logarithmic weight and therefore the \(r\)-fold multiplicative
factor.  Therefore the protected complement \(W\)
contributes
\[
        \alpha_{\rm top}^{\dim_{\RR}W}
        =
        \alpha_{\rm top}^{7}.
\]
The common-rate depth orders the scalar character values without resolving
the seven-dimensional eigenspace.  If \((\mathrm{FLAG})\) is supplied, this
product may additionally be read as seven ordered rank-one quotient factors.
An actual sharp spectral realization further
requires an intertwiner from its eigenspaces to that flag.  In its absence,
the exact Dirac staircase of Theorem~\ref{p3:thm:exact-rg} is a separate
spinorial count and cannot be used to order directions in \(W\).  Degenerate
losses may skip intermediate projector ranks while leaving the additive
endpoint valuation unchanged.

For a frozen \(\mathrm{REG}_\sigma\), the completed form follows by replacing
each primitive topological factor with the scheme-labelled
\(\alpha_\sigma\).  The source theorem
Theorem~\ref{p3:thm:D1-primitive-smooth-uniqueness} supplies the
\(\sigma_{\rm fin}=0\) member only.  The surviving scalar channel is the
infrared remnant described in Proposition~\ref{p3:prop:reeb-selection}.
\end{proof}

\begin{remark}[Generated common-rate readout]
The already-proved Paper~III-C common-rate semigroup is co-realized from the
protected split and certified by \((\mathrm{NEUT}_7)\).  It has one sevenfold-degenerate
non-scalar block.  This readout does not generate seven projector-rank events
and is not a hypothesis of Theorem~\ref{p3:thm:sequential} or an additional
III-D process occurrence.
\end{remark}

\begin{remark}[Why the terminal rank excludes the pseudoscalar]
\label{p3:rem:pseudoscalar-diagnostic-pointer}
An \(SO(3)\)-only average would also retain the orientation pseudoscalar.  The detailed obstruction is
recorded in Appendix~\ref{p3d:app:protected-diagnostics} of Paper~III-G.  The positive chain
uses the already constituted
\(\mathsf C_{\rm IR}^{\rm der}\) scalar remnant.  The excluded
\(SO(3)\)-only pseudoscalar branch would terminate at rank two and is not an
alternative endpoint of the active response.  The actual active obstruction
is stronger and differently typed: its rank-two invariant plane excludes the
scalar endpoint, so it cannot lie between \(P_0\) and \(\Pi_3\).
\end{remark}


\begin{mdframed}[nobreak=true]
\noindent\textbf{Constitution status of the reduced scalar response.}
The reduced scalar coupling \(g_{\eff}\), reciprocal-sector valuation, and
one-parameter scalar-trace response are not consequences of topology or Born
reciprocity alone.  They arise after the finite sector traces in the
spectral-action class \(\Dcl\) are passed through the Paper-III-B role-valued
map \(\mathfrak T_{\rm sec}\), and the resulting response class is restricted
by \((\mathrm{RED})\) and the abstract endpoint
calibration pair \((\mathrm{UV})+(\mathrm{IR})\).  The reduced object is a response curve, not autonomous time evolution.  Native III-D uses no
dimensionful reciprocal-radius calibration in this reduction.
\end{mdframed}

\begin{definition}[Reduced response branch \((\mathrm{RED})\)]
\label{p3:def:reduced-response-branch}
The reduced response branch is a declared response-class domain, not a
GP3-selected universal dynamics.  It extends the already derived
reciprocal-role domain
\[
 \Sigma_{\rm recip}^{\rm II-C/GP2}+\mathrm{G0}+\Dcl
\]
where \(\Sigma_{\rm recip}^{\rm II-C/GP2}\) denotes the premise-minimal
reciprocal closure (including the relation-free product and Born exchange),
by fixing, throughout the reduced leaf, the single common spectral profile
\(f\), its common cutoff/profile scale \(\Lcut\), and any common overall
normalization not absorbed into the scalar response conversion below.  The
branch also fixes the scalar phase/ground-channel projection, shellwise
\(\Spin(4)\)-equivariant coarse graining, and one Born product leaf.

After scalar projection, ordinary sector traces first feed the native finite
two-derivative response through the role-valued map
\(\mathfrak T_{\rm sec}\) of
Paper~III-B, Definition~\ref{p3b:def:sector-trace-born-state}.  Its two
scalar role outputs are
represented by positive typed weights
\[
 a>0,\qquad b>0,
\]
which Born exchange interchanges.  Before the product leaf is fixed they are
independent; on the leaf their product is fixed and their ratio supplies one
coordinate.  No continuum heat-kernel \(a_2\) interpretation is asserted for
these exact finite response weights.  Such an interpretation, when desired,
belongs to a separately stated smooth/semiclassical comparison.

Alternative process profiles, channel resolutions, coarse-graining packages,
or moving product leaves are datum-enlarged response classes rather than
hidden realizations rejected by GP3.

A later interpretation in terms of dimensionful reciprocal radii requires the
separate calibration \((\mathrm{CAL}_{xp})\); the dimensionful phase-cell
equality requires \((\mathrm{CAL}_{\hbar})\).  Neither calibration is used in
the native one-parameter theorem.  The normalized scalar phase generators of
Paper~III-B remain dimensionless comparison representatives and are not
assembled here into a dimensionful \(S^3\times S^1\) metric Dirac operator.
\end{definition}


\section{Born reciprocal roles, calibrated two-root response, and Laurent structure}
\label{paper3d:sec:organized-response}
\subsection{One-parameter reduction and reduced Born reciprocity}
\label{p3:sec:reduced_reciprocity}

The one-parameter reduction is the restricted branch \((\mathrm{RED})\) of
Definition~\ref{p3:def:reduced-response-branch}.  Unitary invariance and block
additivity follow from the closed loop; evenness comes from its generated
\(D^2\) process; valuation continuity is regularity of the occurring profile.
The conditioned-state normalization is theorem-level before
\((\mathrm{UNIT})\); \((\mathrm{UNIT})\) supplies only the physical
response-role adoption, while \((\mathrm{RED})\) supplies the reduced-leaf
data.  The typed
round sectors, scalar projection, and shellwise coarse graining are additional
conditions.  The scalar phase pair contributes no additional metric modulus
because only its normalized comparison generators are retained.

\begin{theorem}[One-parameter scalar-trace reduction on \((\mathrm{RED})\)]
\label{p3:prop:one_parameter_reduction}
\label{p3:thm:one-parameter-reduction}
For the completed package restricted by \((\mathrm{RED})\), every scalar
observable carried by the fixed-profile, fixed-cutoff role-valued sector-trace
output on a fixed
Born product leaf depends on the single Born-odd coordinate
\[
 u=\frac12\log(a/b),
\]
where \(a,b>0\) are the typed scalar role outputs.
Once \(f\), \(\Lcut\), the common normalization, and \(ab\) are fixed, no
transverse scalar role-output coordinate occurs in this reduced datum.  A larger class with an endomorphism, independent channel profile,
anisotropic deformation, moving product scale, additional metric modulus, or
unprojected scalar-circle datum may have additional coordinates.
\end{theorem}

\begin{proof}
\ProofPlan{We separate typed sector symmetry, scalar phase completion, the
channel-profile boundary, and the fixed Born product leaf.}
\begin{proofsteps}
\item \textbf{Typed sector symmetry leaves one scalar role output per reciprocal sector.}
The Paper-I-A/II kinematic package and Paper~III-A reciprocal-role constitution
fix the two typed round sectors.  Shellwise coarse graining is by complete
\(\Spin(4)\)-equivariant shells and preserves those sector types.  In the
scalar two-derivative part of the occurring closed action, the retained information is
therefore one positive scalar role output \(a\) in one sector and one output
\(b\) in the reciprocal sector.  Born exchange swaps them.

\item \textbf{The compact scalar phase pair contributes no new metric modulus.}
Paper~III-B's \(D_{\Sone_t}=-i\partial_\theta\) and
\(D_{\Sone_E}=-i\partial_\phi\) are normalized dimensionless phase-product
comparison generators.  Within \((\mathrm{RED})\), no scalar-circle metric modulus is retained; the
normalized phase generators are comparison representatives.  Introducing a
physical scalar-circle metric would enlarge the datum.  We do
not use the dimensionful equality \(\tper\Eper=2\pi\hbar\), and we do not
construct a dimensionful product Dirac operator from these normalized phase
representatives.

\item \textbf{No hidden channel-resolving coefficient occurs in \(\Dcl\).}
The sectorwise uninserted spectral action admits no natural non-scalar inserted
profile \(\Tr(Bf(D^2/\Lcut^2))\) without additional channel-resolving data.
A grade-resolved or anisotropic action is therefore a datum-enlarged class,
not another coordinate of \((\mathrm{RED})\).

\item \textbf{A fixed Born product leaf has one ratio.}
The leaf condition fixes \(ab=\scal_B^2\).  Hence
\[
 u=\frac12\log(a/b),
 \qquad
 a=\scal_B e^u,
 \qquad
 b=\scal_B e^{-u}.
\]
Every retained scalar-trace observable on that leaf is consequently a function
of this one coordinate.  No dimensionful reciprocal-radius interpretation is
needed for the reduction.
\end{proofsteps}
\end{proof}

\begin{remark}[Scope of the one-parameter branch]
\label{p3:rem:no-hidden-beta-system}
Theorem~\ref{p3:prop:one_parameter_reduction} is global only on a fixed Born
product leaf inside the loop-derived class \(\Dcl\).  Larger theories with
channel projectors, anisotropic deformations, a moving product scale, or
independent cutoff profiles are different datum classes; the detailed
exclusion ledger is deferred to Appendix~\ref{p3d:app:response-diagnostics} of Paper~III-G.
\end{remark}

\begin{definition}[Reciprocal quotient coordinate]

\label{p3:def:reciprocal-quotient-coordinate}
Let the unreduced Born pair be represented on a fixed invariant product leaf by two positive
reciprocal sector response weights
\[
        (a,b)\in\RR_{>0}^2,
        \qquad
        \mathcal B(a,b)=(b,a),
\]
where \(a\) and \(b\) are the two sectoral response weights before quotienting by
Born exchange.  The Born-even response midpoint and Born-odd quotient coordinate are
\[
        \scal_B:=\sqrt{ab},
        \qquad
        u:=\frac12\log\frac{a}{b}.
\]
Equivalently,
\[
        a=\scal_B e^u,
        \qquad
        b=\scal_B e^{-u}.
\]
After one-parameter reduction, the positive reduced response coordinate is the representative
\[
        \scal:=\scal_B e^u .
\]
\end{definition}

\begin{lemma}[Reduced Born reciprocity as an equivariant quotient]
\label{p3:lem:reduced-born-reciprocity}
Using the one-parameter reduction of Theorem~\ref{p3:prop:one_parameter_reduction}, on each fixed product leaf of the
unreduced reciprocal pair \((a,b)\), the quotient-coordinate constructor is
Born-unmarked.  Paper~II-A,
Theorem~\ref{p2a:thm:generated-born-inheritance}, therefore transports the
microscopic Born exchange to the quotient target; the induced action is exactly
\[
        u\longmapsto -u
\]
on the quotient coordinate of Definition~\ref{p3:def:reciprocal-quotient-coordinate}.  Therefore
on the positive reduced response representative it induces
\[
        \scal\longmapsto \frac{\scal_B^2}{\scal}.
\]
The fixed point on that reduced leaf is unique and is
\[
        u=0,
        \qquad
        a=b,
        \qquad
        \scal=\scal_B .
\]
\end{lemma}

\begin{proof}
The quotient map is not an additional dynamical assumption; it is the decomposition of the two
positive reciprocal response weights into their symmetric and antisymmetric logarithmic parts.  The inherited
Born-equivariance follows from Paper~II-A; the following calculation identifies
the target action explicitly.  Applying the microscopic Born exchange gives
\[
        u(\mathcal B(a,b))
        =\frac12\log\frac{b}{a}
        =-\frac12\log\frac{a}{b}
        =-u(a,b).
\]
The response midpoint is Born-even:
\[
        \sqrt{ba}=\sqrt{ab}=\scal_B .
\]
Hence, if \(\scal=\scal_B e^u\), then after Born exchange
\[
        \mathcal B\cdot\scal
        =\scal_B e^{-u}
        =\frac{\scal_B^2}{\scal_B e^u}
        =\frac{\scal_B^2}{\scal}.
\]
The fixed-point equation on the reduced leaf is \(u=-u\), hence \(u=0\).  Equivalently
\(a=b\), and therefore \(\scal=\scal_B\).  This proves both the descent and the uniqueness of the
fixed point on the reduced Born quotient line.
\end{proof}

\begin{theorem}[Born reflection of the seven-step response depth]
\label{p3d:thm:born-seven-depth-reflection}
On the owner-placed weight-one response branch for which
\[
 \frac{\scal(q)}{U}=\alpha_\sigma^{-q},
 \qquad 0\le q\le7,
 \qquad
 \frac{C}{U}=\alpha_\sigma^{-7},
 \qquad \alpha_\sigma>1,
\]
use the reduced Born reciprocity
$\scal\mapsto\scal_B^2/\scal$ with
$\scal_B^2=UC$.  Then
\[
 \boxed{B:q\longmapsto7-q},
 \qquad
 \boxed{q_B=\frac72},
 \qquad
 \boxed{\scal_B=U\alpha_\sigma^{-7/2}}.
\]
Thus the response endpoints $q=0$ and $q=7$ are reciprocal partners and the
response depth $7/2$ is the unique self-dual midpoint.  This is the response
reflection; it is not the protected-octet rank reflection $j\mapsto8-j$.
\end{theorem}

\begin{proof}
By Lemma~\ref{p3:lem:reduced-born-reciprocity},
\[
 B\scal(q)=\frac{UC}{\scal(q)}
 =\frac{U^2\alpha_\sigma^{-7}}{U\alpha_\sigma^{-q}}
 =U\alpha_\sigma^{-(7-q)}=\scal(7-q).
\]
The displayed response character is injective because $\alpha_\sigma>1$, so
the depth transforms as stated.  Its unique fixed point is $q=7/2$, and the
midpoint response follows by substitution.
\end{proof}

\begin{corollary}[Born inversion of the same-owner operational scale]
\label{p3d:cor:born-operational-scale-inversion}
Assume additionally the typed determinant-root ancestry of Paper~III-C on this
same response branch, so that
\[
 \frac{\mu(q)}{\mu_U}=\alpha_\sigma^{-q/4}.
\]
Define
\[
 \mu_B:=\mu_U\alpha_\sigma^{-7/8},
 \qquad
 \ell:=\log(\mu/\mu_B).
\]
Then
\[
 \boxed{B\mu=\frac{\mu_B^2}{\mu}},
 \qquad
 \boxed{B\ell=-\ell},
 \qquad
 \boxed{u=4\ell}
\]
for $u=\log(\scal/\scal_B)$.  A quantum cutoff $k$ inherits these formulas
only after a separately proved same-owner matching $k=c_\mu\mu$ with fixed
Born-even $c_\mu>0$.
\end{corollary}

\begin{proof}
The depth reflection gives
$\mu(7-q)=\mu_U\alpha_\sigma^{-(7-q)/4}$, while
$\mu_B^2/\mu(q)$ is the same expression.  The logarithmic parity follows.
Finally
$u=(7/2-q)\log\alpha_\sigma$ and
$\ell=(7/2-q)\log\alpha_\sigma/4$.
The cutoff statement is only substitution after the declared owner match.
\end{proof}


\begin{definition}[Visible reciprocal sector weights]
\label{p3:def:visible-reciprocal-sector-weights}
After one-parameter reduction and Born quotienting, let
\(\scal_{\rm resp}>0\) be an as-yet undetermined response-conversion scale.
The two microscopic reciprocal sectors then give the positive dimensionless
weights
\[
        X(\scal):=\frac{\scal}{\scal_{\rm resp}},
        \qquad
        Y(\scal):=\frac{\scal_B^2}{\scal_{\rm resp}\,\scal},
\]
The reduced Born involution exchanges them:
\[
        \scal\longmapsto \frac{\scal_B^2}{\scal}
        \quad\Longleftrightarrow\quad
        X\longleftrightarrow Y .
\]
Let \(Z_{\rm role}\), its idempotents \(e_x,e_p\), the role-valued trace
\(\mathfrak T_{\rm sec}\), and the normalized reciprocal role scalarizer \(\lambda_{\rm B}\)
be exactly those of Paper~III-B,
Definition~\ref{p3b:def:sector-trace-born-state}.  In particular, the
idempotents label scalar sector outputs and are not block identities in the
unequal finite matrix carrier.  Set
\begin{equation}
       W_{\rm scal}(\scal):=X(\scal)e_x+Y(\scal)e_p.
       \label{p3:eq:visible-reciprocal-central-element}
\end{equation}
A reduced two-derivative scalar coupling is called a
\emph{reciprocal-sector valuation} when it is the evaluation of this
role-valued element by the normalized reciprocal role scalarizer on \(Z_{\rm role}\).
Equivalently, on independent copies of the role-label algebra it obeys
\[
        G(X_1+X_2,Y_1+Y_2)=G(X_1,Y_1)+G(X_2,Y_2),
        \qquad X_i,Y_i\ge 0 .
\]
The map \(\mathfrak T_{\rm sec}\) first converts ordinary matrix-sector
traces into the scalar outputs assembled in \(W_{\rm scal}\); only then does
the role scalarizer act.  This is a statement about the role-valued finite
two-derivative response, not an assertion that every nonlinear function of
\(\scal\) is additive.
\end{definition}

\begin{lemma}[Normalized reciprocal role scalarizer induces the reciprocal-sector valuation]
\label{p3:lem:spectral-trace-sector-valuation}
In the loop-derived uninserted sectorwise spectral action of
Definition~\ref{p3:def:D1}, ordinary sector traces first produce the
role-valued output \(W_{\rm scal}\in Z_{\rm role}\).  The inherited normalized reciprocal role scalarizer on that role-label algebra is
\begin{equation}
       \lambda_{\rm B}(W_{\rm scal})=\frac{X+Y}{2}.
       \label{p3:eq:reciprocal-central-normalized-trace}
\end{equation}
It is therefore a reciprocal-sector valuation of the exact finite
role-valued response in the sense of
Definition~\ref{p3:def:visible-reciprocal-sector-weights}.
\end{lemma}

\begin{proof}
Paper~III-B's two-stage interface is
\[
 \mathfrak T_{\rm sec}(A_x\oplus A_p)
 =\Tr_{\mathcal H_x}(A_x)e_x
  +\Tr_{\mathcal H_p}(A_p)e_p,
 \qquad
 \mathcal S_{\rm B}=\lambda_{\rm B}\circ\mathfrak T_{\rm sec}.
\]
Thus the fixed scalar-sector outputs \(X,Y\) are assembled as
\(W_{\rm scal}=Xe_x+Ye_p\) before the role scalarizer acts.  A positive
normalized scalarizer on \(Z_{\rm role}\) has
\[
        \lambda(e_x)=a,\qquad \lambda(e_p)=b,
        \qquad a+b=1.
\]
Paper~III-B, Lemma~\ref{p3b:lem:equal-born-role-weights}, gives \(a=b\) from hereditary Born covariance of the licensed scalarization, while unit normalization gives \(a+b=1\); hence \(a=b=1/2\).  Substitution of
\eqref{p3:eq:visible-reciprocal-central-element} proves
\eqref{p3:eq:reciprocal-central-normalized-trace}.  Additivity is the
combination of ordinary sector-trace additivity in \(\mathfrak T_{\rm sec}\)
and linearity of \(\lambda_{\rm B}\) on \(Z_{\rm role}\), so no
Cauchy-extension argument in the response variable is required.

This normalization is not imposed on the unequal finite matrix blocks.  On
the current \(5\oplus15\) seed, the extensive composite instead gives
\[
 \mathcal S_{\rm B}(I_5\oplus0)=\frac52,
 \qquad
 \mathcal S_{\rm B}(0\oplus I_{15})=\frac{15}{2}.
\]
These values are compatible with the half weights of \(\lambda_{\rm B}\)
because the latter belong to the normalized role scalarizer, not to normalized
matrix traces or a prepared state.
\end{proof}
\begin{remark}[Extensive weights versus arbitrary nonlinear profiles]
Here \(X\) and \(Y\) are the visible scalar role outputs obtained after
scalar projection, direct-sum coarse-graining, and the sector traces encoded
by \(\mathfrak T_{\rm sec}\).  The valuation law applies to these extensive
outputs.  If one instead allowed a nonlinear function of the reduced response
coordinate as a new scalar profile, that profile would be an additional
datum-enlarged reduced structure beyond the uninserted direct-sum trace.
\end{remark}


\begin{remark}[General Cauchy/additivity lemma deferred]
Finite-dimensional role-label linearity, together with the sectorwise trace
map, already proves the linear reciprocal-sector valuation used here.  The more general continuous
Cauchy/additivity lemma is supplementary and is deferred to
Appendix~\ref{p3d:app:supplementary-valuations} of Paper~III-G.
\end{remark}

\subsection{Response normalization, endpoint calibrations, and two-root coupling}
\begin{definition}[Scalar response normalization representative]
\label{p3:def:scalar-response-readout-normalization}
For the reduced role-valued element
\(W_{\rm scal}=Xe_x+Ye_p\in Z_{\rm role}\), let
\[
 g_\kappa(\scal):=\kappa\,\lambda_{\rm B}(W_{\rm scal}),
 \qquad \kappa>0.
\]
Before endpoint calibration, only the ratio between the positive scalar
normalization \(\kappa\) and the response-conversion scale is observable in
this one-parameter formula.  Thus \(\kappa\) is a normalization
representative, not a generated physical selector.  We adopt the convenient
representative \(\kappa=2\) and write
\[
 \boxed{g_{\rm eff}(\scal):=g_2(\scal)
 =2\lambda_{\rm B}(W_{\rm scal})=X(\scal)+Y(\scal).}
\]
This convention does not renormalize the extensive matrix action and is not a
third primitive readout role.
\end{definition}

\begin{definition}[Upper and lower root calibrations]
\label{p3:def:uv-ir-calibration-pair}
Let \(0<C<U\) be the declared lower and upper response roots.  With
\(g_{\rm eff}\) already defined by
Definition~\ref{p3:def:scalar-response-readout-normalization}, the upper
calibration \((\mathrm{UV})\) requires \(g_{\rm eff}(U)=1\), while the lower
calibration \((\mathrm{IR})\) requires \(g_{\rm eff}(C)=1\).  The reciprocal role response retains an undetermined Born response
midpoint \(\scal_B>0\); the response-conversion scale
\(\scal_{\rm resp}>0\) has already been declared in
Definition~\ref{p3:def:visible-reciprocal-sector-weights}.
Theorem~\ref{p3:thm:geff_from_reciprocity} proves
\[
 \scal_B^2=CU,
 \qquad
 \scal_{\rm resp}=C+U.
\]
Neither root is identified natively with \(\Rp\), a hierarchy radius, or any
other dimensionful scale.
\end{definition}

\begin{theorem}[Exact two-root effective coupling on \(({\rm RED})\)]
\label{p3:thm:geff_from_reciprocity}
Let \(g_{\eff}(\scal)\) be the reduced dimensionless two-derivative scalar
coupling on \(({\rm RED})\).  The one-parameter reduction, the reduced Born involution, the normalized
role scalarizer, and the response readout normalization give
\[
 g_{\eff}(\scal)
 =\frac{\scal}{\scal_{\rm resp}}
  +\frac{\scal_B^2}{\scal_{\rm resp}\scal}.
\]
For declared roots \(0<C<U\), the calibration pair
\(({\rm UV})+({\rm IR})\),
\[
 g_{\eff}(U)=g_{\eff}(C)=1,
\]
uniquely forces
\[
 \boxed{\scal_B^2=CU,\qquad \scal_{\rm resp}=C+U}
\]
and hence
\[
\boxed{
 g_{\eff}(\scal)
 =\frac{\scal^2+CU}{(C+U)\scal}.
}
\]
\end{theorem}

\begin{proof}
Lemma~\ref{p3:lem:spectral-trace-sector-valuation} gives
\(\lambda_{\rm B}(W_{\rm scal})=(X+Y)/2\), and
Definition~\ref{p3:def:scalar-response-readout-normalization} therefore gives
\[
 g_{\eff}=X+Y
 =\frac{\scal}{\scal_{\rm resp}}
  +\frac{\scal_B^2}{\scal_{\rm resp}\scal}.
\]
The endpoint equations are
\[
 U+\frac{\scal_B^2}{U}=\scal_{\rm resp},
 \qquad
 C+\frac{\scal_B^2}{C}=\scal_{\rm resp}.
\]
Subtracting and using \(U\ne C\) gives
\[
 (U-C)\left(1-\frac{\scal_B^2}{UC}\right)=0,
 \qquad
 \scal_B^2=UC.
\]
Substitution gives \(\scal_{\rm resp}=U+C\), proving uniqueness and the
displayed response.  With this calibration established, its mathematical
consequences are those of Paper~I-D,
Theorem~\ref{p1b:thm:laurent-family}~\cite{Kugelmann2026Paper1D}, with
\(R=\scal\) and \(h=g_{\eff}\).
\end{proof}

\begin{corollary}[Primitive unit-quadratic two-root response]
\label{p3d:cor:primitive-unit-quadratic-response}
On \((\mathrm{RED})+(\mathrm{UV})+(\mathrm{IR})\), define the typed
response polynomial
\[
 \mathscr Q_{C,U}(\scal)
 :=(C+U)\scal\,g_{\rm eff}(\scal).
\]
Then
\[
 \boxed{\mathscr Q_{C,U}(\scal)=\scal^2+CU,\qquad
        \mathscr Q_{C,U}''(\scal)=2.}
\]
Thus the occurring two-root response carries a unit primitive quadratic
coefficient and one zero-order endpoint term, but no linear constituent.
This is a response-role theorem.  It does not by itself place
\(\mathscr Q_{C,U}\) in a Lorentzian gravitational action or identify its
owner with a determinant associated line.
\end{corollary}

\begin{proof}
Multiply the exact formula of
Theorem~\ref{p3:thm:geff_from_reciprocity} by
\((C+U)\scal\), then differentiate twice.  The final scope statement is the
existing response/action type boundary of \((\mathrm{RED})\).
\end{proof}


\begin{corollary}[Separated primitive and finite two-scale response normalizations]
\label{p3d:cor:two-scale-response-normalizations}
On the same calibrated branch define
\[
 \boxed{\mathscr H_{C,U}(\scal):=\scal\,g_{\rm eff}(\scal)
 =\frac{\scal^2+CU}{C+U}
 =\frac{\mathscr Q_{C,U}(\scal)}{C+U}.}
\]
Then
\[
 \boxed{\mathscr H_{C,U}''(\scal)=\frac{2}{C+U}},
 \qquad
 \frac{\mathscr H_{C,U}(\scal)}{\scal}
 =\frac1{C+U}\left(\scal+\frac{CU}{\scal}\right),
\]
and the last normalized ratio is invariant under the reduced Born inversion
\(\scal\mapsto CU/\scal\).
Thus \(\mathscr Q_{C,U}\) and \(\mathscr H_{C,U}\) are not competing
responses: \(\mathscr Q_{C,U}\) is the unit-quadratic primitive obtained by
multiplying by the response-conversion scale \(C+U\), whereas
\(\mathscr H_{C,U}\) is the finite UV/IR interpolation before that
multiplication.  A downstream Einstein--Hilbert action normalization must
therefore be applied to \(\mathscr H_{C,U}\), not silently identified with
the unit primitive \(\mathscr Q_{C,U}\).
\end{corollary}
\begin{proof}
Divide Corollary~\ref{p3d:cor:primitive-unit-quadratic-response} by
\(C+U\), differentiate twice, and use
Theorem~\ref{p3:thm:geff_from_reciprocity}.  For the reciprocal statement,
substitute \(CU/\scal\) in the displayed Laurent sum.
\end{proof}

\begin{corollary}[Full reciprocal primitive and its odd coefficient]
\label{p3d:cor:born-euler-response-pair}
On the same reduced and endpoint-calibrated branch, put \(B=CU\),
\(u_B=\log(\scal/\sqrt B)\), and
\(v_B=\mathrm d\log g_{\rm eff}/\mathrm d\log\scal\). Then
\[
 \mathscr A_B:=(C+U)g_{\rm eff}v_B
   =\scal-B/\scal,\qquad
 \mathrm d\mathscr A_B=(\mathscr Q_{C,U}/\scal)\,\mathrm du_B.
\]
In the scalar-function calculus of Paper~I-D,
Theorem~\ref{p1b:thm:born-euler-anticommutation},
\[
 \mathcal E\mathscr Q_{C,U}=\scal^2-B,\qquad
 \mathcal E(\scal^2-B)=\mathscr Q_{C,U},\qquad
 \mathcal E=\scal\partial_{\scal}-1.
\]
These are complete response identities, including the constant \(B\).
They neither identify \(\mathscr A_B\) with a metric-action coefficient nor
identify the response primitive with that action's Legendre potential.
\end{corollary}
\begin{proof}
Apply Paper~I-D, Theorem~\ref{p1b:thm:laurent-family} and
Theorem~\ref{p1b:thm:born-euler-anticommutation}, with \(R=\scal\).
The action comparison must retain the complete same-action zero-order
combination of Paper~III-F,
Proposition~\ref{p3f:prop:complete-action-legendre-potential}; a second
variation alone does not preserve this constant datum.
\end{proof}


\subsection{Born-invariant Laurent characterization and midpoint relations}
\begin{theorem}[Minimal Born-invariant Laurent characterization]
\label{p3d:thm:minimal-born-laurent-characterization}
Put
\[
 x:=\frac{\scal}{U},\qquad c:=\frac CU\in(0,1),
\]
and consider the real degree-one Laurent class
\[
 L(x)=a x+b x^{-1}+d.
\]
Require:
\begin{enumerate}[label=\textup{(\roman*)}]
\item invariance under the reduced Born exchange
      \(\mathcal B_c:x\mapsto c/x\);
\item the upper endpoint calibration \(L(1)=1\); and
\item no additional linear action constituent, meaning that multiplication
      by the curvature scale does not introduce a term proportional to
      \(\scal\) besides the separately owned Einstein--Hilbert line.
\end{enumerate}
Then the response is unique:
\begin{equation}
\label{p3d:eq:minimal-born-laurent}
 \boxed{
 L_c(x)=\frac{x+c/x}{1+c}
 =\frac{\scal^2+CU}{(C+U)\scal}
 =g_{\rm eff}(\scal).}
\end{equation}
Equivalently, the unnormalized minimal Laurent numerator is
\[
 \mathcal I_c(x):=x+\frac c x,
\]
and the factor \((1+c)^{-1}\) is fixed by endpoint calibration rather than
being an independent interpolation function.
\end{theorem}

\begin{proof}
Born invariance gives
\[
 a\frac c x+b\frac x c+d=ax+\frac b x+d
 \qquad(x>0),
\]
so coefficient comparison forces \(b=ac\).  Multiplication by \(\scal=Ux\)
maps the constant Laurent coefficient to \(d\scal\), which is precisely an
additional linear Einstein--Hilbert-type constituent; condition
\textup{(iii)} therefore gives \(d=0\).  The endpoint condition is then
\(a(1+c)=1\), hence \(a=(1+c)^{-1}\) and
\(b=c(1+c)^{-1}\).  This proves uniqueness and
Eq.~\eqref{p3d:eq:minimal-born-laurent}.  Multiplication by
\((1+c)\scal\) recovers the already derived primitive polynomial
\(\mathscr Q_{C,U}=\scal^2+CU\), so the characterization introduces no
new action term.
\end{proof}

\begin{corollary}[Raw rank-character Laurent secants and protected-octet reflection]
\label{p3d:cor:rank-character-laurent-family}
Suppose an owner-typed hierarchy placement realizes the rank-character ratios
\[
 c_j:=\alpha_\sigma^{-j},\qquad
 C_j:=Uc_j,\qquad j=1,\ldots,7.
\]
For every such realized descendant define
\begin{align}
\label{p3d:eq:rank-character-laurent-family}
 g_j(\scal)
 &:=\frac{\scal^2+UC_j}{(U+C_j)\scal}
   =\frac{x+c_j/x}{1+c_j},\\
 n_j^{\rm L}(\scal)
 &:=\frac{\ln g_j(\scal)}{\ln(\scal/U)}
 \qquad(\scal\ne U),
\end{align}
with the removable value
\[
 n_j^{\rm L}(U)=\frac{1-c_j}{1+c_j}.
\]
Equivalently, with
\[
 A_\sigma:=\ln\alpha_\sigma>0,
 \qquad
 q:=-\frac{\ln(\scal/U)}{A_\sigma},
 \qquad K:=8q,
\]
the same Laurent interpolation is
\begin{equation}
\label{p3d:eq:rank-character-log-clock-form}
 \boxed{
 n_j^{\rm L}(q)
 =\frac{\ln\cosh(jA_\sigma/2)
       -\ln\cosh[(j/2-q)A_\sigma]}
       {qA_\sigma},}
\end{equation}
with its removable value at \(q=0\).  Thus the Born fixed point is
\(q=j/2\), while the exact nonnegative endpoint is \(q=j\), or \(K=8j\).
Then
\[
 g_j(\scal)=\left(\frac{\scal}{U}\right)^{n_j^{\rm L}(\scal)},
 \qquad
 n_j^{\rm L}(C_j)=0,
 \qquad
 \scal_{B,j}=\sqrt{UC_j}=U\alpha_\sigma^{-j/2}.
\]
The raw logarithmic secant has the exact sign table
\begin{equation}
\label{p3d:eq:raw-laurent-index-sign}
 n_j^{\rm L}(\scal)
 \begin{cases}
 >0,&\scal>C_j,\\
 =0,&\scal=C_j,\\
 <0,&0<\scal<C_j,
 \end{cases}
 \qquad
 \lim_{\scal\downarrow0}n_j^{\rm L}(\scal)=-1.
\end{equation}
Thus \([C_j,\infty)\) is the maximal nonnegative domain of the raw secant,
but the formula by itself does not supply a frozen continuation.
Independently, the protected-octet trace reflection on
\(\mathbb R\mathbf1\oplus W_7\) has rank involution
\(j\mapsto8-j\) and therefore the endpoint pair
\[
 \boxed{j_{\rm first}=1,\qquad j_{\rm term}=7}
\]
and
\begin{equation}
\label{p3d:eq:one-seven-complement-product}
 C_1C_7
 =U^2\alpha_\sigma^{-8}
 =\bigl(U\alpha_\sigma^{-4}\bigr)^2.
\end{equation}
This is a rank and endpoint symmetry.  It neither selects the coefficient
dressing used in a metric action nor asserts the false pointwise relation
\(n_1^{\rm L}+n_7^{\rm L}=1\).
\end{corollary}

\begin{proof}
Theorem~\ref{p3d:thm:minimal-born-laurent-characterization} identifies
the response at \(C=C_j\).  Paper~I-D,
Theorem~\ref{p1b:thm:laurent-family}~\cite{Kugelmann2026Paper1D}, then
gives the logarithmic-clock form, removable value, Born fixed point,
strict sign table and maximal nonnegative domain, with
\(R=\scal\), \(A=A_\sigma\) and \(n_j=n_j^{\rm L}\).
The limit at zero follows from
\(g_j(\scal)\sim UC_j/[(U+C_j)\scal]\).
Paper~III-C's protected decomposition has ranks one and seven.  Under
their common octet owner, complementary ranks sum to eight, so the
nontrivial pair containing the terminal rank-seven constituent is
\((1,7)\).  Multiplicativity of the positive rank character gives
Eq.~\eqref{p3d:eq:one-seven-complement-product}.  This rank symmetry
does not imply a functional sum rule or assign the action dressing.
\end{proof}

\begin{theorem}[Calibration-free midpoint and endpoint-slope consistency of rank-character responses]
\label{p3d:thm:v285-calibration-free-hierarchy-dip}
On a single owner-placed positive-character rank-response branch of
Corollary~\ref{p3d:cor:rank-character-laurent-family}, let
\(L_\sigma=\log\alpha_\sigma>0\), \(j\in\{1,\ldots,7\}\),
\(c_j=e^{-jL_\sigma}\in(0,1)\), and
\[
 x=\frac{\scal}{U}>0,\qquad
 g_j(x)=\frac{x+c_j/x}{1+c_j}.
\]
The unique global minimum, the two logarithmic endpoint slopes, and their
normalization-free identity are
\begin{align}
\label{p3d:eq:v285-dip-and-slopes}
 x_{\min,j}&=\sqrt{c_j},&
 m_j:=g_j(x_{\min,j})
 &=\operatorname{sech}\!\left(\frac{jL_\sigma}{2}\right),\\
 v_j^+&:=\left.\frac{d\log g_j}{d\log\scal}\right|_{\scal=U}
 =\tanh\!\left(\frac{jL_\sigma}{2}\right),&
 v_j^-&:=\left.\frac{d\log g_j}{d\log\scal}\right|_{\scal=C_j}
 =-v_j^+,\\
 \boxed{m_j^2+(v_j^+)^2}&=1,&
 \boxed{L_\sigma}&=\frac2j\operatorname{arcosh}(m_j^{-1}).
\end{align}
Consequently, if two distinct levels \(j,k\) actually occur in the
\emph{same} determinant/response owner, their dimensionless minima satisfy
\[
 \boxed{\frac1j\operatorname{arcosh}(m_j^{-1})
       =\frac1k\operatorname{arcosh}(m_k^{-1})},
 \qquad
 m_k=\operatorname{sech}\!\left(\frac{k}{j}
       \operatorname{arcosh}(m_j^{-1})\right).
\]
These are consistency constraints, not new calibration parameters.  They
require neither an absolute value of \(U\) nor an SI scale.  Their
\emph{physical} use requires actual response readouts at the indicated
levels, plus a generated common-owner ancestry identification; neither
those occurrences nor a common temporal flow is supplied by this theorem.
\end{theorem}
\begin{proof}
The exact Laurent family is Corollary~\ref{p3d:cor:rank-character-laurent-family}.
Differentiate with respect to \(x\):
\[
 g_j'(x)=\frac{1-c_j/x^2}{1+c_j},\qquad
 g_j''(x)=\frac{2c_j}{(1+c_j)x^3}>0.
\]
The unique critical point is \(x=\sqrt{c_j}\), where
\(m_j=2\sqrt{c_j}/(1+c_j)\).  Substituting
\(c_j=e^{-jL_\sigma}\) gives the hyperbolic-secant expression.
The logarithmic derivative is
\[
 \frac{d\log g_j}{d\log\scal}
 =\frac{x^2-c_j}{x^2+c_j}.
\]
Evaluating at \(x=1\) and \(x=c_j\), respectively, yields
\(v_j^\pm=\pm(1-c_j)/(1+c_j)
 =\pm\tanh(jL_\sigma/2)\).
The identity \(\operatorname{sech}^2 z+\tanh^2 z=1\) proves the
normalization-free constraint.  Since \(L_\sigma>0\), the inverse of
\(\cosh\) on \([0,\infty)\) yields the unique reconstruction of
\(L_\sigma\).  Equating the recovered logarithm for two levels proves
the cross-level relation.  All steps use only the declared rank-character
response, not a physical owner-placement theorem.
\end{proof}

\begin{corollary}[Born-midpoint curvature-shell coordinate]
\label{p3d:cor:born-midpoint-shell-coordinate}
On the realized rank-character leaf of
Corollary~\ref{p3d:cor:rank-character-laurent-family}, let
\[
 \mathcal B_j(\scal):=\frac{UC_j}{\scal},
 \qquad
 K_R(\scal):=16\log_{\alpha_\sigma}\!\frac{U}{\scal}.
\]
Then
\begin{equation}
\label{p3d:eq:born-midpoint-shell-reflection}
 \boxed{
 K_R\!\left(\mathcal B_j(\scal)\right)=16j-K_R(\scal),
 \qquad
 K_R(\scal_{B,j})=8j,}
\end{equation}
where
\[
 \scal_{B,j}=\sqrt{UC_j}=U\alpha_\sigma^{-j/2}
\]
is the unique Born fixed point on that leaf.  In particular the first realized
cell has
\[
 \boxed{\scal_{B,1}=U\alpha_\sigma^{-1/2},\qquad K_R(\scal_{B,1})=8.}
\]
The corresponding normalized Born-even slack is
\begin{equation}
\label{p3d:eq:born-midpoint-shell-slack}
 \boxed{
 \mathfrak S_j(\scal)
 :=\frac12\left(\frac{\scal}{\scal_{B,j}}
 +\frac{\scal_{B,j}}{\scal}\right)-1
 =\cosh\!\left[\frac{A_\sigma}{16}\bigl(K_R-8j\bigr)\right]-1\ge0,}
\end{equation}
with equality only at the Born midpoint.

The symbol \(K_R\) is a curvature-shell coordinate on the realized positive
response variable.  It is not Paper~III-C's structural common-rate clock
\(\Ksh=8q_{\rm sp}\), and no cosmological epoch name is assigned here.
\end{corollary}

\begin{proof}
Using \(C_j=U\alpha_\sigma^{-j}\),
\[
 \frac{U}{\mathcal B_j(\scal)}
 =\frac{\scal}{C_j}
 =\alpha_\sigma^j\frac{\scal}{U}.
\]
Taking \(16\log_{\alpha_\sigma}\) gives the reflection law.  Evaluating at
\(\scal_{B,j}=U\alpha_\sigma^{-j/2}\) gives \(K_R=8j\).  Finally write
\(u=\log(\scal/\scal_{B,j})=-(A_\sigma/16)(K_R-8j)\); the displayed slack is
\(\cosh u-1\), so strict positivity away from the fixed point follows from
strict convexity of \(\cosh\).
\end{proof}

\subsection{Positive-part Laurent junction and finite carrier dressing}
\begin{theorem}[Canonical positive-part Laurent projection and endpoint trace]
\label{p3d:thm:canonical-nonnegative-laurent-projection}
For a real scalar \(y\), let
\begin{equation}
\label{p3d:eq:positive-part-projector}
 \boxed{
 \Pi_+(y):=\frac{y+\sqrt{y^2}}{2}=\max\{y,0\}.}
\end{equation}
It is the unique Euclidean metric projection of the real line onto the closed
convex cone \(\mathbb R_{\ge0}\):
\[
 \Pi_+(y)=\underset{z\ge0}{\operatorname{argmin}}\,(z-y)^2.
\]
For every realized Laurent rank descendant define
\begin{equation}
\label{p3d:eq:laurent-positive-part-projection}
 \boxed{
 n_j^+(\scal)
 :=\Pi_+\!\left(n_j^{\rm L}(\scal)\right)
 =\frac{n_j^{\rm L}(\scal)
 +\sqrt{(n_j^{\rm L}(\scal))^2}}{2}.}
\end{equation}
Then
\begin{equation}
\label{p3d:eq:nonnegative-laurent-index-properties}
 0\le n_j^+(\scal)\le1,\qquad
 n_j^+(\scal)=0\ \Longleftrightarrow\ 0<\scal\le C_j,
\end{equation}
while \(n_j^+=n_j^{\rm L}\) on \([C_j,\infty)\).  Thus the projection is
exactly zero, rather than the reflected positive value
\(\lvert n_j^{\rm L}\rvert\), wherever the raw Laurent index is negative.
It is continuous but, for the simple zero at \(C_j\), not \(C^1\).

On the active Laurent domain \(\scal\ge C_j\), the nonlinear response remains
\begin{equation}
\label{p3d:eq:active-laurent-response-and-trace}
 \boxed{
 \mathcal H_j^{\rm act}(\scal)
 :=\scal\left(\frac{\scal}{U}\right)^{n_j^+(\scal)}
 =\frac{\scal^2+UC_j}{U+C_j},\qquad \scal\ge C_j,}
 \qquad
 \boxed{\operatorname{ev}_{C_j}\mathcal H_j^{\rm act}=C_j.}
\end{equation}
The last expression is a boundary trace.  It does not define
\(\scal(\scal/U)^0=\scal\) as a freely varied nonlinear action constituent
for \(0<\scal<C_j\).  For the sole purpose of defining a distributional junction on the positive
real axis, assemble the \emph{continuous glued representative}
\begin{equation}
\label{p3d:eq:laurent-trace-glued-representative}
 \boxed{\mathcal H_j^{\rm junc}(\scal):=
 \begin{cases}
 C_j,&0<\scal\le C_j,\\
 (\scal^2+UC_j)/(U+C_j),&\scal>C_j.
 \end{cases}}
\end{equation}
Its two traces agree at \(C_j\), while its first derivative has jump
\(2C_j/(U+C_j)\). On \((0,\infty)\), in the sense of distributions,
\begin{equation}
\label{p3d:eq:endpoint-trace-positive-hessian-measure}
 \boxed{d\bigl((\mathcal H_j^{\rm junc})'\bigr)
 =\frac{2}{U+C_j}\mathbf1_{(C_j,\infty)}\,d\scal
  +\frac{2C_j}{U+C_j}\,\delta_{C_j}.}
\end{equation}
This is a nonnegative measure. The terminal trace \(C_j\) by itself is
constant and has zero distributional Hessian. The glued representative is
\emph{bookkeeping}, not a new physical off-shell action or a proof of an
actual crossing history.
\end{theorem}

\begin{proof}
If \(y\ge0\), the unconstrained minimizer \(z=y\) is admissible.  If
\(y<0\), the strictly convex objective is increasing on \([0,\infty)\), so
its unique constrained minimizer is \(z=0\).  Since
\(\sqrt{y^2}=\lvert y\rvert\) for real \(y\), this proves
Eq.~\eqref{p3d:eq:positive-part-projector}.  Apply Paper~I-D,
Theorem~\ref{p1b:thm:laurent-family}~\cite{Kugelmann2026Paper1D}, with
\(R=\scal\) and \(C=C_j\): its exact sign and range give
Eqs.~\eqref{p3d:eq:laurent-positive-part-projection}
and~\eqref{p3d:eq:nonnegative-laurent-index-properties}.
On every \(\scal\ge C_j\), \(n_j^+=n_j^{\rm L}\), and the logarithmic
definition gives \(\scal(\scal/U)^{n_j^{\rm L}}=\scal g_j(\scal)\)
for \(\scal\ne U\), with equality at \(U\) by continuity.
Thus Eq.~\eqref{p3d:eq:active-laurent-response-and-trace} holds also for
\(\scal>U\), beyond the displayed interval of the 1B active identity.
Evaluation at \(\scal=C_j\) gives \(C_j\).  The constant left branch of
Eq.~\eqref{p3d:eq:laurent-trace-glued-representative} has derivative zero;
the active right branch has derivative \(2\scal/(U+C_j)\).
The branches agree in value at \(C_j\), so the distributional first
derivative has no delta term. Its ordinary traces are \(0\) and
\(2C_j/(U+C_j)\), respectively. Differentiating this piecewise
first derivative against a compactly supported smooth test function on
\((0,\infty)\) yields the regular density and the positive Dirac
mass of Eq.~\eqref{p3d:eq:endpoint-trace-positive-hessian-measure}.
The distribution does \emph{not} arise from differentiating the constant
terminal trace alone.
\end{proof}

\begin{theorem}[Least-support Laurent dressing index]
\label{p3d:thm:least-support-n2-dressing}
Fix a decreasing-curvature clock
\(q=-\ln(\scal/U)/\ln\alpha_\sigma\) and a completed inflationary end
\(q_{\rm end}\) satisfying
\[
 \frac12<q_{\rm end}<1.
\]
Within the already realized positive-part Laurent family, impose only the
following dressing-role requirements:
\begin{enumerate}[label=\textup{(D\arabic*)}]
\item the Born turnover center \(q=j/2\) lies strictly after inflation,
      \(j/2>q_{\rm end}\);
\item the dressing is exactly absent no later than the protected midpoint
      \(K=32\), equivalently its endpoint \(q=j\) obeys \(j\le4\); and
\item among complete candidates satisfying \textup{(D1)--(D2)}, GP3 retains
      the unique least active support.
\end{enumerate}
Then
\begin{equation}
\label{p3d:eq:least-support-n2-selection}
 \boxed{j_{\rm dress}=2,\qquad
        C_{\rm dress}=C_2=U\alpha_\sigma^{-2},\qquad
        \{R\in(0,U]:n_2^+(R)\ne0\}=(C_2,U].}
\end{equation}
Its Born center is \(q=1\), or \(K=8\), and its exact freeze is \(q=2\),
or \(K=16\).  The protected-octet rank reflection still gives the separate
\(1\leftrightarrow7\) endpoint symmetry; it is not the selector used here.
Likewise \(C_1^2/U=C_2\) is a corroborating equality of hierarchy values, not
by itself a theorem identifying tensor grade with response depth.
\end{theorem}

\begin{proof}
Because \(q_{\rm end}>1/2\), condition \textup{(D1)} excludes \(j=1\) and
forces \(j\ge2\).  Condition \textup{(D2)} forces \(j\le4\), so the complete
candidate set is exactly \(\{2,3,4\}\).  Since
\(C_4<C_3<C_2<U\), their nontrivial active supports are strictly nested:
\[
 (C_2,U]\subsetneq(C_3,U]\subsetneq(C_4,U].
\]
Therefore the unique least support is the \(j=2\) member.  The center and
freeze follow from Eq.~\eqref{p3d:eq:rank-character-log-clock-form} and the
positive-part zero set.  Neither the octet reflection nor the value identity
\(C_1^2/U=C_2\) enters the minimization, proving the scope statements.
\end{proof}

\begin{remark}[Occurrence boundary of the \(n_2\) selector]
Theorem~\ref{p3d:thm:least-support-n2-dressing} compares realizations only
after the coefficient-dressing role and its two support requirements have
been stated.  It does not make a coefficient dressing occur in a Lorentzian
metric action.  That action-role occurrence remains an explicit Paper--III-E
premise; its least-support tie-break applies only to its own
optional comparison rule and is distinguished from the active-carrier rule of
Theorem~\ref{p3d:thm:v2219-carrier-four-dressing}.
\end{remark}


\begin{theorem}[Carrier-constrained ultraviolet dressing selection]
\label{p3d:thm:v2219-carrier-four-dressing}
Assume the dressing index is in the previously certified comparison set
$j\in\{2,3,4\}$, but replace the least-active-support condition (D3) by
the \emph{first actual carrier rule}: the coefficient source is represented
by the first nonzero spectral grade of the same protected finite carrier,
$\mathcal Q_{\rm act}$. On the existing positive PQTI coefficient
line take its canonical positive character
$\rho_Q(s):=\exp(s\log Q_\sigma)$, $Q_\sigma=8\sqrt{\alpha_\sigma}>1$,
and evaluate it on the selected Laurent exponent (without asserting a
metric-action insertion). Then the unique selected rank and its dependent
mathematical coefficient pullback are
\[
 \boxed{j_{\rm dress}=4,\qquad P_2=P_3=0,\quad P_4\ne0,
 \qquad D_4(R)=Q_\sigma^{n_4^+(R)}.}
\]
Its reciprocal Born involution and fixed point are
\[
 \mathcal B_4(R)=\frac{UC_4}{R},\qquad
 R_{B,4}=\sqrt{UC_4}=U\alpha_\sigma^{-2}=C_2.
\]
On the dimension-weighted physical scale dictionary,
$C_2=E_{32}^{\,2}$, and rank four is also the unique protected-octet
reflection fixed grade $J(4)=4$.  The Laurent midpoint is a balance
condition, \emph{not} the scalaron frequency equality $m_s=E_{32}$.
\end{theorem}
\begin{proof}
The first actual active grade above zero is four by
Proposition~\ref{p3c:prop:v2219-first-actual-carrier}. Intersecting
\(\{2,3,4\}\) with \(\operatorname{spec}\mathcal Q_{\rm act}\setminus\{0\}\)
gives exactly \(\{4\}\). Since \(Q_\sigma>0\), positive functional calculus constructs
\(\rho_Q(s)=\exp(s\log Q_\sigma)=Q_\sigma^s\) on the generated
one-dimensional coefficient line. It is multiplicative and continuous,
with \(\rho_Q(1)=Q_\sigma\); any other such character has continuous
additive logarithm and thus coincides with this one.
Its scalar evaluation at the already defined \(n_4^+(R)\) therefore
yields \(D_4(R)=\rho_Q(n_4^+(R))\).
This uses the same character mathematics later formalized in Paper~III-F,
but does not import the later depth-two pullback as an upstream premise.
This is an exact character pullback, not a proof that \(P_4\) has
entered the physical gravitational action. The unique positive fixed point
of the reciprocal involution solves \(R^2=UC_4\). The native energy
identity and \(J(4)=4\) follow from the independent hierarchy and octet
involution.
\end{proof}
\begin{remark}[Scope of the optional least-support selection]
Theorem~\ref{p3d:thm:least-support-n2-dressing} still proves $j=2$
\emph{if} its optional least-support criterion (D3) is imposed.  It no
longer selects the active physical UV coefficient: a formal Laurent index
with no actual carrier cannot satisfy the first-actual-carrier selection
premise.  The incidence of $P_4$ into the gravitational coefficient remains
a common-action coupling certificate to be checked independently.
\end{remark}


\begin{remark}[Normalization-representative invariance]
\label{p3d:rem:response-normalization-invariance}
If one keeps a general \(\kappa>0\) in
Definition~\ref{p3:def:scalar-response-readout-normalization}, then before
choosing the representative \(\kappa=2\) the response is
\[
 g_\kappa(\scal)
 =\frac{\kappa}{2\scal_{\rm resp}}
  \left(\scal+\frac{\scal_B^2}{\scal}\right).
\]
The two endpoint equations imply
\[
 \scal_B^2=CU,
 \qquad
 \scal_{\rm resp}^{(\kappa)}=\frac\kappa2(C+U),
\]
and therefore every positive normalization representative gives the same
calibrated invariant curve
\[
 g_\kappa(\scal)
 =\frac{\scal^2+CU}{(C+U)\scal}.
\]
Thus the factor \(2\) is a convenient convention, not a surviving physical
selector.  The notation \(\scal_{\rm resp}=C+U\) used in the main text is the
\(\kappa=2\) representative.
\end{remark}


\begin{remark}[Logical status of the reduced coupling]
\label{p3:rem:geff_status}
Within \(({\rm RED})\), the reduced involution
\(\scal\mapsto \scal_B^2/\scal\) and the reciprocal-sector valuation give
the form \(X+Y\) after normalized role-label conversion.  The calibration
pair fixes \(\scal_B^2=CU\) and \(\scal_{\rm resp}=C+U\).  A nonlinear
self-dual response is a different reduced action datum, not something
excluded by GP1--GP3.
\end{remark}

\subsection{Normalized Born-symmetric curves and endpoint factorization}
\begin{corollary}[Exact Born-symmetric one-parameter response curve]
\label{p3:cor:fully_derived_one_parameter_trajectory}
On each reduced Born quotient leaf define
\[
        u:=\ln(\scal/\scal_B).
\]
Then the reduced response curve is the single global line
\[
        g_{\eff}(u)=\frac{2\scal_B}{C+U}\cosh u,
\]
with Born exchange acting by \(u\mapsto -u\).  The instantaneous logarithmic index is
\[
        \upsilon(u)=\frac{d\ln g_{\eff}}{du}=\tanh u,
\]
so
\[
        \frac{d\upsilon}{d\ln\scal}=1-\upsilon^2.
\]
The \(\tanh\) and Riccati formulas are algebraic identities of this curve.
They do not supply an autonomous time evolution or beta function without an
additional law specifying how the curve is traversed.
\end{corollary}

\begin{proof}
Theorem~\ref{p3:thm:geff_from_reciprocity} gives
\[
 g_{\eff}(\scal)
 =\frac{\scal}{C+U}+\frac{\scal_B^2}{(C+U)\scal}.
\]
Substitute
\[
  \scal=\scal_B e^u .
\]
This gives the hyperbolic-cosine form.  Born exchange is
Lemma~\ref{p3:lem:reduced-born-reciprocity}, hence \(u\mapsto -u\).

Differentiating
\[
  \ln(2\scal_B\cosh u/(C+U))
\]
gives \(\upsilon=\tanh u\).  Since \(u=\ln(\scal/\scal_B)\), differentiation
with respect to \(\ln\scal\) gives
\[
  \frac{d\upsilon}{d\ln\scal}=1-\upsilon^2.
\]
The only coordinate in the declared \((\mathrm{RED})\) fixed-leaf class is
\(u\).  A second coordinate requires datum that enlarges this class, for
example channel resolution, a moving product scale, anisotropy, or another
metric/profile modulus.
\end{proof}

\begin{remark}[Immediate consequence for the running index]
\label{p3:rem:riccati_from_geff}
Corollary~\ref{p3:cor:fully_derived_one_parameter_trajectory} immediately yields the Riccati identity for the instantaneous index later recorded in Theorem~\ref{p3:thm:riccati}.
\end{remark}


\begin{remark}[Dimensionless hierarchy handoff]
\label{p3d:rem:dimensionless-hierarchy-handoff}
The native hierarchy data imported from Paper~III-C are the dimensionless
positive character and its neutral dual scalar-depth ratios
\[
 H_{56}^{(\sigma)}=\alpha_\sigma^{7},
 \qquad
 \frac{X_w(q)}{X_w(0)}=\alpha_\sigma^{-wq}.
\]
The second formula denotes the neutral character and is generated without
\((\mathrm{HT})\).  It does not identify the abstract response-root owner
with the determinant-dual owner, or \(U\) with \(\Rp\) or another
dimensionful standard.  Once a response composite contains a generated
nonzero ancestry morphism from the upper-root owner to the determinant-dual
owner, however, Paper~III-C's owner-inheritance theorem fixes its lower root
and intermediate ratios without a further normalization.  The residual
typed placement and absolute calibration are
deferred to Appendix~\ref{p3d:app:hierarchy-placement} of Paper~III-G.
\end{remark}

\begin{remark}[Anchored hierarchy and its physical uses]
\label{p3d:rem:anchored-reg-hierarchy-import}
On the occurring Hurwitz/Hopf branch,
Theorem~\ref{p3c:thm:occurring-hurwitz-reg-origin} fixes
\(\mathrm{REG}_0\) and \(\alpha_0\).  The generic
rank-character Laurent formulas retain their unanchored
\(\alpha_\sigma\) version.  An absolute scale, physical
running or pole map, and an actual history are separate
inputs to claims consuming them.  In particular a
hierarchy-depth difference is not a physical clock duration.
\end{remark}


\begin{theorem}[Exact two-root endpoint factorization]
\label{p3:thm:exact-reciprocal-freeze-interface}
For the native response roots \(0<C<U\),
\[
 g_{\eff}(\scal)=\frac{\scal^2+CU}{(C+U)\scal}
\]
and
\begin{equation}
\label{p3:eq:exact-freeze-equation}
 \boxed{g_{\eff}(\scal)-1
 =\frac{(\scal-C)(\scal-U)}{(C+U)\scal}.}
\end{equation}
Thus \(g_{\eff}(C)=g_{\eff}(U)=1\) and the Born fixed point is
\(\scal_B=\sqrt{CU}\).  The secant value at the lower root is established
only after the secant index is defined in the next section.  These are native
two-root identities and require no hierarchy-scale placement.
\end{theorem}

\begin{proof}
Subtract one from Theorem~\ref{p3:thm:geff_from_reciprocity} and factor the
numerator.  The roots and fixed point follow immediately.
\end{proof}

\begin{remark}[Forward transport normalization]
\label{p3:rem:forward-transport-normalization}
Any downstream use of the exact response must retain
\[
 \scal_{\rm resp}\,\scal\,g_{\eff}(\scal)=\scal^2+CU,
 \qquad \scal_{\rm resp}=C+U.
\]
Replacing \(C+U\) by either endpoint separately is not exact.
\end{remark}

\newpage


\subsection{Born conjugation of the Hopf--Clifford scale character}
\label{p3d:sec:v222-phase-born-duality}

The unitary companion of the positive rank character is constructed in
Paper~III-C, Theorem~\ref{p3c:thm:v222-rank-character-unitary-companion}.
The statement below is entirely on the \emph{reduced hierarchy scale}
$\scal$; it is not an assertion that a spacetime Ricci scalar and
$\scal$ are already identified in a given gravitational action.

\begin{theorem}[Exact Born conjugation and rank-seven antiphase]
\label{p3d:thm:v222-born-phase-conjugation}
On the owner-placed response branch of
Theorem~\ref{p3d:thm:born-seven-depth-reflection}, write
\[
 L_\sigma=\log\alpha_\sigma>0,\qquad
 C_\sigma=Ue^{-7L_\sigma},\qquad
 \chi_\sigma(\scal)
 =\exp\!\left(\frac{2\pi i}{L_\sigma}\log\frac{\scal}{U}\right).
\]
For the reduced Born involution $\mathcal B\scal=UC_\sigma/\scal$,
\begin{equation}
\label{p3d:eq:v222-born-conjugation}
 \boxed{\chi_\sigma(\mathcal B\scal)
       =\overline{\chi_\sigma(\scal)}}.
\end{equation}
Thus $F_\sigma=\operatorname{Re}\chi_\sigma$ is Born-even and
$S_\sigma=\operatorname{Im}\chi_\sigma$ is Born-odd.
At the unique Born-fixed scale $\scal_B=\sqrt{UC_\sigma}$,
\[
 \boxed{\chi_\sigma(\scal_B)=-1,\qquad
        F_\sigma(\scal_B)=-1,\qquad S_\sigma(\scal_B)=0.}
\]
All completed rank depths $j\in\{0,\ldots,7\}$ are phase returns:
$\chi_\sigma(U\alpha_\sigma^{-j})=1$.
\end{theorem}

\begin{proof}
Write $\scal=U\alpha_\sigma^{-q}$.
The inherited reduced Born involution is $q\mapsto7-q$.
Consequently
\[
 \chi_\sigma(\mathcal B\scal)
 =e^{-2\pi i(7-q)}=e^{2\pi iq}
 =\overline{e^{-2\pi iq}}
 =\overline{\chi_\sigma(\scal)}.
\]
Taking real and imaginary parts proves the stated parity.  The fixed
point $q=7/2$ has character $e^{-7\pi i}=-1$; at any integer $j$ the
value is $e^{-2\pi i j}=1$.
\end{proof}

\begin{theorem}[Primitive contact phase and the completed Born twist]
\label{p3d:thm:v222-primitive-phase-twist}
Assume $(\mathrm{REG})$, so $L_\sigma=L_{\rm top}+\delta_{\rm loc}$
with $L_{\rm top}=4\pi^2$.  In addition to the completed character
above, retain the \emph{undressed} primitive character
\[
 \chi_{\rm top}(\scal)
   :=\exp\!\left(\frac{2\pi i}{L_{\rm top}}
                         \log\frac{\scal}{U}\right).
\]
At the stabilized rank depths
$\scal_j=Ue^{-j(L_{\rm top}+\delta_{\rm loc})}$ one has
\begin{equation}
\label{p3d:eq:v222-primitive-phase-drift}
 \boxed{\chi_{\rm top}(\scal_j)
       =e^{-i j\delta_{\rm loc}/(2\pi)}}.
\end{equation}
For the \emph{completed} Born reflection
$\mathcal B\scal=UC_\sigma/\scal$, the primitive character obeys the
fixed-twist law
\begin{equation}
\label{p3d:eq:v222-primitive-born-twist}
 \boxed{\chi_{\rm top}(\mathcal B\scal)
    =e^{-i7\delta_{\rm loc}/(2\pi)}
                         \overline{\chi_{\rm top}(\scal)}}.
\end{equation}
At the completed Born midpoint this gives
\[
 \chi_{\rm top}(\scal_B)
    =-e^{-i7\delta_{\rm loc}/(4\pi)},\qquad
 F_{\rm top}(\scal_B)
    =-\cos\!\left(\frac{7\delta_{\rm loc}}{4\pi}\right).
\]
These equations are exact algebraic consequences of the positive
hierarchy character and do not require a gravitational action.
\end{theorem}

\begin{proof}
Direct substitution gives
\[
 \frac{2\pi}{L_{\rm top}}\log\frac{\scal_j}{U}
   =-2\pi j-\frac{j\delta_{\rm loc}}{2\pi}.
\]
For the Born reflection,
$\log(\mathcal B\scal/U)=\log(C_\sigma/U)-\log(\scal/U)$;
therefore
\[
 \chi_{\rm top}(\mathcal B\scal)
 =\exp\!\left[-2\pi i7
    \frac{L_{\rm top}+\delta_{\rm loc}}{L_{\rm top}}\right]
      \overline{\chi_{\rm top}(\scal)}.
\]
The integer $7$ contributes a unit factor, leaving
Eq.~\eqref{p3d:eq:v222-primitive-born-twist}.  At
$\scal_B=Ue^{-7L_\sigma/2}$ the same substitution gives the
midpoint formulas.
\end{proof}

\begin{remark}[A winding theorem is not a physical phase-selector]
\label{p3d:rem:v222-phase-nonaction-firewall}
On the completed interval $U\to C_\sigma$ the completed character
winds seven times and returns to $1$.  The \emph{primitive} character
on the same completed interval ends at
$e^{-i7\delta_{\rm loc}/(2\pi)}\ne1$; that open phase curve does
not itself define an integer winding number without an additional
closure convention.  The primitive and completed circles must not be
conflated.  Moreover, the normalized contact residue and Clifford rank
fix the periodic \emph{representation}, but a physical phase selection
requires a separately generated operator/coupling.  The framework
continues to derive active depths $0,4,5,7$ without $(\mathrm{FLAG})$.
Neither the cosine nor the sine is inserted into the gravity action.
Any comparison to the four-dimensional de Sitter Ricci curvature,
including a condition $7\delta_{\rm loc}=\log2$, is a new endpoint
matching hypothesis rather than a consequence of this theorem.
\end{remark}

\subsection{Conditional curvature-clock chart and capped comparison}
\label{p3d:sec:v223-cap-clock}
\begin{proposition}[Phase curvature chart and finite Clifford endpoint]
\label{p3d:prop:v223-curvature-clock}
Use the completed logarithmic phase on its stated shell, with
$L=\log\alpha=4\pi^2+\delta$, $\delta=1/(6\sqrt2)$.
On the candidate comparison identifying its scale variable with
$R=R[g]>0$ and its upper anchor with $P$, define
\[
 q(R)=\frac{\log(P/R)}{L},\quad C_j=P\alpha^{-j},\quad
 R_s=P\alpha^{-7/2},\quad x=R/R_s=\alpha^{7/2-q}.
\]
Completed phase returns occur at $q=j$; the protected Clifford endpoint
is $C_7>0$.  For retained cap anchors
\[
 B_U=\alpha^{1/2+\delta/2},\quad
 B_I=\alpha^{7/2-\delta/2},\quad \Delta=B_I-B_U,
 \quad \sigma=B_IR_s,
\]
the profile $B(R)=B_U+2\Delta/(1+\sqrt{1+x^2})$ obeys
\[
 B(q)=B_U+\frac{2\Delta}{1+\sqrt{1+\alpha^{7-2q}}}.
\]
Writing $d=1-B_U/B_I$ and $\epsilon_B=1-B(C_7)/B_I$ gives
\[
 \epsilon_B=\frac{d\alpha^{-7}}{(1+\sqrt{1+\alpha^{-7}})^2}
 \le\frac d4\alpha^{-7}.
\]
For $D=\sigma^2/B^2$ its terminal relative error from $R_s^2$ is
$D(C_7)/R_s^2-1=(1-\epsilon_B)^{-2}-1$.
\end{proposition}
\begin{proof}
Substitute $R=P\alpha^{-q}$ and rationalize
$1-2/(1+\sqrt{1+x^2})=x^2/(1+\sqrt{1+x^2})^2$ at $x^2=\alpha^{-7}$.
The source formula follows from $\sigma=B_IR_s$.
\end{proof}
Numerically the coefficient bound is approximately $1.05403\,10^{-121}$
and the source bound approximately $2.10806\,10^{-121}$.  They are
nonzero controlled endpoint errors.  A domain endpoint does not itself
introduce a freeze outside that domain or new junction conditions.

\begin{remark}[Logarithmic depth is not an auxiliary hyperbolic chart]
\label{p3d:rem:v223-log-asinh}
On this candidate branch
\[
 \ell_{\rm hier}=\log(R/R_s)=L(7/2-q),\qquad
 \zeta=\operatorname{arsinh}(R/R_s),\qquad
 e^{\ell_{\rm hier}}=\sinh\zeta.
\]
Thus $\ell_{\rm hier}$ and $\zeta$ must not be identified.  In
particular $R=0$ is not the Clifford endpoint and is not forced to lie
at $q=7/2$ by the phase hierarchy.  The latter location belongs to a
different optional odd-coordinate identification.  It is not an
obstruction to the declared positive curvature clock.
Here $q=q(R[g])$ is derived: $q_R=-1/(LR)$ and $q_{RR}=1/(LR^2)$.
It has no independent Euler equation or externally prescribed clock
force.  All chain-rule contacts belong to the metric variation.
\end{remark}
The native phase theorem remains a representation on a reduced scale.
Its identification with metric Ricci curvature, its upper dimensional
normalization and the cap anchors above are stated comparison premises,
not consequences of winding alone.  The primitive and completed
characters retain their distinct periods and twist.  No phase cosine
is inserted into the action.  Under the declared reciprocal comparison
$q\mapsto7-q$, $x\mapsto1/x$, the cap chart transforms by
Theorem~\ref{p1b:thm:v223-born-pencil}; it does not transform by $t\mapsto-t$.


\section{Rank-seven Laurent jets and local-to-global scope}
\label{paper3d:sec:organized-laurent}
\subsection{Reciprocal and spectral response indices}

\begin{definition}[Declared two-root response and running indices]
Within \(({\rm RED})+({\rm UV})+({\rm IR})\), let \(0<C<U\) and define
\begin{equation}
\label{p3:eq:geff_def}
 g_{\eff}(\scal)
 :=\frac{\scal^2+CU}{(C+U)\scal}.
\end{equation}
The instantaneous index is
\begin{equation}
\upsilon(\scal)
:=\frac{d\ln g_{\eff}}{d\ln\scal}.
\label{p3:eq:upsilon_def}
\end{equation}
For \(\scal\ne U\), the scale-averaged secant index based at the exact upper
root is
\begin{equation}
\label{p3:eq:n_secant_def}
 n_{\rm sec}(\scal)
 :=\frac{\ln g_{\eff}(\scal)}{\ln(\scal/U)}.
\end{equation}
At \(\scal=U\) it is defined by its continuous limit,
\[
 n_{\rm sec}(U):=\frac{U-C}{U+C}.
\]
Thus \(g_{\eff}(\scal)=(\scal/U)^{n_{\rm sec}(\scal)}\) away from the base and by
continuous extension at the base.
\end{definition}

\begin{remark}[Status of the reduced variables]
\label{p3:rem:reduced_variables_status}
The variable \(\upsilon\) is the instantaneous logarithmic slope, whereas
\(n_{\rm sec}\) is a secant exponent.  The Riccati law below is an identity obeyed by
\(\upsilon\), not an autonomous evolution equation for \(n_{\rm sec}\).  The secant
index is not additive under shell concatenation.  If successive intervals
have logarithmic lengths \(\lambda_i\) and secant indices \(n_{{\rm sec},i}\), then
\[
 n_{{\rm sec},12}
 =\frac{n_{{\rm sec},1}\lambda_1+n_{{\rm sec},2}\lambda_2}
        {\lambda_1+\lambda_2}.
\]
The additive object is the logarithmic transport
\(n_{\rm sec}\lambda\).  This response exponent is not the active entropy
index of the hierarchy-scale state used later.
\end{remark}

\begin{remark}[Exponent-space character evaluation is not secant concatenation]
\label{p3d:rem:v263-exponent-character-pullback}
The preceding nonadditivity statement concerns concatenation of physical
logarithmic intervals.  It does not obstruct evaluation of an independently
constituted positive one-parameter coefficient character on the generated
Laurent exponent.  Concretely, if
\[
 \rho:(\mathbb R,+)\longrightarrow\mathbb R_{>0},
 \qquad \rho(s+t)=\rho(s)\rho(t),
\]
is already generated on its own coefficient owner, then
\[
 R\longmapsto \rho\!\left(n_j^+(R)\right)
\]
is only the pullback of that character along the scalar exponent readout.
No statement
\[
 n_j^+(q_1+q_2)=n_j^+(q_1)+n_j^+(q_2)
\]
is implied or required.  Thus the additive parameter of a coefficient
representation and the nonadditive physical secant coordinate remain
separately typed.
\end{remark}

\begin{theorem}[Riccati identity for the instantaneous index]
\label{p3:thm:riccati}
Let \(\scal_B:=\sqrt{CU}\).  The instantaneous index satisfies
\begin{equation}
\boxed{\frac{d\upsilon}{d\ln \scal} = 1 - \upsilon^2}
\end{equation}
with solution
\begin{equation}
\label{p3:eq:upsilon_solution}
\upsilon(\scal)
=\frac{\scal^2-\scal_B^2}{\scal^2+\scal_B^2}
=\tanh\!\bigl(\ln(\scal/\scal_B)\bigr).
\end{equation}
\end{theorem}

\begin{proof}
The calibrated response of Theorem~\ref{p3:thm:geff_from_reciprocity}
is \(h\) in Paper~I-D,
Theorem~\ref{p1b:thm:laurent-family}~\cite{Kugelmann2026Paper1D}.
With \(R=\scal\), \(v=\upsilon\) and \(\scal_B=\sqrt{CU}\), that
theorem gives both the Riccati identity and
Eq.~\eqref{p3:eq:upsilon_solution}.
\end{proof}

\begin{proposition}[Closed form and continuous base value for the secant index]
\label{p3:prop:n_closed}
For \(\scal>0\), \(\scal\ne U\),
\begin{equation}
\boxed{
n_{\rm sec}(\scal)
=
\frac{\ln\!\big[(\scal^2+CU)/((C+U)\scal)\big]}
     {\ln(\scal/U)}.
}
\end{equation}
At \(\scal=U\), the removable value is
\[
 n_{\rm sec}(U)=\frac{U-C}{U+C}.
\]
\end{proposition}

\begin{proof}
Substituting the calibrated response \eqref{p3:eq:geff_def} into
\eqref{p3:eq:n_secant_def} identifies \(n_{\rm sec}\) with \(n\) in
Paper~I-D, Theorem~\ref{p1b:thm:laurent-family}~\cite{Kugelmann2026Paper1D},
for \(R=\scal\).  Its smooth removable value at \(U\) is the one displayed.
\end{proof}

\begin{proposition}[Maximal nonnegative secant branch]
\label{p3:def:active-reduced-window}
For the continuous secant index of Proposition~\ref{p3:prop:n_closed},
\[
 \boxed{\{\scal>0:n_{\rm sec}(\scal)\ge0\}=[C,\infty).}
\]
This set is the intrinsic nonnegative secant branch
\(\mathcal I_{\rm sec}\).  The upper root \(U\) is an interior removable base
point of the secant quotient, not a boundary or singularity of the response.
\end{proposition}

\begin{proof}
For the calibrated pair \(0<C<U\), Paper~I-D,
Theorem~\ref{p1b:thm:laurent-family}~\cite{Kugelmann2026Paper1D}, gives
\(n_{\rm sec}(\scal)\ge0\) exactly when \(\scal\ge C\), and smoothness
at \(\scal=U\).  This identifies the local branch
\(\mathcal I_{\rm sec}\) with \([C,\infty)\).
\end{proof}

\begin{theorem}[Exact secant range on the intrinsic nonnegative branch]
\label{p3:thm:exact-calibrated-secant-range}
\label{p3:thm:entropy-domain-obstruction-active-branches}
On the intrinsic nonnegative secant branch
\(\mathcal I_{\rm sec}=[C,\infty)\), the secant index is continuous and
strictly increasing, with
\[
 n_{\rm sec}(C)=0,
 \qquad
 n_{\rm sec}(U)=\frac{U-C}{U+C},
 \qquad
 \lim_{\scal\to\infty}n_{\rm sec}(\scal)=1.
\]
Consequently
\[
 0\le n_{\rm sec}(\scal)<1
 \quad\text{for every finite }\scal\ge C.
\]
The whole nonnegative secant branch therefore lies in the numerical entropy
range of the protected octet.  This range comparison does not identify the
secant exponent with the hierarchy-active entropy index.
\end{theorem}

\begin{proof}
Paper~I-D, Theorem~\ref{p1b:thm:laurent-family}~\cite{Kugelmann2026Paper1D},
with \(R=\scal\), gives smoothness, strict monotonicity, both finite
endpoint values and \(0\le n_{\rm sec}<1\) on \([C,\infty)\).
Since \(g_{\eff}(\scal)\sim\scal/(C+U)\) as \(\scal\to\infty\),
the quotient of logarithms tends to one.  The displayed numerical range
is therefore contained in the protected-octet entropy range; this
comparison of ranges makes no identification of the two indices.
\end{proof}

\begin{definition}[Generated neutral visibility family \((\mathrm{NEUT}_7)\)]
\label{p3:def:neutral-visibility-readout}
On the protected split derived by
\(\mathsf C_{\rm IR}^{\rm der}\), the theorem-status certification
\((\mathrm{NEUT}_7)\) denotes the already-generated and co-realized Paper~III-C
basis-neutral common-rate semigroup
\[
 V_\tau=P_0+e^{-\tau}P_W.
\]
Paper~I-C, Proposition~\ref{p1b:prop:trace-coordinate}
\cite{Kugelmann2026Paper1C}, is the mathematical owner of this common-rate
semigroup and its trace-defect coordinate; the III-C construction supplies
the generated role/provenance used here.  Thus the owner citation does not
add a second process occurrence.  III-D does not
reconstruct the semigroup or use it to generate the projectors \(P_0,P_W\).
Its extensive value
\(Z(\tau)=\operatorname{Tr}V_\tau=1+7e^{-\tau}\) is a closed heat loop.
The entropy branch first forms the boundary state
\(\rho_\tau=V_\tau/Z(\tau)\) and then applies the generated state-decorated
loop \(-\operatorname{Tr}(\rho_\tau\log\rho_\tau)\).  Thus
\((\mathrm{NEUT}_7)\) is a generated mixed-role construction in \(\Rgen\),
not a primitive third role.  Whenever an occurring composite contains this
frozen family, Paper~III-C
Corollary~\ref{p3c:cor:corealized-protected-visibility-family} realizes all
of these constituents without a separate neutral-state occurrence tag.
\end{definition}

\begin{theorem}[Unique basis-neutral entropy lift]
\label{p3:thm:basis-neutral-entropy-lift}
On \(\mathsf C_{\rm IR}^{\rm der}\), use the generated family certified by
\((\mathrm{NEUT}_7)\).
Let \(P_0\) and \(P_W=1-P_0\) be the scalar and seven-dimensional
non-scalar projectors, and use the imported \((\mathrm{NEUT}_7)\) semigroup
\[
 V_\tau=P_0+e^{-\tau}P_W,
 \qquad
 \rho_\tau:=\frac{V_\tau}{1+7e^{-\tau}},
 \qquad \tau\in[0,\infty].
\]
Its entropy is
\begin{equation}
\label{p3:eq:neutral-entropy-lift}
 S(\rho_\tau)
 =\ln(1+7e^{-\tau})
  +\frac{7\tau e^{-\tau}}{1+7e^{-\tau}}.
\end{equation}
It decreases strictly from \(\ln8\) to \(0\).  Hence for every
\(\scal\in[C,\infty)\) there is a unique
\(\tau_{\rm sec}(\scal)\in[0,\infty]\) satisfying
\[
 S(\rho_{\tau_{\rm sec}(\scal)})
 =n_{\rm sec}(\scal)\ln8.
\]
This proves existence and uniqueness of an auxiliary secant-matched state
within the \((\mathrm{NEUT}_7)\) family.  That family does not
constitute the \(1\oplus7\) split; it acts on the split already derived by
\(\mathsf C_{\rm IR}^{\rm der}\).  No physical shell-transport
identification is asserted in this theorem.  The auxiliary matching sends
\(\scal=C\) to the scalar state and
\(\scal\to\infty\) to the maximally mixed octet.  The upper response root
\(\scal=U\) is an interior point of this response/entropy branch; identifying
it with the maximally mixed endpoint would require an additional endpoint
map and is not asserted.
\end{theorem}

\begin{proof}
The spectrum of \(\rho_\tau\) consists of
\[
 p_0=\frac1{1+7e^{-\tau}},
 \qquad
 p_W=\frac{e^{-\tau}}{1+7e^{-\tau}}
\]
with multiplicities one and seven.  Direct substitution into
\(-\Tr(\rho\ln\rho)\) gives
\eqref{p3:eq:neutral-entropy-lift}.  Its derivative is
\[
 \frac{d}{d\tau}S(\rho_\tau)
 =-\frac{7\tau e^{-\tau}}{(1+7e^{-\tau})^2}<0
 \qquad(\tau>0).
\]
At \(\tau=0\), \(\rho_0=I_8/8\) and \(S=\ln8\); as
\(\tau\to\infty\), \(\rho_\tau\to P_0\) and \(S\to0\).
Theorem~\ref{p3:thm:entropy-domain-obstruction-active-branches} places
\(n_{\rm sec}(\scal)\ln8\) in this exact range, so continuity and strict
monotonicity give the unique \(\tau_{\rm sec}(\scal)\).  No forward
Paper~III-E theorem is used in this proof.
\end{proof}

\begin{remark}[Chronological status of the physical hierarchy-active state]
\label{p3d:rem:mresp-derived-later-state}
The state \(\rho_{\tau_{\rm sec}(\scal)}\) above is an auxiliary
secant-matched comparison state only.  At the III-D stage no physical
response-scale state is selected by that matching.  On the later independently
constituted Paper~III-E \((\mathrm{MRESP})\) branch, the occurring response
composite contains the generated visibility family and its own scale map, so
its hierarchy-active state \(\sigma_{\rm act}(\scal)\) is
\emph{derived later} by constituent co-realization.  This does not identify
\(\sigma_{\rm act}\) with the secant-matched state and does not derive the
later state-to-Hessian transport law, which remains part of the declared
response class.
\end{remark}


\subsection{Fixed-depth rank-seven Laurent jet}
\label{p3d:sec:v173-n7-local-jet}
Put $C:=C_7$ and
\[
 \mathscr H_7(R):=\frac{R^2+UC}{U+C},\qquad
 n_7(R):=\Pi_+\!\left(n_7^{\rm L}(R)\right).
\]
The exact secant identity is
$\mathscr H_7(R)=R(R/U)^{n_7(R)}$ on the active nonnegative branch, but
$n_7$ is a coefficient/depth readout and is not differentiated by a local
configuration variation.

\begin{theorem}[Actual excess-source derivative and exact local-character certificate]
\label{p3d:thm:v2211-excess-character-certificate}
Fix an active background $\bar R>C$, its fixed scheme/depth coefficient
$n=n_7(\bar R)$, $w=1+n$, and
$\Delta=\mathscr H_7(\bar R)-C>0$. Let $E_{\bar R}(r)$ be the
\emph{actual} $C^1$ local Euler/source excess readout obtained from an
independently supplied complete action/source word, on a positive local
configuration interval about $r=\bar R$, with scheme/depth and the
terminal coefficient $C$ held fixed. Its diagonal-value port is
$E_{\bar R}(\bar R)=\Delta$. Define the evaluable Ward defect
\begin{equation}
\label{p3d:eq:v2211-excess-ward-residual}
 \boxed{\mathcal W_{\bar R}(r)
 :=r\partial_rE_{\bar R}(r)-wE_{\bar R}(r).}
\end{equation}
Then the complete source is exactly
\begin{equation}
\label{p3d:eq:v2211-excess-ward-reconstruction}
 \boxed{E_{\bar R}(r)=
 \Delta(r/\bar R)^w+
 r^w\int_{\bar R}^{r}s^{-w-1}\mathcal W_{\bar R}(s)\,ds.}
\end{equation}
Consequently the pure excess-character germ holds if and only if
$\mathcal W_{\bar R}=0$ on that interval. For only its consumed local
susceptibility the exact weaker necessary and sufficient test is
\begin{equation}
\label{p3d:eq:v2211-excess-local-jet-certificate}
 \boxed{\chi_{\rm act}(\bar R)
 :=\frac{\partial_rE_{\bar R}(\bar R)}{\bar R}
 =\frac{w\Delta+\mathcal W_{\bar R}(\bar R)}{\bar R^2}.}
\end{equation}
We denote an actual normalized diagonal-value port together with an evaluated
$\mathcal W_{\bar R}(\bar R)=0$ by $(\mathrm{EXJET}_7)$; the stronger
interval condition is $(\mathrm{EXCHAR}_7)$. These tags certify a result
of the actual source calculation; they are not axioms setting its defect to
zero. If the independently owned one-variable coefficient is $f_{\bar R}$
and its complete Euler readout is $r f'_{\bar R}-f_{\bar R}$, then
\[
 E_{\bar R}=r f'_{\bar R}-f_{\bar R}-C,\qquad
 \mathcal W_{\bar R}=r^2f''_{\bar R}
 -w(rf'_{\bar R}-f_{\bar R}-C).
\]
The calculation must retain actual operator, density/localizer, field, source
and boundary derivatives. Reconstructing $f$ from the proposed $\chi$ and
then checking this equation is not an independent certificate.
\end{theorem}
\begin{proof}
Differentiate $r^{-w}E_{\bar R}(r)$ at fixed $w$ and $\bar R$:
\[
 \partial_r(r^{-w}E_{\bar R})
 =r^{-w-1}\mathcal W_{\bar R}.
\]
Integrating from $\bar R$ proves the complete formula and the interval
if-and-only-if statement. Evaluation of the defining defect at $\bar R$
gives the local-jet identity. Finally
$\partial_r(rf'_{\bar R}-f_{\bar R})=rf''_{\bar R}$ gives the displayed
actual-action version. No diagonal derivative of $n_7(R)$ is taken.
\end{proof}

\begin{theorem}[Conditional fixed-depth Laurent excess jet]
\label{p3d:thm:v173-fixed-depth-laurent-excess-jet}
On the actual local source domain of
Theorem~\ref{p3d:thm:v2211-excess-character-certificate}, suppose
$(\mathrm{EXJET}_7)$ has been evaluated successfully. Its local
susceptibility is then
\begin{equation}
\label{p3d:eq:v173-n7-local-susceptibility}
 \boxed{\chi_7^{\rm jet}(\bar R)
 =\frac{(1+n_7(\bar R))(\bar R^2-C^2)}{(U+C)\bar R^2}.}
\end{equation}
Under the stronger evaluated $(\mathrm{EXCHAR}_7)$ the unique excess germ
is
\[
 \Delta^{\rm loc}_{7,\bar R}(r)
 =\Delta_7(\bar R)(r/\bar R)^{1+n}.
\]
The supplied susceptibility extends continuously by $\chi_7^{\rm jet}(C)=0$
and obeys
\[
 0<\chi_7^{\rm jet}(R)\le\frac2{U+C}\quad(R>C),\qquad
 \lim_{R\to\infty}\chi_7^{\rm jet}(R)=\frac2{U+C},
\]
and
\begin{equation}
\label{p3d:eq:v173-n7-terminal-jet}
 \boxed{\chi_7^{\rm jet}(C+\delta)
 =\frac{2\delta}{C(U+C)}+O(\delta^2).}
\end{equation}
Without the actual source certificate these are a valid candidate-response
model and its exact consequences, rather than a derivation of a physical
local Hessian from the finite-scale secant alone.
\end{theorem}
\begin{proof}
The zero local Ward defect gives $\chi=w\Delta/\bar R^2$.
Substitute $\Delta=(\bar R^2-C^2)/(U+C)$ to obtain the formula.
The interval certificate gives the stronger unique germ by the preceding
theorem. The bounds use $0\le n_7<1$ and $0<1-C^2/R^2<1$; its upper
limit follows from $n_7\to1$. Smoothness of the active Laurent secant at
$C$, with $n_7(C)=0$, gives $n_7(C+\delta)=O(\delta)$ and the stated
expansion. No value/character theorem evaluates the actual Ward defect.
\end{proof}

\begin{proposition}[Exact nonuniqueness before the excess-source port]
\label{p3d:prop:v2211-excess-jet-counterfamily}
Put $H=\mathscr H_7(\bar R)$. The germs
\[
 E_A(r)=(H-C)(r/\bar R)^w,\qquad
 E_B(r)=H(r/\bar R)^w-C
\]
have the same diagonal excess, fixed coefficient $n$, and terminal constant,
but their defects are $\mathcal W_A=0$ and $\mathcal W_B=wC$.
Their inverse-Euler jets differ by $wC/\bar R^2$.
Moreover for every fixed $\beta>0$,
\[
 E_\beta(r)=\Delta(r/\bar R)^\beta,\qquad
 \chi_\beta(\bar R)=\frac{\beta(\bar R^2-C^2)}{(U+C)\bar R^2}
\]
is a single positive local character with the same diagonal value and
terminal Hessian decay. The member $\beta=2$ also has the printed
high-curvature limit. Thus the exact secant, tangent-sort distinction,
positivity and terminal decay do not by themselves select $w=1+n$ as the
\emph{excess} representation. An actual Ward/source evaluation is required.
\end{proposition}
\begin{proof}
Direct differentiation gives $rE_A'=wE_A$ and
$rE_B'=w(E_B+C)$. For $E_\beta$ the defect is
$(\beta-w)E_\beta$, giving the stated jets. Their common diagonal data
therefore leave the local derivative unresolved. These functions show the
missing information; they are not promoted to additional physical actions.
\end{proof}

\begin{remark}[Secant, visibility and local jet are distinct]
\label{p3d:rem:v173-three-n7-sorts}
The exact secant exponent $n_7^{\rm L}$, the normalized unresolved trace
$\nu_7$, and the local susceptibility $\chi_7^{\rm jet}$ are three distinct
readouts.  The first fixes the finite-scale Laurent character, the second is
a visibility trace, and the third is the fixed-depth local configuration
second jet.  No equality among these sorts is used.
\end{remark}

\subsection{Complete action jets and the local/global response distinction}
\label{p3d:sec:v2216-complete-action-jets}

A fixed-background configuration derivative and a derivative after pulling
background coefficients back to curvature are different operations. The
following tests retain this distinction without modifying the Laurent
response or its terminal data. They apply in one common physical pairing
after the operator, density, localizer, source and boundary maps have been
supplied. Co-realization alone does not supply those maps; see
Theorem~\ref{p3a:thm:v2216-common-action-consumption}.

\begin{theorem}[Complete action jets and exact matching residual]
\label{p3d:thm:v2216-complete-action-jets}
Let $I=(C,R_+)$ with $C>0$, let $B$ be a regular finite coefficient chart,
and let $A\in C^2(B\times I)$ be an independently supplied scalar action
coefficient in a fixed physical unit and density chart. Let $b:I\to B$ be
the actual $C^2$ coefficient map and put
$f_{\rm act}(R)=A(b(R),R)$. Then
\begin{align}
 f_{\rm act}''&=A_{rr}+\mathcal G,\notag\\
 \mathcal G&=2A_{rb}[b']+A_{bb}[b',b']+A_b[b''],
 \label{p3d:eq:v2216-full-pullback}
\end{align}
where all partial derivatives are evaluated at $(b(R),R)$.
For a supplied continuous target $K$, integrable at $C$, define
\[
 f_K(R)=R-C+\int_C^R(R-t)K(t)\,dt.
\]
Suppose $f_{\rm act},f_{\rm act}'$ extend to $C$ and the following
integral exists. With
$\delta_0=f_{\rm act}(C)$, $\delta_1=f_{\rm act}'(C)-1$ and
$\mathcal M=f_{\rm act}''-K$, the exact residual is
\begin{equation}
 \boxed{f_{\rm act}-f_K=\delta_0+\delta_1(R-C)+
          \int_C^R(R-t)\mathcal M(t)\,dt.}
 \label{p3d:eq:v2216-action-residual}
\end{equation}
Thus equality holds on $I$ if and only if both endpoint ports agree and
$\mathcal M=0$. For $K=A_{rr}(b(R),R)$, the latter test is exactly
$\mathcal G=0$.

In a scalar background chart define the frozen diagonal jets
$a(R)=A(R,R)$, $p(R)=A_r(R,R)$ and $k(R)=A_{rr}(R,R)$. They are the
complete two-jet of one function exactly when
\[
 a'=p,\qquad p'=k,
 \quad\hbox{equivalently}\quad A_b(R,R)=A_{br}(R,R)=0.
\]
Agreement only of the diagonal Hessian with $k$ is the weaker condition
$2A_{br}+A_{bb}=0$ on that diagonal.
\end{theorem}
\begin{proof}
Differentiate $A(b(R),R)$ twice to obtain
\eqref{p3d:eq:v2216-full-pullback}. Twice integrating
$(f_{\rm act}-f_K)''=\mathcal M$ with its two endpoint values gives
\eqref{p3d:eq:v2216-action-residual} and its equivalence. For the last
statement the chain rule gives $a'=A_r+A_b$ and
$p'=A_{rr}+A_{rb}$. Necessity follows from the definition of a two-jet;
conversely the two identities make $a$ the required function. The weaker
condition follows by differentiating $a$ twice directly.
\end{proof}

\begin{remark}[Density and two-endpoint tests of the same action]
\label{p3d:rem:v2217-density-matching-import}
The density derivatives in Paper~III-C,
Proposition~\ref{p3c:prop:v2217-determinant-density-jets}, are included in
$A$ and its complete pullback above. Paper~III-E,
Theorem~\ref{p3e:thm:v2217-uv-ir-moment-nonuniqueness}, supplies the exact
two-endpoint moment and chart-gluing tests. Neither determinant scaling nor
smooth UV/IR endpoint data substitute for evaluating the independent
same-action residual; the existing global EXJET obstruction is retained.
\end{remark}


\begin{remark}[Full physical matching retains more than the scalar port]
The theorem compares one scalar coefficient. Equality of complete physical
actions additionally requires every retained metric, matter, mixed-source
and boundary variation, together with the same domain and pairing. For
$\Gamma(q)=\mathcal S(\iota(q))$ the general second-variation rule is
\[
 D^2\Gamma[u,v]=D^2\mathcal S[D\iota u,D\iota v]
                    +D\mathcal S[D^2\iota(u,v)].
\]
Thus a field/source map carries a second contact even when its linearized
map is known. Multiplying or dividing an action by a moving density carries
the corresponding product-rule terms. Coefficients that are external
marks remain fixed in a physical variation; substituting $b=b(R)$ is
licensed only by an actual pullback. Endpoint matching and a supplied
Hessian give uniqueness of a reconstruction, not an origin certificate for
an independently owned action.
\end{remark}

\begin{theorem}[Terminal global action obstruction for a bounded EXJET character]
\label{p3d:thm:v2216-exjet-global-obstruction}
Let $f\in C^2([C,R_1])$ satisfy $f(C)=0$, $f'(C)=1$, and let $w$ be
continuous. If its \emph{same global} Euler excess
$E(R)=Rf'(R)-f(R)-C$ obeys
\[
 RE'(R)-w(R)E(R)=0
\]
at every point, then $E\equiv0$ and $f''\equiv0$.

In particular let $w=1+n_7$, $D=U+C$, and
\[
 \Delta(R)=\frac{R^2-C^2}{D},\qquad
 \chi_7(R)=\frac{w(R)\Delta(R)}{R^2}.
\]
The terminally normalized integral coefficient and its excess are
\[
 f_{\rm gl}=R-C+\int_C^R(R-t)\chi_7(t)\,dt,\qquad
 E_{\rm gl}=\int_C^R t\chi_7(t)\,dt.
\]
On every finite active interval $R>C$,
\begin{equation}
 \boxed{\mathcal W_{\rm gl}:=R^2\chi_7-wE_{\rm gl}
     =w(\Delta-E_{\rm gl})>0.}
 \label{p3d:eq:v2216-global-exjet-residual}
\end{equation}
Consequently the original bounded-character EXJET relation cannot be imposed
on this same global metric action at every background. Local certification
of a separate frozen source family remains meaningful, but its consumption
as a global Hessian requires a separately supplied second-jet transport.
\end{theorem}
\begin{proof}
The endpoint data give $E(C)=0$. Since $w/R$ is continuous for $C>0$,
the integrating factor gives
$E(R)=E(C)\exp\{\int_C^R w(t)\,dt/t\}=0$. The identity
$E'=Rf''$ then proves the first statement. On the native active branch,
$0\le n_7<1$ and $0<\chi_7<2/D$. Integration against $t>0$ gives
\[
 0<E_{\rm gl}(R)<\int_C^R\frac{2t}{D}\,dt=\Delta(R).
\]
Substitution of $R^2\chi_7=w\Delta$ proves the strict residual. This
obstruction does not differentiate a fixed depth during a local variation:
it tests whether the entire family can be one ordinary global function.
\end{proof}

\begin{corollary}[Raw Laurent Euler values do not glue with the local Hessian]
\label{p3d:cor:v2216-raw-euler-nongluing}
No $C^2$ coefficient on a nonempty undressed active interval can satisfy
both
\[
 Rf'-f=\frac{R^2+UC}{U+C},\qquad f''=\chi_7.
\]
\end{corollary}
\begin{proof}
Differentiating the first identity forces $f''=2/(U+C)$, whereas the
strict inequality in the preceding proof gives $\chi_7<2/(U+C)$.
An affine correction changes neither second derivative.
\end{proof}

\subsection{Finite coefficient tests of complete running-response jets}
\label{p3d:sec:v226-response-coefficients}
\begin{proposition}[A finite native jet identity has complete scalar tests]
\label{p3d:prop:v226-response-coefficients}
Let $A(r),\widetilde A(r)$ be actual $C^k$ native response matrices in the
same generated level-one momentum frame. Their jets through order $k$ agree at
$r_*$ exactly when
\[
 \left.\partial_r^\ell\operatorname{Tr}
       [E_j(\widetilde A(r)-A(r))]\right|_{r_*}=0,
 \quad j=1,\ldots,137,\quad 0\le\ell\le k.
\]
In a transported unitary frame use the covariant derivatives of
Corollary~\ref{p1b:cor:v226-source-jets}.
\end{proposition}
\begin{proof}
Apply the fixed linear dual inverse of
Theorem~\ref{p1b:thm:v226-dual-coefficients} to every derivative of the
actual difference. Vanishing of all coefficients is necessary and sufficient.
\end{proof}
A tested scalar trace can close a smaller consumed quotient; it need not
reconstruct an unused matrix. Conversely a scalar profile match does not
prove this complete native jet identity. Source incidence, density,
localizer and boundary derivatives remain in the actual matrix words.
This finite test leaves the current Laurent/capped profile and selected
excess-character obligations on their existing physical domains.


\subsection{Global energy comparison from local source jets}
\label{p3d:sec:v227-running-frequency}
\begin{corollary}[Finite energy tests close the consumed local running jets]
\label{p3d:cor:v227-running-frequency}
At a fixed running/source point, two actual finite Hermitian response
families in one fixed kinetic frame have identical consumed jets through
order $k$ as functions of all meromorphic energies if those jets agree at
$\lfloor(k+1)(m_1+m_2)/2\rfloor+1$ distinct upper-half-plane energies.
\end{corollary}
\begin{proof}
Apply Corollary~\ref{p1b:cor:v227-response-jets} to the actual running
source derivatives. The degree bound is in the spectral energy variable.
\end{proof}
This strengthens the fixed-energy coefficient tests of the preceding
section. It does not extrapolate arbitrary $C^k$ curvature dependence
away from the source point, identify the hierarchy index with energy,
or provide an uncomputed capped/Laurent profile. A curvature-dependent
energy frame consumes its actual frame/clock derivatives before applying
the comparison theorem in a common normalization.


\section{Positive mathematical providers and fixed-momentum sources}
\label{paper3d:sec:organized-providers}
\subsection{Positive mathematical providers for the Born second jet}
\label{p3d:sec:v2216-positive-born-provider}

The conditional response has two explicit positive mathematical providers.
They establish existence and useful approximation results. Neither identifies
the physical action by fitting it to the proposed susceptibility. In this
section $0<C<U$, $D=U+C$, and $n=n_7$ on its active nonnegative branch.
The physical constructor and complete-source requirements of
Theorem~\ref{p3d:thm:v2216-complete-action-jets} remain in force.

\begin{theorem}[Positive resolvent measure and terminal-null Schur provider]
\label{p3d:thm:v2216-positive-born-provider}
Put $a_0=U^2$, $b_0=UC$, $c_0=C^2$, and define
\[
 F_{\rm B}(x)=\frac{\log[(x+b_0)/(a_0+b_0)]}{\log(x/a_0)}.
\]
The value at $x=a_0$ is removable. There is a unique atom-free probability
measure $\mu$ on $(0,\infty)$ such that
\begin{equation}
 F_{\rm B}(x)=\int\frac{x}{x+s}\,\mu(ds),\qquad
 \chi_7(R)=\frac2D\int\frac{R^2-c_0}{R^2+s}\,\mu(ds).
 \label{p3d:eq:v2216-born-resolvent}
\end{equation}
Its density is
\[
 \frac{d\mu}{ds}=
 \begin{cases}
 \displaystyle\frac{\log[(a_0+b_0)/(b_0-s)]}
 {s[(\log(s/a_0))^2+\pi^2]},&0<s<b_0,\\[6pt]
 \displaystyle\frac{\log[s(a_0+b_0)/(a_0(s-b_0))]}
 {s[(\log(s/a_0))^2+\pi^2]},&s>b_0.
 \end{cases}
\]
It has $F_{\rm B}(c_0)=1/2$ and infinite first moment.

Let $h_\infty=2/D$, $\mathcal H=L^2(\mu)$ and
\[
 \mathcal V=\{z:\int(c_0+s)|z|^2\,d\mu<\infty\}.
\]
For $R\ge C$ the closed nonnegative form on
$\mathbb R\oplus\mathcal V$
\[
 Q_R(H,z)=\frac{h_\infty H^2}{2}
 +\frac12\int(R^2+s)z^2\,d\mu
 +H\sqrt{h_\infty}\int\sqrt{c_0+s}\,z\,d\mu
\]
has exact stationary reduction
$\inf_z Q_R(H,z)=\chi_7(R)H^2/2$. Its interior remains gapped at $C$,
while the full form has the terminal null direction
$z_C=-H\sqrt{h_\infty}/\sqrt{c_0+s}$.
\end{theorem}
\begin{proof}
Use principal logarithms on $\mathbb C\setminus(-\infty,0]$ and put
$G(z)=F_{\rm B}(z)/z$. At $z=-s+i0$ its Stieltjes inversion density
$-\operatorname{Im}G(-s+i0)/\pi$ is exactly the displayed positive
function. For $s<b_0$ the numerator logarithm is real and negative; for
$s>b_0$ it has imaginary part $\pi$. Division by
$\log(s/a_0)+i\pi$ gives the two formulas. Their numerators are positive,
in the second case because $s(a_0+b_0)>a_0(s-b_0)$.

A keyhole Cauchy contour gives $G(z)=\int(z+s)^{-1}\mu(ds)$.
Its inner circle vanishes because $G(z)=O(1/(z\log z))$ near zero;
the outer circle vanishes because $G(z)=O(1/z)$ at infinity. The
logarithmic singularity at $-b_0$ has zero indentation contribution, and
the point $a_0$ is removable. There are therefore no omitted atoms or
affine terms. Stieltjes inversion proves uniqueness. Multiplication by
$x$ gives the first representation, and $F_{\rm B}(x)\to1$ plus
monotone convergence gives total mass one. Algebra in the native secant
gives $1+n(R)=2F_{\rm B}(R^2)$ and hence the second representation.
Direct substitution gives $F_{\rm B}(c_0)=1/2$. The density is asymptotic
to $\log(1+b_0/a_0)/(s[\log(s/a_0)]^2)$ at infinity, proving that its
first moment diverges.

Weighted Cauchy--Schwarz makes the loading a bounded functional on
$\mathcal V$. Completing the square gives
\[
 Q_R=\frac{\chi_7(R)H^2}{2}
 +\frac12\int(R^2+s)
 \left(z+\frac{H\sqrt{h_\infty}\sqrt{c_0+s}}{R^2+s}\right)^2d\mu.
\]
The shift belongs to $\mathcal V$, since
$\int(c_0+s)/(R^2+s)\,d\mu\le1$. This triangular shift is bounded and
invertible in the form norm on each bounded $R$ interval; the displayed
square together with the underlying Hilbert norm is complete. The form is
therefore closed and nonnegative. Its unique stationary shift proves the
reduction and the stated null direction. The interior multiplication
operator is bounded below by $C^2>0$ throughout.
\end{proof}

\begin{corollary}[Positive finite approximations and exact terminal slope]
\label{p3d:cor:v2216-born-positive-approximation}
Partition a central interval $[s_-,s_+]$ into cells of length at most
$\epsilon(c_0+\inf I_j)$, keep both tails as cells, and assign each its
actual mass $m_j=\mu(I_j)$. Positive finite nodes yield
\[
 \chi_N(R)=h_\infty\sum_jm_j\frac{R^2-c_0}{R^2+s_j},\qquad
 \sup_{R\ge C}|\chi_N-\chi_7|
 \le h_\infty(\eta+\epsilon/4),
\]
where $\eta$ is the total tail mass. Each positive-mass cell can use its
harmonic node
\[
 \frac1{c_0+s_j}=\frac1{m_j}\int_{I_j}\frac{d\mu}{c_0+s}.
\]
Then probability, terminal zero and
$\chi_N'(C)=\chi_7'(C)=2/[C(U+C)]$ hold exactly.
\end{corollary}
\begin{proof}
For $X=R^2-c_0\ge0$,
$|\partial_s[X/(X+c_0+s)]|\le1/[4(c_0+s)]$ uniformly in $X$.
The central oscillation is at most $\epsilon/4$ and each tail oscillation
at most one. Sum against the actual cell masses. The harmonic mean lies
inside its cell and is finite even on the unbounded upper tail. It
preserves $\int(c_0+s)^{-1}d\mu=1/(2c_0)$, which gives the derivative
identity by differentiating the finite and integral response kernels.
\end{proof}

\begin{remark}[Provider and measure scope]
The loading $\sqrt{h_\infty}\sqrt{c_0+s}$ is in the form dual, not in
$\mathcal H$, because the first moment diverges. The theorem therefore
does not posit a bounded off-diagonal Hilbert matrix. Nor does it posit a
trace-class Radon Gaussian covariance on this continuous $L^2$ space.
Classical Schur reduction and a Gaussian log-determinant are different
operations. Each finite conditional Gaussian has a positive interior pivot,
but its unrestricted joint endpoint Gaussian requires a specified null
quotient. No finite fixed quadratic pencil gives the exact logarithmic
function: its response would be rational in $R^2$, whereas
$F_{\rm B}$ has a genuine logarithmic branch at $-UC$. This excludes
that finite-pencil replacement, not finite nonlinear providers. The measure
is a dilation of a scalar response and is not a particle-mass distribution
or a physically promoted infinite carrier.
\end{remark}

\begin{theorem}[Finite nonlinear Born second-jet provider]
\label{p3d:thm:v2216-finite-born-second-jet}
Put $L=\log(U/C)$, $a=L/2$, $s=\log(R/\sqrt{UC})$ and let
$\Psi(0)=0$, $\Psi'(x)=x+\log\cosh x$. The native secant is
\[
 n(R)=\frac{\log\cosh s-\log\cosh a}{s-a},
\]
with its removable value at $s=a$. On
$\mathbb Re_0\oplus W$, $\dim W=7$, take the normalized trace
$\tau_8=\operatorname{Tr}/8$ and set
\[
 A_R=aP_0+sP_W,\qquad B(v)=e_0v^{\mathsf T}+ve_0^{\mathsf T},
\]
\[
 \mathcal F_R(v)=\beta_R\tau_8
   [\Psi(A_R+\kappa_RB(v))-\Psi(A_R)],\qquad \beta_R\ge0.
\]
This is an $O(W)$-invariant, even, nonnegative finite scalar word with
\begin{equation}
 \mathcal F_R(0)=0,\quad D_v\mathcal F_R(0)=0,\quad
 D_v^2\mathcal F_R(0)=\frac{\beta_R\kappa_R^2}{4}(1+n(R))I_W.
 \label{p3d:eq:v2216-finite-born-hessian}
\end{equation}
Writing its scalar orbit function as $\varphi_R(|v|^2)$, the frozen
local curvature germ
\[
 f_R(r)=a_Rr+b_R+\varphi_R(\log^2(r/R))
\]
has
\begin{equation}
 f_R''(R)=\frac{\beta_R\kappa_R^2}{4R^2}(1+n(R)),\qquad
 f_R'''(R)=-\frac3R f_R''(R).
 \label{p3d:eq:v2216-finite-born-third-jet}
\end{equation}
Its second jet is the proposed $\chi_7$ exactly when the independent
source-Gram comparison is $\beta_R\kappa_R^2=4\Delta(R)$. With that
comparison, the original EXJET diagonal-value and local Ward ports further
require $b_R=-C-\Delta(R)$.
\end{theorem}
\begin{proof}
The only two moving eigenvalues are
\[
 \lambda_\pm=\frac{a+s}{2}\pm
       \sqrt{\frac{(a-s)^2}{4}+\kappa_R^2|v|^2};
\]
the other six remain $s$. The relative trace is therefore
$\beta_R[\Psi(\lambda_+)+\Psi(\lambda_-)-\Psi(a)-\Psi(s)]/8$.
It is even and invariant, and nonnegative because
$\Psi''=1+\tanh x>0$ and the moving pair has fixed sum with increasing
separation. Its Hessian is
\[
 \frac{\beta_R\kappa_R^2}{4}
       \frac{\Psi'(s)-\Psi'(a)}{s-a}I_W.
\]
The divided difference is $1+n(R)$, including its removable value at
$s=a$. The first two configuration jets follow since the orbit word is
even in $z=\log(r/R)$, with $z'=1/R$, $z''=-1/R^2$ at $r=R$.
The third derivative is $3D_z^2\mathcal F_R(0)z'z''$, proving
\eqref{p3d:eq:v2216-finite-born-third-jet}. At the center the Euler
excess is $Rf_R'-f_R-C=-b_R-C$. Equating it to $\Delta$ fixes $b_R$;
the Gram identity then gives $R^2f_R''=w\Delta$, the local Ward test.
\end{proof}

\begin{proposition}[Finite-word profile, equivariance and higher-jet limits]
\label{p3d:prop:v2216-finite-born-limits}
The finite construction proves a second-jet existence result, not a unique
action profile or equality of complete source germs with $f_{\rm gl}$.
Indeed:
\begin{enumerate}
\item The strictly convex profiles
\[
 \Psi_\lambda(x)=\tfrac12x^2+
       \lambda\int_0^x\log\cosh t\,dt,\qquad -1\le\lambda\le1,
\]
all have $\Psi_\lambda(x)+\Psi_\lambda(-x)=x^2$ but give a divided
difference $1+\lambda n$. Positivity and reciprocal reflection therefore
do not select $\lambda=1$.
\item No nonzero scalar-to-$W$ lift is $O(W)$-equivariant when the scalar
has the trivial action. The scalar orbit prescription above is invariant,
but its use as a physical curvature source remains an incidence to prove.
\item With the matching Gram, the complete finite logarithmic germs cannot
be restrictions of a single $C^3$ global coefficient having Hessian
$\chi_7$, even after arbitrary affine corrections to each germ.
\end{enumerate}
\end{proposition}
\begin{proof}
The first statement follows from
$\Psi_\lambda''=1+\lambda\tanh x>0$ at finite $x$ and oddness of the
integral of the even function $\log\cosh t$. For the second, an
equivariant lift of a fixed scalar must be fixed by every orthogonal
transformation; the only such vector is zero. Taking an orbit function
avoids choosing a vector and does not prove a field map.

For the last claim write $F(s)=\log\cosh s$ and
$n=[F(s)-F(a)]/(s-a)$. Strict convexity gives
\[
 n_s=\frac{F'(s)-[F(s)-F(a)]/(s-a)}{s-a}>0
\]
for $s\ne a$, with positive limit $F''(a)/2$. Thus $n_R>0$ and
both positive factors in
$\chi_7=(1+n)(1-C^2/R^2)/D$ increase, so $\chi_7'(R)>0$ for $R>C$.
A global coefficient with this Hessian has third derivative
$\chi_7'>0$. The frozen germ instead has $f_R'''=-3\chi_7/R<0$ by
\eqref{p3d:eq:v2216-finite-born-third-jet}. Affine corrections change
neither jet. At $C$ a zero Gram makes the entire frozen nonlinear word
zero, whereas the global third derivative is $2/(CD)>0$.
\end{proof}

\begin{remark}[Native profile and terminal source comparison]
The required $\Psi''=1+\tanh x$ increases strictly. A positive proper-time
mixture on this same spectral coordinate, with fixed real shift, has
$\widetilde\Psi''(x)=\int t^2e^{-tx}\,d\rho(t)$, $\rho\ge0$,
which is nonincreasing wherever it exists. Thus this finite word is not
identified with the native heat constructor by a positive heat mixture.
Changing the operator, allowing signed subtractions or supplying another
effective constructor would change the claim and requires its own license.
At $R=C$, reciprocal reflection gives exactly
$\mathcal F_C(v)=\beta_C\kappa_C^2|v|^2/8$: the secant factor is one,
so its terminal response zero comes from the source Gram. A possible
factorization $\kappa_R^2=1-C^2/R^2$, $\beta_R=4R^2/D$ meets that
mathematical Gram identity, but its matrix lift is not $C^1$ at $C$.
The invariant trace is regular in $\kappa_R^2|v|^2$; a physical quotient
domain and the affine value/slope ports still have to be evaluated.
\end{remark}

\begin{corollary}[Unique scalar transgression and its physical premise]
\label{p3d:cor:v2216-second-jet-transgression}
An independently generated scalar Hessian $K$ in the actual curvature
coordinate has the unique terminal-normalized primitive
\[
 \mathcal T_C[K](R)=R-C+(R-C)^2
       \int_0^1(1-u)K(C+u(R-C))\,du.
\]
It has $\mathcal T_C[K]''=K$, value zero and slope one at $C$.
For the matched finite second jet this is $f_{\rm gl}$. Its construction
consumes only that second jet and supplied endpoint data; it is not an
identity of the finite source germs, a profile-selection theorem or a
retroactive origin proof for the heat action.
\end{corollary}
\begin{proof}
The substitution $t=C+u(R-C)$ gives
$\int_C^R(R-t)K(t)\,dt$. Differentiate twice. The difference of any two
such primitives has zero second derivative and is affine; the endpoint
data set its two coefficients to zero. The nonidentity follows from
Proposition~\ref{p3d:prop:v2216-finite-born-limits} and
Theorem~\ref{p3d:thm:v2216-exjet-global-obstruction}.
\end{proof}


\subsection{Fixed-momentum finite source words and a composed interpolation}
\label{p3d:sec:v2218-fixed-momentum-source}

The generated level-one momentum seed and radius remain fixed in
this section: $\mathcal H_1$ and $R_{\rm mom}$ are not varied.
This is a finite-seed response construction, not a proof of universal
physical full-endpoint multiplicity-one.  The depth-dependent coefficients below are response/readout data
on that fixed structure, not a change of momentum representation or a
spatial-radius homothety.  A different middle profile is allowed when
testing inflation and late Einstein gravity with its vacuum scale.
Consequently the exact-profile exclusions above are not general
exclusions of a common interpolating action.  They do, however, require
the complete finite value and its source contacts to be evaluated before
a frozen susceptibility is identified with a metric coefficient.

\begin{definition}[Exact finite Born orbit value]
\label{p3d:def:v2218-finite-orbit-value}
Keep $0<C<U$, $D=U+C$, $H=\log(U/C)$, $a=H/2$, and the existing
profile $\Psi$ of Theorem~\ref{p3d:thm:v2216-finite-born-second-jet}.
For $s\in\mathbb R$ and $y\ge0$ define
\begin{align}
 d(s,y)&=\sqrt{(a-s)^2/4+y},\qquad
 \lambda_\pm(s,y)=(a+s)/2\pm d(s,y),\nonumber\\
 \Phi(s,y)&=\frac18\bigl[
 \Psi(\lambda_+)+\Psi(\lambda_-)-\Psi(a)-\Psi(s)\bigr].
 \label{p3d:eq:v2218-finite-orbit-value}
\end{align}
Thus the complete finite source word is
$\mathcal F_R(v)=\beta_R\Phi(s_R,\kappa_R^2|v|^2)$, with
$s_R=\log(R/\sqrt{UC})$.  No additional spectral profile has been
introduced.  The normalized trace is still the real eight-channel
readout of that theorem; its identification with an ordinary physical
finite-Hilbert action trace is not supplied by this notation.
\end{definition}

\begin{proposition}[Complete orbit contacts and terminal reflection]
\label{p3d:prop:v2218-finite-orbit-contacts}
Let an actual scalar source section supply twice differentiable
$\beta(R)$ and $y(R)$, and put
$\mathscr W(R)=\beta(R)\Phi(s_R,y(R))$.  The complete derivatives are
\begin{align}
 \mathscr W'&=\beta'\Phi+
       \beta(\Phi_s/R+\Phi_y y'),\nonumber\\
 \mathscr W''&=\beta''\Phi+
 2\beta'(\Phi_s/R+\Phi_y y')\nonumber\\
 &\quad+\beta\left[
 \frac{\Phi_{ss}-\Phi_s}{R^2}
 +\frac{2\Phi_{sy}y'}R+\Phi_{yy}(y')^2+\Phi_y y''\right].
 \label{p3d:eq:v2218-complete-orbit-jets}
\end{align}
At $y=0$, write
$w(s)=[\Psi'(s)-\Psi'(a)]/(s-a)=1+n(s)$.  Then
\begin{gather}
 \Phi(s,0)=\Phi_s(s,0)=\Phi_{ss}(s,0)=0,
 \qquad \Phi_y(s,0)=w(s)/8,\nonumber\\
 \Phi_{yy}(s,0)=
 \frac{\Psi''(s)+\Psi''(a)-2w(s)}{8(s-a)^2},
 \qquad
 \boxed{\Phi(-a,y)=y/8.}
 \label{p3d:eq:v2218-orbit-zero-jets}
\end{gather}
The removable value of $\Phi_{yy}(a,0)$ is $\Psi''''(a)/48$.
\end{proposition}
\begin{proof}
The two moving eigenvalues give
\eqref{p3d:eq:v2218-finite-orbit-value}; the other six cancel.
The chain rule uses $s_R'=1/R$ and $s_R''=-1/R^2$ and gives
\eqref{p3d:eq:v2218-complete-orbit-jets}.  Differentiating the symmetric
eigenvalue pair with respect to $y$ at zero gives the displayed divided
differences.  Taylor expansion at $s=a$ gives the removable value.
At $s=-a$ the pair is
$\pm\sqrt{a^2+y}$, so
$\Psi(x)+\Psi(-x)=x^2$ proves the exact last identity.
Symmetric spectral calculus makes the orbit value smooth in $y$ at the
removable eigenvalue coincidence.  Thus these formulas also apply to
the regular invariant quotient when a chosen square-root matrix lift
itself is not differentiable at a terminal zero.
\end{proof}

\begin{proposition}[The natural frozen diagonal and free elimination]
\label{p3d:prop:v2218-diagonal-cancellation}
For the frozen germ of
Theorem~\ref{p3d:thm:v2216-finite-born-second-jet}, write its nonlinear
part as
\[
 \mathscr W(b,r)=\mathcal F_b(e\log(r/b)),\qquad |e|=1,
 \quad K_{\rm loc}(b)=\partial_r^2\mathscr W(b,r)|_{r=b}.
\]
At every regular background, including all coefficient derivatives,
\begin{equation}
 \boxed{\mathscr W(b,b)=0,\qquad
 \nabla^2_{(b,r)}\mathscr W|_{r=b}
 =K_{\rm loc}(b)
 \begin{pmatrix}1&-1\\-1&1\end{pmatrix}.}
 \label{p3d:eq:v2218-diagonal-cancellation}
\end{equation}
The complete pullback along $b=r=R$ therefore has zero nonlinear
Hessian.  Moreover, minimizing this isolated nonnegative word over
$v$ gives zero for every $b$; when $\beta_b\kappa_b^2>0$ the minimum
is attained only at $v=0$.
\end{proposition}
\begin{proof}
For every $b$, both $\mathcal F_b(0)$ and its first source derivative
are zero.  Coefficient derivatives consequently vanish at that source.
The first derivatives of $\log(r/b)$ at the diagonal are $1/b$ and
$-1/b$.  Their tensor square gives the matrix in
\eqref{p3d:eq:v2218-diagonal-cancellation}; contraction with $(1,1)$
is zero.  Strict convexity of $\Psi$ gives strict positivity of its
trace second variation along each nonzero off-diagonal source line.
The relative word is consequently positive away from zero when its
prefactor and source scaling are nonzero.  If either vanishes the
entire word is zero, with no unique minimizing source asserted.
\end{proof}

The preceding proposition concerns the nonlinear source word, not
independently supplied affine background terms.  A different actual
source section, a constraint, state preparation or elimination in a
larger common action can change the result.  Those data must be
evaluated; the positive frozen probe response cannot replace them.

\begin{definition}[Fixed-momentum visibility-norm candidate]
\label{p3d:def:v2218-fixed-momentum-candidate}
Fix a positive hierarchy scheme, write
$\alpha=\alpha_{\rm hier}(\sigma_{\rm fin})>1$,
$C=U\alpha^{-7}$ and $Q=8\sqrt\alpha$, and retain the existing
upper quadratic coupling $Q$ as a constant in the following composed
candidate.  For $C\le R\le U$, set
\begin{gather}
 \zeta(R)=\frac{\log(R/C)}{H},\qquad
 k(R)=1-C^2/R^2,\nonumber\\
 y(R)=k(R)\zeta(R)^2,\qquad
 \beta_Q(R)=4QR^2/D,\nonumber\\
 \boxed{f_Q(R)=R-C+
     \frac{4QR^2}{D}\Phi\bigl(s_R,k(R)\zeta(R)^2\bigr).}
 \label{p3d:eq:v2218-fixed-momentum-candidate}
\end{gather}
The scalar orbit norm is constructible from the already generated
visibility process: $q=7(1-e^{-\tau})$ and $R=U\alpha^{-q}$ give
\begin{equation}
 \boxed{\zeta=e^{-\tau},\qquad
 \zeta^2=\frac17\operatorname{Tr}_W(V_\tau^2).}
 \label{p3d:eq:v2218-visibility-source-norm}
\end{equation}
The sufficient Gram factorization is
$\kappa_R^2=k(R)$, $\beta_R=\beta_Q(R)$, and hence
$\beta_R\kappa_R^2=4Q\Delta(R)$,
$\Delta(R)=(R^2-C^2)/D$.  It is the existing possible factorization
with one fixed upper coupling, not a consequence that the Gram
comparison selects $\beta_R$ and $\kappa_R$ separately.  Extending
that coupling as a constant and using the visibility norm as this
finite source section define the candidate; their physical incidence
is not inferred from the availability of the ingredients.
\end{definition}

No nonzero scalar-to-$W$ equivariant vector lift has been introduced.
The invariant orbit value is independent of a chosen vector
representative, but this does not prove an actual field coupling.
Likewise, using the existing finite $\Psi$ word does not identify it
with the selected ordinary heat action.  These distinctions keep
\eqref{p3d:eq:v2218-fixed-momentum-candidate} separate from the
previously reconstructed active Laurent branch and its physical
matching conditions.

For direct evaluation of all source contacts in the candidate, one may
use \eqref{p3d:eq:v2218-complete-orbit-jets} with
\begin{align}
 y'&=\frac{2C^2\zeta^2}{R^3}
           +\frac{2k\zeta}{HR},\nonumber\\
 y''&=-\frac{6C^2\zeta^2}{R^4}
       +\frac{8C^2\zeta}{HR^4}
       +\frac{2k}{R^2}(H^{-2}-\zeta/H),\nonumber\\
 \beta_Q'/\beta_Q&=2/R,
 \qquad \beta_Q''/\beta_Q=2/R^2.
 \label{p3d:eq:v2218-fixed-candidate-source-jets}
\end{align}

\begin{theorem}[Terminal activation and controlled ultraviolet value]
\label{p3d:thm:v2218-fixed-candidate-endpoints}
The invariant candidate \eqref{p3d:eq:v2218-fixed-momentum-candidate}
is smooth on its one-sided terminal domain and satisfies, with
$x=R-C$,
\begin{equation}
 \boxed{\begin{gathered}
 f_Q(R)=R-C+\frac{Qx^3}{H^2CD}+O(x^4),\\
 f_Q(C)=0,\quad f_Q'(C)=1,\quad f_Q''(C)=0,\quad
 f_Q'''(C)=\frac{6Q}{H^2CD}.
 \end{gathered}}
 \label{p3d:eq:v2218-fixed-candidate-terminal}
\end{equation}
For any band on which the moving pair has
$\lambda_-(s_R,y(R))\ge m_0>0$, put
$\delta=(1+e^{2m_0})^{-1}$.  Its complete nonlinear value obeys
\begin{equation}
 \boxed{(1-\delta)\frac{Q(R^2-C^2)\zeta^2}{D}
 \le f_Q-(R-C)
 \le\frac{Q(R^2-C^2)\zeta^2}{D}.}
 \label{p3d:eq:v2218-fixed-candidate-UV-bound}
\end{equation}
In particular, $0\le\sqrt{k}\,\zeta\le1$ gives
$\lambda_-\ge\min(a,s_R)-1$, providing a uniform value bound on
any band with $\min(a,s_R)>1$.  This is a finite-band estimate, not
a limit beyond the physical ceiling.
\end{theorem}
\begin{proof}
At the terminal, $\zeta=x/(HC)+O(x^2)$ and
$k=2x/C+O(x^2)$, so
$y=2x^3/(H^2C^3)+O(x^4)$.  Since
$s_C=-a$, \eqref{p3d:eq:v2218-orbit-zero-jets} gives
$\Phi(s_R,y)=y/8+O(xy)+O(y^2)$.
Multiplication by $\beta_Q(C)=4QC^2/D$ proves
\eqref{p3d:eq:v2218-fixed-candidate-terminal}.  Smoothness uses the
invariant variable $y$, not the non-differentiable square-root lift.

For the second claim let $M=\operatorname{diag}(a,s_R)$ and let $J$
be the off-diagonal two-by-two source with entries $\sqrt y$.
The whole segment $M+\theta J$, $0\le\theta\le1$, has spectrum
at least $m_0$: its endpoints do, and positivity of the shifted
matrices is convex.  On that spectral interval
$2(1-\delta)\le\Psi''\le2$.  The divided-difference formula for
the trace Hessian therefore bounds its second variation by these
constants times $\operatorname{Tr}J^2=2y$.
The first trace variation at $\theta=0$ is zero.  Twice integrating
in $\theta$ gives
\[
 2(1-\delta)y\le
 \operatorname{Tr}[\Psi(M+J)-\Psi(M)]\le2y.
\]
Multiplication by $\beta_Q/8$ and
$\beta_Qy/4=Q(R^2-C^2)\zeta^2/D$ proves the bound.
The last estimate is the elementary eigenvalue perturbation bound
$\lambda_{\min}(M+J)\ge\min(a,s_R)-\|J\|$.
\end{proof}

The ultraviolet estimate is a bound on the finite value.  Its use for
metric slow-roll derivatives consumes the complete jets
\eqref{p3d:eq:v2218-complete-orbit-jets} and
\eqref{p3d:eq:v2218-fixed-candidate-source-jets}, rather than
differentiating an error bound without further control.  Also, the
coefficient is proportional to $\zeta^2$, not an exactly constant
$R^2$ coefficient.  The location and stability of its finite
stationary points must be evaluated for this composed candidate.

\begin{proposition}[A direct finite-value dressing has a first-jet jump]
\label{p3d:prop:v2218-finite-value-junction}
Let $\mathscr W_0=(f_Q-(R-C))/Q$ and consider instead the direct
finite-value multiplication
\[
 \widetilde f(R)=R-C+D_2(R)\mathscr W_0(R),\qquad
 D_2(R)=Q^{n_2^+(R)}.
\]
At $B=C_2=U\alpha^{-2}$, the full first derivative has the nonzero jump
\begin{equation}
 \boxed{\widetilde f'(B+)-\widetilde f'(B-)
 =\log Q\,\frac{U-B}{2\log\alpha\,B(U+B)}
       \mathscr W_0(B)>0.}
 \label{p3d:eq:v2218-finite-value-junction}
\end{equation}
Thus this direct multiplication is not an ordinary $C^1$ metric
coefficient without an actual junction prescription.  The result concerns
the full nonzero finite value; it does not change the separate
susceptibility-integral construction in which $D_2$ multiplies a
Hessian.
\end{proposition}
\begin{proof}
At $B$, $n_2^{\rm L}=0$.  Differentiating the printed Laurent secant
gives
$n_2^{\rm L\,\prime}(B)=(U-B)/[2\log\alpha\,B(U+B)]>0$.
The positive-part branch has derivative zero below $B$ and that value
above it.  Since $D_2(B)=1$, only $D_2'\mathscr W_0$ jumps in
the product derivative.  The finite word is strictly positive there:
$C<B<U$, $\zeta(B)>0$, $k(B)>0$, and $\Psi$ is strictly convex.
This proves the formula, irrespective of the numerical size of the jump.
\end{proof}

\begin{remark}[Physical status of the composed candidate]
\label{p3d:rem:v2218-fixed-candidate-status}
The results above derive a finite scalar constructor and its complete
contacts from existing fixed-momentum data.  They neither select a
physical source section from all admissible sections nor establish a
new action-origin theorem.  A physical identification must still
evaluate the actual source/trace and Gram factorization, its metric
and density transport, the comparison of the finite $\Psi$ word with
the complete native operator/effective action, and its endpoint and
vacuum normalization.  Quantum elimination additionally retains its
actual determinant and measure terms; no such determinant is added to
a classical stationary word by notation.  A source orbit computable
from $V_\tau$ and a common action that physically consumes that orbit
are distinct claims.  The candidate does not replace the existing
active Laurent branch or assert a cosmological history connecting its
regimes.
\end{remark}

\section{Geometric response limits and gravitational action incidence}
\label{paper3d:sec:response-incidence}
\subsection{Grade-four carrier versus gravitational action incidence}
\label{p3d:sec:v229-actual-carrier-port}
\begin{remark}[The first actual carrier is not yet the consumed gravitational port]
\label{p3d:rem:v229-actual-gravitational-incidence}
The protected spectrum $\operatorname{spec}\mathcal Q_{\rm act}
=\{0,4,5,7\}$ gives $P_2=P_3=0$ and $P_4\ne0$.
Theorem~\ref{p3d:thm:v2219-carrier-four-dressing} selects $j=4$
within $\{2,3,4\}$ \emph{under its declared first-actual-carrier
rule}.  This is a valid finite geometric selection, not a derivation
that the physical metric coefficient consumes $P_4$.  That further
claim requires a typed, nonzero, preterminal insertion of the actual
$P_4$ sector into the gravitating action and the complete variation
of its curvature and interior blocks, in the sense of
Definition~\ref{p2a:def:v229-action-incidence-certificate}.
The two-port Cap Schur calculation admits distinct independent,
Born-antisymmetric and shared-symmetric incidence patterns; their
stationary returns are respectively $B_I(1-u^2)r^2$ and $B_Ir^2$.
A generic positivity/terminal-zero theorem does not choose between
them unless its required-output premise is proved for the actual
physical owner.  The resulting gravitational susceptibility is an
observable-facing question, not an index fitted to a desired action.
\end{remark}


\begin{corollary}[The exact grade-four heat-response test has an incidence input]
\label{p3d:cor:v286-protected-grade-four-incidence-test}
On the protected branch $P_4\ne0$, a fixed finite grade-four
compression $A_4>0$ and an actual, licensed scalar curvature/source
tangent $B_4$ on that same block define the positive mathematical
heat Gram of
Theorem~\ref{p3b:thm:v286-finite-heat-incidence-gram}. In the
locally affine, stationary finite-heat owner with complete interior
tangents $C_{4,i}$, its reduced grade-four second response is strictly
positive if and only if
\[
 B_4\notin\operatorname{span}_{\mathbb R}\{C_{4,i}\}.
\]
Its value is the Gram-distance square and determinant ratio of that
theorem.  If the action operator/metric incidence is nonlinear,
Proposition~\ref{p3b:prop:v286-heat-incidence-contact-obstruction}
requires its complete second contact.  Thus $P_4\ne0$ determines
neither a nonzero $B_4$ nor its nonzero \emph{reduced} action response,
and the abstract finite test alone does not identify the physical
coefficient $f_{RR}$.
\end{corollary}
\begin{proof}
Compress to the specified finite block and apply the positive Gram
and Schur formulas of
Theorem~\ref{p3b:thm:v286-finite-heat-incidence-gram}.
The strict-positivity equivalence follows from
orthogonal projection in that Gram. For a nonlinear insertion the
additional $-\tfrac12\operatorname{Tr}(e^{-A_4}X''_{4})$ term can
cancel the positive Gram by
Proposition~\ref{p3b:prop:v286-heat-incidence-contact-obstruction}. Actual action/density and variation
licensing are the separate requirements of
Definition~\ref{p2a:def:v229-action-incidence-certificate}.
\end{proof}


\subsection{Petersen finite invariants and the running-source firewall}
\label{p3d:sec:v319-petersen-running}
\begin{proposition}[Geometric ratios do not determine a beta function]
\label{p3d:prop:v319-petersen-no-running}
The signed Petersen eigenvalues $\pm\sqrt5,0$, the two-core polar
ratios $2:1/2$, and the complementary cycle flux ratio $\phi$
are invariants of a fixed finite marked geometry.  They do not depend
on a renormalization scale $\mu$ and, without a licensed action and a
specified renormalization prescription, imply neither
$\mu\,d g/d\mu$ nor a running mass or an anomalous dimension.
A uniform whole-orbit quadratic or quartic directional source is
isotropic by Theorem~\ref{p3b:thm:v319-petersen-moment-six}; its
first allowed nonconstant finite icosahedral directional invariant
within this moment family is the sextic $H_6$.
\end{proposition}
\begin{proof}
All numerical factors are eigenvalues or invariant ratios of finite
matrices with fixed entries in $\mathbb Q(\sqrt5)$.  Scale
derivatives vanish for those fixed matrices.  The assertion for
uniform directional sources is the exact moment identity in Paper~III-B.
No equality with an RG coupling or a momentum-dependent action
coefficient has been proved.
\end{proof}
\begin{remark}[Full-jet condition]
A marked finite source inserted into a nonlinear response must retain
its first and second background contacts.  In particular a nonzero
protected finite Gram response cannot substitute for the physical
reduced action Hessian: the complete-contact obstruction of
Proposition~\ref{p3b:prop:v286-heat-incidence-contact-obstruction}
and the common-action rule of
Theorem~\ref{p3d:thm:v2216-complete-action-jets} still apply.
\end{remark}

\subsection{Static radius homothety and physical curvature compatibility}
\label{p3d:sec:v229-radius-curvature}
\begin{proposition}[Static geometric compatibility of the proposed homothety]
\label{p3d:prop:v229-static-radius-test}
If the round radius $a_R=a_0/G(R)$ required by
Theorem~\ref{p1b:thm:v229-coupled-homothety} is also identified with
the physical static $S^3$ radius, its curvature obeys
$R=R_0[1+(R/R_s)^2]$, $R_0=6/a_0^2$.
This cannot hold on an open interval of $R$.
\end{proposition}
\begin{proof}
Static product curvature is $6/a_R^2$. Substitution gives a nonzero
quadratic polynomial in $R$, whose zero set contains no open interval.
\end{proof}
For closed FLRW, $R=6(\dot H+2H^2+a^{-2})$ in proper time, so the
static substitution is unavailable. An internal or auxiliary round
factor is likewise not automatically this physical spatial factor.
The comparison homothety does not vary the framework's fixed primitive
momentum radius. Its actual incidence with physical curvature remains
part of the same-action comparison, with all moving source contacts.


\section{Conclusions}
\label{paper3d:sec:organized-conclusion}
\label{p3d:sec:conclusion}
Paper~III-D has derived the protected rank valuation and exact two-root
response on the current fixed-seed constitutional stack.  The generated one-copy momentum seed \(\mathcal H_1\) is imported
upstream and is not re-selected here.  Universal full-endpoint
multiplicity-one remains conditional on the GP3 comparison.  The protected rank endpoint is
the theorem-level output
\(\mathsf C_{\rm IR}^{\rm der}\).  The matched
contact tensor gives the additive rank valuation.  Smooth completion retains
the explicit fixed-shell family
\(\mathfrak A_{\rm REG}=\{\mathrm{REG}_\sigma\}\); each frozen member gives
its own scheme-labelled completed weight, while \(\mathrm{REG}_0\) is only
the source prescription.  The positive determinant character and its
generated neutral dual character supply the dimensionless multiplicative
hierarchy directly.  The resulting native output is
dimensionless:
\[
 H_{56}^{(\sigma)}=\alpha_\sigma^{7}.
\]
The common-rate scalar coordinate gives seven ordered hierarchy levels,
including six finite hits and one asymptotic endpoint.  Seven ordered
rank-one projector events additionally require \((\mathrm{FLAG})\); without
it the scalar hierarchy, active depths \(q=0,4,5,7\), endpoint rank drop, and
exponent seven remain canonical.

For one fixed protected support, the exact \(8/k\)
conditioned-state density is theorem-level.  When the support is
orbit-invariant in the III-C generated active family, its occurrence is
inherited rather than resupplied; a particular rank-three member requires the
marked diagonal, and an arbitrary additional support remains a G0 boundary
occurrence.
The coefficient \(8/k\) is rank-indexed but the density remains support-indexed.
\((\mathrm{UNIT})\) only types that already derived state as a
boundary-conditioned response, not an extensive loop action.  The canonical
dimensionless spectral filtration is inherited from Paper~III-C without
\((\mathrm{ACT}_{\rm Cl})\); only an alternative or physically actualized
external-scale support law requires that interface.  Inside the
reduced class \((\mathrm{RED})\), the native variables are the two positive
scalar role outputs; no dimensionful reciprocal-radius or
scalar-circle metric is constructed.  The normalized role scalarizer and the two
abstract endpoint calibrations give the normalization-representative invariant
curve
\[
 g_{\rm eff}(\scal)=\frac{\scal^2+CU}{(C+U)\scal},
 \qquad
 \scal_B^2=CU.
\]
In the adopted \(\kappa=2\) representative,
\(\scal_{\rm resp}=C+U\); another positive response normalization is absorbed
into the corresponding conversion scale and gives the same calibrated curve.
The hyperbolic-cosine, \(\tanh\), Riccati, factorization, and secant formulas
are exact identities of this two-root response curve, not autonomous evolution
laws.  Identifying \(U\) with an absolute physical standard is an optional
later calibration.  Once the occurring response contains the required
generated owner-ancestry morphism and \(U\) is its upper reference section,
the ratios
\(C/U=\alpha_\sigma^{-7}\) and
\(X_w(q)/X_w(0)=\alpha_\sigma^{-wq}\) are inherited from the generated
neutral character rather than independently matched.  Their physical
orientation is decreasing only on \(\mathfrak A_{\rm REG}^{+}\).  On the
owner-placed branch with \(0<C<U\), this positive-character condition is not
an extra input: Proposition~\ref{p3d:prop:ordered-owner-joint-locus} proves it
is the exact joint co-realizability locus.  The critical and negative-character
members remain REG realizations elsewhere but cannot realize that ordered
owner-placed response.

On the common-rate neutral family \((\mathrm{NEUT}_7)\), the
closed heat loop is normalized to a boundary state and then passed through
the generated entropy functional, producing a unique basis-neutral
secant-matched comparison state.  The adopted response representative
\(g_{\rm eff}=2\lambda_{\rm B}(W_{\rm scal})\)
is scalar postprocessing; the factor \(2\) is normalization convention rather
than a physical selector.  When the later Paper-III-E \((\mathrm{MRESP})\)
composite occurs, constituent co-realization supplies its hierarchy-active
normalized state at that response's own scale parameter as a genuine
\emph{derived-later} specialization.  No equality with the secant-matched
comparison state is inferred, and the later state-to-Hessian coupling remains
an explicit response-class law.  Physical reciprocal-role statements consume the
premise-minimal Paper-II-C/GP2 reciprocal closure; \(\mathsf P_{\rm mc}\) is
only an optional G0 full-profile test certificate.
\((\mathrm{BR}_{\rm ind})\) and \((\mathrm{IND}_{xp})\) are upstream
certification labels, not native III-D antecedents.  Any later dimensionful
reciprocal-radius interpretation additionally carries its own calibration
boundary.  No optional \(\mathsf P_{\rm sli}\)/fixed-action momentum selector
is a native III-D antecedent.
The physical action, conditioned response, response-scale token, and
endpoint or hierarchy calibrations are all typed by G0; the family
\(\mathfrak A_{\rm REG}\) remains explicit analytic scheme data rather than a
consequence of G0, and no current discriminator selects its source member.
\printbibliography
\end{document}
